% La création d'une activité agropastorale — Allemand 1ère A — Cours signé OnBuch+
% Programme officiel MINESEC. Compiler avec : tectonic ALA08-la-creation-d-une-activite-agropastorale.tex
\newcommand{\DOCMATIERE}{Allemand}
\newcommand{\DOCNIVEAU}{1ère A}
\newcommand{\DOCMODULE}{Chapitre 2 : Vie économique}
\newcommand{\DOCLECON}{ALA08}
\newcommand{\DOCTITRE}{La création d'une activité agropastorale}
% OnBuch+ — préambule « cours signés » ENRICHI (compilé avec tectonic / XeLaTeX)
% Système visuel « Pop » d'OnBuch+ : fond crème (jamais blanc), bordures encre
% épaisses, ombres « dures » décalées, titres Archivo Black en capitales.
% Version MATHÉMATIQUES : graphes (pgfplots/tikz), boîtes pédagogiques
% (définition, propriété, méthode, attention, le savais-tu, exercice résolu,
% exercice, TP, bilan), page de garde et signature OnBuch+ ancrée en bas.
%
% Macros attendues dans main.tex AVANT \input{preamble} :
%   \DOCMATIERE  \DOCNIVEAU  \DOCMODULE  \DOCLECON (numéro)  \DOCTITRE
\documentclass[11pt]{article}

\usepackage{fontspec}
\usepackage[ngerman,french]{babel} % français = langue principale ; l'allemand via \all{..} / textebox
\usepackage[a4paper,margin=2cm,top=3.4cm,bottom=2.8cm,headheight=1.4cm,footskip=1.3cm]{geometry}
\usepackage{amsmath,amssymb,amsfonts}
\usepackage{mathtools} % \xmapsto, \coloneqq... souvent utilisés par les modèles
\usepackage{cancel} % simplifications barrées (bilans de forces)
% \abs{z} (module, valeur absolue) et \norm{u} (norme), utilisés par les modèles
\providecommand{\abs}{}\let\abs\relax\DeclarePairedDelimiter{\abs}{\lvert}{\rvert}
\providecommand{\norm}{}\let\norm\relax\DeclarePairedDelimiter{\norm}{\lVert}{\rVert}
\usepackage[table]{xcolor} % [table] : fournit aussi \rowcolors (alternance de lignes)
\usepackage{pagecolor}
\usepackage{tikz}
\usetikzlibrary{arrows.meta,positioning,calc,shapes.geometric,shapes.misc,shapes.arrows,decorations.pathmorphing,decorations.pathreplacing,decorations.markings,patterns,shadows,mindmap,trees,fit,backgrounds,matrix,intersections,angles,quotes}
\usepackage{pgfplots}
\usepackage[version=4]{mhchem} % équations nucléaires \ce{^{235}_{92}U}
\usepackage[european,siunitx]{circuitikz} % schémas électriques
\pgfplotsset{compat=1.18}
\usepgfplotslibrary{fillbetween,groupplots} % aires entre courbes (calcul intégral)
\usepackage{enumitem}
\usepackage{fancyhdr}
% [most] charge toutes les bibliothèques tcolorbox (listings, minted, raster,
% documentation...) et gonfle inutilement la mémoire TeX sur un document long
% (~300 boîtes) : on ne charge que ce qui est réellement utilisé.
\usepackage[skins,breakable]{tcolorbox}
\usepackage{graphicx}
\usepackage{colortbl}
\usepackage{booktabs}
\usepackage{tabularx}
\usepackage{diagbox} % case d'en-tête diagonale des tableaux à double entrée
% Colonnes extensibles centrée / gauche / droite (C, L, R), souvent utilisées par les modèles
\newcolumntype{C}{>{\centering\arraybackslash}X}
\newcolumntype{L}{>{\raggedright\arraybackslash}X}
\newcolumntype{R}{>{\raggedleft\arraybackslash}X}
\usepackage{array}
\usepackage{multicol}
\usepackage{multirow}
\usepackage{siunitx}
% parse-numbers=false : accepte aussi \SI{2,5 \times 10^{-3}}{...}
\sisetup{locale=FR,output-decimal-marker={,},group-minimum-digits=4,parse-numbers=false}
\DeclareSIUnit\ohm{\ensuremath{\symup{\Omega}}} % U+2126 absent de la police
\DeclareSIUnit\division{div} % réglages d'oscilloscope (ms/div, V/div)
\DeclareSIUnit\tr{tr} % tours (vitesse de rotation, ex. \SI{1500}{\tr\per\minute})
\DeclareSIUnit\molar{M} % concentration molaire (ex. \SI{3}{\molar}, \SI{1}{\micro\molar})
\DeclareSIUnit\an{an} % année (le modèle confond parfois avec \year, un compteur TeX) — ex. \SI{5}{\kilo\meter\per\an}
\DeclareSIUnit\FCFA{FCFA} % franc CFA (ex. \SI{500}{\FCFA\per\kilo\gram})
\DeclareSIUnit\degreeBrix{\textdegree Brix} % teneur en sucre d'un sirop/jus (réfractométrie)
\DeclareSIUnit\month{mois} % le modèle utilise parfois \month, un compteur TeX, comme unité de durée
\usepackage{qrcode}
% \degree : commande fréquemment utilisée par les modèles pour un angle ou une
% température en dehors de \SI{}{} (non fournie par les paquets chargés).
\providecommand{\degree}{^{\circ}}
% \qed (fin de démonstration), utilisé par les modèles sans amsthm
\providecommand{\qed}{\hfill$\square$}
% \barre : double barre des valeurs interdites dans un tableau de variations (tkz-tab)
\providecommand{\barre}{\ensuremath{\|}}
% environnement proof (amsthm non chargé) utilisé par les modèles
\ifdefined\proof\else\newenvironment{proof}[1][Démonstration]{\par\noindent\textit{#1.}\ }{\qed\par}\fi
\usepackage{lastpage}
\usepackage{hyperref}
\hypersetup{hidelinks}

% ============================================================
% Palette « Pop » OnBuch+ — mêmes hex que la classe P (plus_ui.dart)
% ============================================================
\definecolor{poporange}{HTML}{F59321}
\definecolor{popdark}{HTML}{E07A0C}
\definecolor{popcream}{HTML}{FFF6E8}
\definecolor{popink}{HTML}{1C1714}
\definecolor{poppurple}{HTML}{7A5AE0}
\definecolor{popgreen}{HTML}{2E9E6A}
\definecolor{poppink}{HTML}{E0457B}
\definecolor{popblue}{HTML}{2F7DE1}
\definecolor{popgold}{HTML}{C98A12}
\definecolor{popmuted}{HTML}{8A7F76}
\definecolor{popline}{HTML}{EADFD2}
\definecolor{popgray}{HTML}{8A7F76}
\colorlet{popgrey}{popgray}
% Alias tolérants (le modèle nomme parfois les couleurs différemment)
\colorlet{popwhite}{white}
\colorlet{popblack}{popink}
\colorlet{popred}{poppink}
\colorlet{popyellow}{popgold}
% Teintes claires (fonds de tableaux/encadrés) — le modèle nomme parfois
% les couleurs avec un suffixe L (ex. popblueL) sans les avoir définies.
\colorlet{poporangeL}{poporange!20}
\colorlet{popdarkL}{popdark!20}
\colorlet{poppurpleL}{poppurple!20}
\colorlet{popgreenL}{popgreen!20}
\colorlet{poppinkL}{poppink!20}
\colorlet{popblueL}{popblue!20}
\colorlet{popgoldL}{popgold!20}
\colorlet{popredL}{popred!20}
\colorlet{popgrayL}{popgray!20}
% Teintes claires pour fonds de tableaux / zones
\colorlet{popblueL}{popblue!10!white}
\colorlet{poporangeL}{poporange!14!white}
\colorlet{popgreenL}{popgreen!12!white}
\colorlet{poppurpleL}{poppurple!12!white}
\colorlet{poppinkL}{poppink!12!white}
\colorlet{popgoldL}{popgold!14!white}

\pagecolor{popcream}

% ============================================================
% Typographie
% ============================================================
\defaultfontfeatures{Ligatures=TeX}
\setmainfont{PlusJakartaSans-Medium.ttf}[Path=../../../fonts/,BoldFont=PlusJakartaSans-Bold.ttf]
% Maths en sans-serif (Fira Math) avec chiffres et lettres droites de Plus
% Jakarta, pour que les formules se fondent dans le texte.
\usepackage[math-style=ISO,bold-style=ISO]{unicode-math}
\setmathfont{FiraMath-Regular.otf}[Path=../../../fonts/]
\setmathfont{PlusJakartaSans-Medium.ttf}[Path=../../../fonts/,range={up/num,up/{Latin,latin},\mathup}]
\usepackage{icomma} % virgule décimale sans espace en mode math
% Caractères mathématiques Unicode absents des polices de texte (Plus Jakarta,
% Archivo Black) : rendus par leur équivalent mathématique, en texte comme en
% formule — sinon ils s'affichent en carré vide (ex. « Arithmétique dans ℤ »).
\usepackage{newunicodechar}
\newunicodechar{ℝ}{\ensuremath{\mathbb{R}}}
\newunicodechar{ℤ}{\ensuremath{\mathbb{Z}}}
\newunicodechar{ℂ}{\ensuremath{\mathbb{C}}}
\newunicodechar{ℕ}{\ensuremath{\mathbb{N}}}
\newunicodechar{ℚ}{\ensuremath{\mathbb{Q}}}
\newunicodechar{𝔻}{\ensuremath{\mathbb{D}}}
\newunicodechar{α}{\ensuremath{\alpha}}
\newunicodechar{β}{\ensuremath{\beta}}
\newunicodechar{γ}{\ensuremath{\gamma}}
\newunicodechar{δ}{\ensuremath{\delta}}
\newunicodechar{ε}{\ensuremath{\varepsilon}}
\newunicodechar{θ}{\ensuremath{\theta}}
\newunicodechar{λ}{\ensuremath{\lambda}}
\newunicodechar{μ}{\ensuremath{\mu}}
\newunicodechar{σ}{\ensuremath{\sigma}}
\newunicodechar{τ}{\ensuremath{\tau}}
\newunicodechar{φ}{\ensuremath{\varphi}}
\newunicodechar{ψ}{\ensuremath{\psi}}
\newunicodechar{ω}{\ensuremath{\omega}}
\newunicodechar{Σ}{\ensuremath{\Sigma}}
% majuscules produites par \MakeUppercase dans les titres (α-aminés -> Α)
\newunicodechar{Α}{\ensuremath{\alpha}}
\newunicodechar{Β}{\ensuremath{\beta}}
\newunicodechar{Γ}{\ensuremath{\Gamma}}
\newunicodechar{Δ}{\ensuremath{\Delta}}
\newunicodechar{Ε}{\ensuremath{\varepsilon}}
\newunicodechar{Θ}{\ensuremath{\Theta}}
\newunicodechar{Λ}{\ensuremath{\Lambda}}
\newunicodechar{Μ}{\ensuremath{\mu}}
\newunicodechar{Τ}{\ensuremath{\tau}}
\newunicodechar{Φ}{\ensuremath{\Phi}}
\newunicodechar{Ψ}{\ensuremath{\Psi}}
\newunicodechar{Ω}{\ensuremath{\Omega}}
\newunicodechar{↦}{\ensuremath{\mapsto}}
\newunicodechar{∈}{\ensuremath{\in}}
\newunicodechar{∉}{\ensuremath{\notin}}
\newunicodechar{∧}{\ensuremath{\wedge}}
\newunicodechar{⇔}{\ensuremath{\Leftrightarrow}}
\newunicodechar{⟺}{\ensuremath{\Longleftrightarrow}}
\newunicodechar{⇒}{\ensuremath{\Rightarrow}}
\newunicodechar{⟹}{\ensuremath{\Longrightarrow}}
\newunicodechar{∘}{\ensuremath{\circ}}
\newunicodechar{∪}{\ensuremath{\cup}}
\newunicodechar{∩}{\ensuremath{\cap}}
\newunicodechar{≡}{\ensuremath{\equiv}}
\newunicodechar{⊂}{\ensuremath{\subset}}
\newunicodechar{⊥}{\ensuremath{\perp}}
\newunicodechar{∃}{\ensuremath{\exists}}
\newunicodechar{∀}{\ensuremath{\forall}}
\newunicodechar{∛}{\ensuremath{\sqrt[3]{\,}}}
\newunicodechar{∜}{\ensuremath{\sqrt[4]{\,}}}
\newunicodechar{⃗}{\ensuremath{{}^{\to}}}
\newunicodechar{ⁿ}{\ifmmode^{n}\else\textsuperscript{\ensuremath{n}}\fi}
\newunicodechar{ˣ}{\ifmmode^{x}\else\textsuperscript{\ensuremath{x}}\fi}
\newunicodechar{ᵇ}{\ifmmode^{b}\else\textsuperscript{\ensuremath{b}}\fi}
\newunicodechar{ʳ}{\ifmmode^{r}\else\textsuperscript{\ensuremath{r}}\fi}
\newunicodechar{ᵘ}{\ifmmode^{u}\else\textsuperscript{\ensuremath{u}}\fi}
\newunicodechar{ʸ}{\ifmmode^{y}\else\textsuperscript{\ensuremath{y}}\fi}
\newunicodechar{ᵃ}{\ifmmode^{a}\else\textsuperscript{\ensuremath{a}}\fi}
\newunicodechar{ᵉ}{\ifmmode^{e}\else\textsuperscript{\ensuremath{e}}\fi}
\newunicodechar{ᵏ}{\ifmmode^{k}\else\textsuperscript{\ensuremath{k}}\fi}
\newunicodechar{ᵖ}{\ifmmode^{p}\else\textsuperscript{\ensuremath{p}}\fi}
\newunicodechar{ᴺ}{\ifmmode^{N}\else\textsuperscript{\ensuremath{N}}\fi}
\newunicodechar{ᶜ}{\ifmmode^{c}\else\textsuperscript{\ensuremath{c}}\fi}
\newunicodechar{ʰ}{\ifmmode^{h}\else\textsuperscript{\ensuremath{h}}\fi}
\newunicodechar{ᵠ}{\ifmmode^{q}\else\textsuperscript{\ensuremath{q}}\fi}
\newunicodechar{ₙ}{\ifmmode_{n}\else\textsubscript{\ensuremath{n}}\fi}
\newunicodechar{ₐ}{\ifmmode_{a}\else\textsubscript{\ensuremath{a}}\fi}
\newunicodechar{ᵢ}{\ifmmode_{i}\else\textsubscript{\ensuremath{i}}\fi}
\newunicodechar{ᵣ}{\ifmmode_{r}\else\textsubscript{\ensuremath{r}}\fi}
\newunicodechar{ₖ}{\ifmmode_{k}\else\textsubscript{\ensuremath{k}}\fi}
\newunicodechar{ᵦ}{\ifmmode_{\beta}\else\textsubscript{\ensuremath{\beta}}\fi}
\newunicodechar{ₓ}{\ifmmode_{x}\else\textsubscript{\ensuremath{x}}\fi}
\newfontfamily\popdisplay{ArchivoBlack-Regular.ttf}[Path=../../../fonts/]
\newfontfamily\poplogo{SpaceGrotesk-Bold.ttf}[Path=../../../fonts/]
\newfontfamily\popsemi{PlusJakartaSans-SemiBold.ttf}[Path=../../../fonts/]
\newfontfamily\popxbold{PlusJakartaSans-ExtraBold.ttf}[Path=../../../fonts/]
\newfontfamily\popxboldls{PlusJakartaSans-ExtraBold.ttf}[Path=../../../fonts/,LetterSpace=3.0]

\setlist[enumerate]{leftmargin=1.7em,itemsep=5pt,topsep=4pt}
\setlist[itemize]{leftmargin=1.7em,itemsep=4pt,topsep=4pt,label={\color{poporange}\textbullet}}
\setlength{\parindent}{0pt}
\setlength{\parskip}{5pt}
\renewcommand{\arraystretch}{1.25}

% pgfplots : style Pop par défaut
\pgfplotsset{
  every axis/.append style={axis line style={popink,line width=1pt},tick style={popink},
    grid style={popline,line width=0.5pt},label style={font=\small},tick label style={font=\footnotesize}},
  popcurve/.style={poporange,line width=1.8pt,no markers,smooth},
}
% Même style disponible dans un \draw TikZ simple (hors pgfplots)
\tikzset{popcurve/.style={poporange,line width=1.8pt}}
\tikzset{popgrid/.style={popline,very thin}}
\colorlet{popcurve}{poporange} % le nom du style est parfois utilisé comme couleur

% ============================================================
% Pill / squiggle
% ============================================================
\newcommand{\poppill}[3]{%
  \tikz[baseline=-0.55ex]\node[draw=popink,line width=1.2pt,rounded corners=2.6pt,%
    fill=#2,inner xsep=6.5pt,inner ysep=3pt]{%
    \popxboldls\fontsize{7}{7}\selectfont\color{#3}\MakeUppercase{#1}};%
}
\newcommand{\popsquiggle}[1]{%
  \tikz[baseline=-0.4ex]{
    \draw[#1,line width=2.6pt,line cap=round]
      plot [smooth] coordinates {(0,0.09) (0.34,0.01) (0.68,0.21) (1.02,0.01) (1.36,0.10)};
  }%
}
% Mot-clé surligné (marqueur orange sous le texte)
\newcommand{\cle}[1]{{\popxbold\color{popdark}#1}}

% ============================================================
% En-tête / pied
% ============================================================
\pagestyle{fancy}
\fancyhf{}
\renewcommand{\headrulewidth}{0pt}
\renewcommand{\footrulewidth}{0pt}
\lhead{\raisebox{-0.15ex}{\poplogo\fontsize{13}{13}\selectfont\color{popink}OnBuch\color{poporange}+}}
\rhead{\poppill{\DOCMATIERE\ \textbullet\ \DOCNIVEAU\ \textbullet\ Leçon \DOCLECON}{white}{popink}}
\lfoot{\popxbold\fontsize{7.6}{7.6}\selectfont\color{popmuted}\MakeUppercase{Cours signé OnBuch+}\ \textbullet\ {\popsemi\DOCTITRE}}
\rfoot{\tikz[baseline=-0.6ex]\node[rounded corners=8.5pt,fill=popink,minimum height=17pt,minimum width=17pt,inner xsep=5pt,inner sep=0]{\popxbold\fontsize{8}{8}\selectfont\color{popcream}\thepage\,/\,\pageref*{LastPage}};}

% ============================================================
% Titres
% ============================================================
\newcommand{\chaptitre}[2]{%
  \begin{tcolorbox}[enhanced,colback=poporange,colframe=popink,boxrule=2.6pt,arc=11pt,
      drop shadow=popink,left=15pt,right=15pt,top=12pt,bottom=11pt,before skip=3pt,after skip=18pt]
    \poppill{#2}{popink}{popcream}\par\vspace{9pt}
    {\raggedright\hyphenpenalty=10000\exhyphenpenalty=10000\popdisplay\fontsize{22}{24}\selectfont\color{popcream}\MakeUppercase{#1}\par}
  \end{tcolorbox}
}
% Section numérotée : pastille orange avec le numéro + titre Archivo
\newcounter{popsec}
\newcommand{\coursec}[1]{%
  \stepcounter{popsec}\setcounter{popsub}{0}%
  \par\vspace{16pt}\noindent
  \begin{minipage}{\linewidth}
  \tikz[baseline=-0.7ex]\node[draw=popink,line width=1.6pt,fill=poporange,rounded corners=4pt,minimum size=22pt,inner sep=0]{\popdisplay\fontsize{12}{12}\selectfont\color{popcream}\thepopsec};%
  \hspace{8pt}{\popdisplay\fontsize{15}{17}\selectfont\color{popink}\MakeUppercase{#1}}\par
  \vspace{4pt}\noindent{\color{poporange}\rule{36pt}{3.6pt}}
  \end{minipage}\par\nopagebreak\vspace{8pt}
}
\newcounter{popsub}
\newcommand{\courssub}[1]{%
  \stepcounter{popsub}%
  \par\vspace{9pt}\noindent{\popxbold\fontsize{11.5}{13}\selectfont\color{popink}{\color{poporange}\thepopsec.\thepopsub}\hspace{6pt}#1}\par\nopagebreak\vspace{4pt}
}

% ============================================================
% Boîtes pédagogiques — même recette : fond blanc, bordure épaisse colorée,
% ombre dure encre, badge plein. Titre optionnel [#1].
% ============================================================
\tcbset{popbase/.style={breakable,enhanced,colback=white,boxrule=2.3pt,arc=9pt,
  shadow={3pt}{-3pt}{0pt}{popink},left=14pt,right=14pt,top=11pt,bottom=11pt,before skip=11pt,after skip=15pt,
  fontupper=\popsemi\fontsize{10.6}{14.5}\selectfont\color{popink},
  fontlower=\popsemi\fontsize{10.6}{14.5}\selectfont\color{popink}}}
% Coche dessinée (le glyphe ✓ n'existe pas dans les polices du document)
\newcommand{\popcheck}[1]{\tikz[baseline=-0.5ex]\draw[#1,line width=1.6pt,line cap=round,line join=round](-0.13,0.0)--(-0.04,-0.09)--(0.13,0.1);}
\providecommand{\coche}{\popcheck{popgreen}}
\providecommand{\resultat}[1]{\par\smallskip\noindent\fcolorbox{popgreen}{popgreenL}{\parbox{\dimexpr\linewidth-2\fboxsep-2\fboxrule\relax}{\textbf{Résultat.} #1}}\par\smallskip}
\newcommand{\popboxhead}[3]{\poppill{#1}{#2}{white}\ifx\relax#3\relax\else\hspace{6pt}{\popxbold\fontsize{10.6}{12}\selectfont\color{popink}#3}\fi\par\vspace{7pt}}

\newtcolorbox{aretenir}[1][]{popbase,colframe=popgreen,before upper={\popboxhead{À retenir}{popgreen}{#1}}}
% Texte en allemand (document de lecture, dialogue, extrait) : cadre
% d'étude. Un titre optionnel (source, auteur) s'affiche dans l'en-tête.
\newtcolorbox{textebox}[1][]{popbase,colframe=popblue,before upper={\popboxhead{Text}{popblue}{#1}\begin{otherlanguage}{ngerman}},after upper={\end{otherlanguage}}}
% Marques ✓ / ✗ (forme correcte / fautive) souvent utilisées par les modèles
\providecommand{\cross}{\textcolor{poppink}{\ensuremath{\times}}}
\providecommand{\xmark}{\cross}\providecommand{\crossmark}{\cross}\providecommand{\cmark}{\ensuremath{\checkmark}}
% Ligne de réponse / justification à compléter (inventée par les modèles)
\providecommand{\justification}[1]{\par\noindent\textit{Justification~:} \dotfill\par}
% \textipa (package tipa non chargé) : la police ne couvre pas l'API -> simple passage
\providecommand{\textipa}[1]{#1}
% Soulignement (ulem non chargé) : \uline, \ul
\providecommand{\uline}[1]{\underline{#1}}\providecommand{\ul}[1]{\underline{#1}}
% \alert (beamer) inventée par les modèles : texte mis en évidence
\providecommand{\alert}[1]{\textbf{\textcolor{poppink}{#1}}}
% variantes ASCII d'ordinaux écrites en macro
\providecommand{\iere}{\textsuperscript{re}}\providecommand{\ieme}{\textsuperscript{e}}\providecommand{\ere}{\textsuperscript{re}}\providecommand{\eme}{\textsuperscript{e}}\providecommand{\er}{\textsuperscript{er}}\providecommand{\ie}{\textsuperscript{e}}
% Ordinaux français écrits en macro par les modèles : 1\ière, 3\ième
\newcommand{\ière}{\textsuperscript{re}}\newcommand{\ième}{\textsuperscript{e}}
% \exemple (inventée par les modèles) : étiquette « Exemple : »
\providecommand{\exemple}{\textbf{Exemple~:}\ }
% \square utilisable aussi hors mode math (cases à cocher)
\AtBeginDocument{\let\squareMath\square\renewcommand{\square}{\ensuremath{\squareMath}}}
% Mot ou phrase en allemand à l'intérieur d'un paragraphe français
\newcommand{\all}[1]{\ifmmode\text{\textit{#1}}\else\begingroup\selectlanguage{ngerman}\itshape #1\endgroup\fi}
\let\esp\all % alias toléré
\newtcolorbox{exemplebox}[1][]{popbase,colframe=poporange,before upper={\popboxhead{Exemple}{poporange}{#1}}}
\newtcolorbox{definition}[1][]{popbase,colframe=popblue,before upper={\popboxhead{Définition}{popblue}{#1}}}
\newtcolorbox{propriete}[1][]{popbase,colframe=poppurple,before upper={\popboxhead{Propriété}{poppurple}{#1}}}
\newtcolorbox{methode}[1][]{popbase,colframe=popgold,before upper={\popboxhead{Méthode}{popgold}{#1}}}
\newtcolorbox{attention}[1][]{popbase,colframe=poppink,before upper={\popboxhead{Attention}{poppink}{#1}}}
\newtcolorbox{experience}[1][]{popbase,colframe=popblue,colback=popblueL,before upper={\popboxhead{Expérience}{popblue}{#1}}}
% « Le savais-tu ? » : carte violette pleine (contexte, culture, Cameroun)
\newtcolorbox{savaistu}[1][]{popbase,colback=poppurple,colframe=popink,
  fontupper=\popsemi\fontsize{10.4}{14.2}\selectfont\color{white},
  before upper={\poppill{Le savais-tu\,?}{white}{poppurple}\ifx\relax#1\relax\else\hspace{6pt}{\popxbold\color{white}#1}\fi\par\vspace{7pt}}}
% Exercice résolu : énoncé en haut, \tcblower puis la solution sur fond crème
\newcounter{popexo}\newcounter{popexr}
\newtcolorbox{exoresolu}[1][]{popbase,colframe=poporange,colbacklower=popcream,
  segmentation style={popink,line width=1.2pt,dashed},
  before upper={\stepcounter{popexr}\popboxhead{Exercice résolu \thepopexr}{poporange}{#1}},
  before lower={\poppill{Solution}{popink}{popcream}\par\vspace{6pt}}}
% Exercice (non résolu en place) — la correction est dans la partie Corrigés
\newtcolorbox{exercice}[1][]{popbase,colframe=popink,
  before upper={\stepcounter{popexo}\popboxhead{Exercice \thepopexo}{popink}{#1}}}
% Corrigé
\newtcolorbox{corrige}[1][]{popbase,colframe=popgreen,colback=popgreenL,
  before upper={\popboxhead{Corrigé}{popgreen}{#1}}}
% Objectifs / prérequis / bilan
\newtcolorbox{objectifs}{popbase,colframe=popink,colback=poporangeL,
  before upper={\popboxhead{Objectifs de la leçon}{popink}{}}}
\newtcolorbox{prerequis}{popbase,colframe=popink,colback=popblueL,
  before upper={\popboxhead{Prérequis}{popblue}{}}}
\newtcolorbox{bilan}{popbase,colframe=popink,colback=white,boxrule=2.8pt,
  before upper={\popboxhead{Fiche bilan}{poporange}{L'essentiel à connaître pour l'examen}}}
% Figure encadrée avec légende
\newtcolorbox{popfigure}[1][]{enhanced,breakable=false,colback=white,colframe=popline,boxrule=1.4pt,arc=8pt,
  left=10pt,right=10pt,top=10pt,bottom=8pt,before skip=10pt,after skip=12pt,halign=center,
  lower separated=false,halign lower=center,
  fontlower=\popsemi\fontsize{9}{11.5}\selectfont\color{popmuted},
  #1}
\newcommand{\legende}[1]{\par\vspace{6pt}{\popsemi\fontsize{9}{11.5}\selectfont\color{popmuted}#1}}

% ============================================================
% Signature OnBuch+
% ============================================================
\newcommand{\poplogoinline}[1]{{\poplogo\fontsize{#1}{#1}\selectfont\color{popink}OnBuch\color{poporange}+}}
\newcommand{\signatureonbuch}{%
  \begin{tcolorbox}[enhanced,colback=popink,colframe=popink,boxrule=2.2pt,arc=10pt,
      drop shadow=poporange,left=14pt,right=14pt,top=10pt,bottom=10pt,before skip=0pt,after skip=0pt]
    \tikz[baseline=-0.7ex]\node[circle,fill=popgreen,draw=popcream,line width=1.3pt,minimum size=22pt,inner sep=0]{\popcheck{white}};%
    \hspace{9pt}\begin{minipage}[c]{0.62\linewidth}
      {\popxbold\fontsize{12}{13}\selectfont\color{popcream}Signé\ \textbullet\ {\poplogo OnBuch\color{poporange}+}}\par\vspace{2pt}
      {\raggedright\popsemi\fontsize{8}{10.5}\selectfont\color{popline}\MakeUppercase{Équipe pédagogique OnBuch+}\\\MakeUppercase{Conforme au programme officiel MINESEC}\par}
    \end{minipage}\hfill
    {\poplogo\fontsize{22}{22}\selectfont\color{popcream}OnBuch\color{poporange}+}
  \end{tcolorbox}%
}

% ============================================================
% Page de garde
% #1 titre · #2 matière · #3 niveau · #4 module · #5 n° leçon · #6 durée ·
% #7 description / objectifs (texte ou itemize)
% ============================================================
\newcommand{\pagegarde}[7]{%
  \thispagestyle{empty}
  \newgeometry{a4paper,margin=1.7cm,top=1.6cm,bottom=1.4cm}%
  \begin{tikzpicture}[remember picture,overlay]
    \fill[poporange!18] ([xshift=-3.2cm,yshift=-3.6cm]current page.north east) circle (4.6cm);
    \fill[poppurple!14] ([xshift=2.4cm,yshift=7.6cm]current page.south west) circle (3.8cm);
  \end{tikzpicture}%
  {\poplogo\fontsize{30}{30}\selectfont\color{popink}OnBuch\color{poporange}+}\hfill
  \raisebox{0.6ex}{\poppill{Cours signé \textbullet\ Premium}{popink}{popcream}}\par
  \vspace{1pt}\popsquiggle{poppurple}\par
  \vspace{8pt}
  \begin{tcolorbox}[enhanced,colback=poporange,colframe=popink,boxrule=2.8pt,arc=13pt,
      drop shadow=popink,left=18pt,right=18pt,top=12pt,bottom=13pt,after skip=10pt]
    \poppill{#2 \textbullet\ #3}{popink}{popcream}\hspace{5pt}\poppill{#4}{white}{popink}\par\vspace{8pt}
    {\popdisplay\fontsize{14}{14}\selectfont\color{popink}LEÇON #5}\par\vspace{4pt}
    {\raggedright\hyphenpenalty=10000\exhyphenpenalty=10000\popdisplay\fontsize{25}{27}\selectfont\color{popcream}\MakeUppercase{#1}\par}
    \vspace{8pt}
    {\popxbold\fontsize{9.5}{9.5}\selectfont\color{popink}\MakeUppercase{Durée conseillée : #6}}
  \end{tcolorbox}
  \begin{tcolorbox}[enhanced,colback=white,colframe=popink,boxrule=2.2pt,arc=10pt,
      drop shadow=popink,left=16pt,right=16pt,top=10pt,bottom=10pt,after skip=8pt]
    \poppill{Dans cette leçon}{poppurple}{white}\par\vspace{5pt}
    \popsemi\fontsize{9.3}{12.6}\selectfont\color{popink}#7
  \end{tcolorbox}
  \noindent\begin{minipage}[t]{0.487\linewidth}
    \begin{tcolorbox}[enhanced,colback=popgreen,colframe=popink,boxrule=2.2pt,arc=10pt,
        drop shadow=popink,left=11pt,right=11pt,top=10pt,bottom=10pt,height=3.1cm,valign=center]
      \tcbox[colback=white,colframe=popink,boxrule=1.3pt,arc=5pt,boxsep=0pt,nobeforeafter,
        left=3pt,right=3pt,top=3pt,bottom=3pt,valign=center]{\qrcode[height=1.9cm]{https://play.google.com/store/apps/details?id=com.luvvix.onbuch&hl=fr}}%
      \hspace{8pt}\begin{minipage}[c]{0.5\linewidth}
        {\popdisplay\fontsize{11}{12.5}\selectfont\color{white}\MakeUppercase{Télécharge OnBuch}}\par\vspace{4pt}
        {\raggedright\popsemi\fontsize{8.4}{11}\selectfont\color{white}Gratuit \textbullet\ Google Play\\Tuteur IA, annales, concours, résultats\par}
      \end{minipage}
    \end{tcolorbox}
  \end{minipage}\hfill
  \begin{minipage}[t]{0.487\linewidth}
    \begin{tcolorbox}[enhanced,colback=poppurple,colframe=popink,boxrule=2.2pt,arc=10pt,
        drop shadow=popink,left=11pt,right=11pt,top=10pt,bottom=10pt,height=3.1cm,valign=center]
      \tikz[baseline=-0.7ex]\node[circle,fill=white,draw=popink,line width=1.3pt,minimum size=1.6cm,inner sep=0]{\popdisplay\fontsize{24}{24}\selectfont\color{poppurple}+};%
      \hspace{8pt}\begin{minipage}[c]{0.6\linewidth}
        {\popdisplay\fontsize{11}{12.5}\selectfont\color{white}\MakeUppercase{Passe à OnBuch+}}\par\vspace{4pt}
        {\raggedright\popsemi\fontsize{8.4}{11}\selectfont\color{white}Tous les cours signés, Tuteur IA illimité, fiches PDF — onglet OnBuch+ de l'app\par}
      \end{minipage}
    \end{tcolorbox}
  \end{minipage}
  \vspace{8pt}
  \popcontenu
  \vfill
  \signatureonbuch
  \restoregeometry
  \newpage
}

% Rangée « Ce cours contient » (page de garde)
\newcommand{\poptile}[3]{% #1 couleur #2 gros texte #3 légende
  \begin{tcolorbox}[enhanced,colback=#1,colframe=popink,boxrule=2pt,arc=9pt,shadow={2.5pt}{-2.5pt}{0pt}{popink},
      left=6pt,right=6pt,top=9pt,bottom=9pt,halign=center,valign=center,height=1.75cm,width=\linewidth,nobeforeafter]
    {\popdisplay\fontsize{12.5}{13}\selectfont\color{white}\MakeUppercase{#2}}\par\vspace{3pt}
    {\centering\popsemi\fontsize{7.8}{9.5}\selectfont\color{white}#3\par}
  \end{tcolorbox}}
\newcommand{\popcontenu}{%
  \par\noindent{\popxboldls\fontsize{7.5}{7.5}\selectfont\color{popmuted}CE COURS SIGNÉ CONTIENT}\par\vspace{7pt}
  \noindent\begin{minipage}[t]{0.235\linewidth}\poptile{popblue}{Cours}{complet, illustré, pas à pas}\end{minipage}\hfill
  \begin{minipage}[t]{0.235\linewidth}\poptile{poporange}{Méthodes}{et exercices résolus}\end{minipage}\hfill
  \begin{minipage}[t]{0.235\linewidth}\poptile{poppink}{Exercices}{type examen + corrigés}\end{minipage}\hfill
  \begin{minipage}[t]{0.235\linewidth}\poptile{popgreen}{Fiche bilan}{l'essentiel à retenir}\end{minipage}\par}

% Sommaire
\newtcolorbox{sommaire}{enhanced,colback=white,colframe=popink,boxrule=2.2pt,arc=10pt,
  drop shadow=popink,left=18pt,right=18pt,top=13pt,bottom=13pt,before skip=6pt,after skip=20pt,
  before upper={\poppill{Sommaire}{poppurple}{white}\par\vspace{9pt}\popsemi\fontsize{10.6}{14.5}\selectfont\color{popink}}}

% Signature de fin de cours — poussée tout en bas de la dernière page
\newcommand{\findecours}{%
  \par\vfill\nopagebreak
  \begin{center}\popsquiggle{poporange}\end{center}\vspace{4pt}
  \signatureonbuch
}
\AtBeginDocument{\def\llbracket{\mathopen{[\mkern-3.5mu[}}\def\rrbracket{\mathclose{]\mkern-3.5mu]}}} % intervalles d'entiers (glyphe ⟦ absent de la police)
% Glyphes absents de Fira Math : redéfinis à partir de glyphes disponibles
\AtBeginDocument{%
  \def\setminus{\mathbin{\backslash}}\let\smallsetminus\setminus
  \def\vdots{\mathord{\vbox{\baselineskip3.2pt\lineskiplimit0pt\kern4pt\hbox{.}\hbox{.}\hbox{.}}}}%
  \def\ddots{\mathinner{\mkern1mu\raise6.5pt\hbox{.}\mkern2mu\raise3.6pt\hbox{.}\mkern2mu\raise0.7pt\hbox{.}\mkern1mu}}%
  \def\triangle{\mathord{\scalebox{0.9}{$\symup{\Delta}$}}}%
  \def\nabla{\mathord{\raisebox{\depth}{\scalebox{1}[-1]{$\symup{\Delta}$}}}}%
  \def\checkmark{\popcheck{popdark}}%
  \def\star{\mathbin{\ast}}\def\bigstar{\mathord{\ast}}%
  \def\diamond{\mathbin{\scalebox{0.6}{$\Diamond$}}}%
  \def\bigcup{\mathop{\mathchoice{\raisebox{-0.3em}{\scalebox{1.7}{$\cup$}}}{\scalebox{1.25}{$\cup$}}{\cup}{\cup}}\displaylimits}%
  \def\bigcap{\mathop{\mathchoice{\raisebox{-0.3em}{\scalebox{1.7}{$\cap$}}}{\scalebox{1.25}{$\cap$}}{\cap}{\cap}}\displaylimits}%
  \def\lesseqqgtr{\mathrel{\lessgtr}}\def\bigoplus{\mathop{\oplus}\displaylimits}%
}
\providecommand{\ssi}{si et seulement si} % abréviation fréquente
\AtBeginDocument{\providecommand{\wideparen}{\overparen}} % arc de cercle

%% FIGFIX-BEGIN v4 — correctifs globaux des figures (docs/onbuch-generation/tools/figfix)
\usepackage{adjustbox}\usepackage{etoolbox}
% toute figure TikZ/pgfplots plus large que la ligne est réduite à la largeur disponible
\BeforeBeginEnvironment{tikzpicture}{\begin{adjustbox}{max width=\linewidth}}
\AfterEndEnvironment{tikzpicture}{\end{adjustbox}}
% légendes pgfplots lisibles par-dessus les courbes
\pgfplotsset{every axis/.append style={legend style={fill=white,fill opacity=0.92,text opacity=1,draw=black!15,inner sep=3pt}}}
% étiquettes d'arêtes (syntaxe edge["..."]) avec fond blanc
\tikzset{every edge quotes/.append style={fill=white,inner sep=1.2pt}}
%% FIGFIX-END
\begin{document}

\pagegarde{La création d'une activité agropastorale}{Allemand}{1ère A}{Chapitre 2 : Vie économique}{ALA08}{2 h}{Cette leçon propose la lecture et l'exploitation d'un texte sur la création d'une activité agropastorale. Elle permet d'acquérir le vocabulaire de la ferme et de l'élevage, de maîtriser le futur I et la construction finale « um ... zu », et de produire un court texte de projet agricole.
\vspace{4pt}
\begin{multicols}{2}\small
\begin{itemize}
\item À la fin de cette leçon, je sais comprendre globalement puis en détail un texte allemand sur la création d'une activité agropastorale.
\item À la fin de cette leçon, je sais employer correctement le vocabulaire thématique (die Farm, das Vieh, der Mais, das Feld, anbauen, züchten) avec articles et pluriels.
\item À la fin de cette leçon, je sais former et conjuguer le futur I (werden + infinitif) pour parler de projets.
\item À la fin de cette leçon, je sais exprimer la finalité avec « um ... zu + infinitif ».
\item À la fin de cette leçon, je sais distinguer futur I et présent à valeur de futur.
\item À la fin de cette leçon, je sais répondre à des questions de compréhension en phrases complètes.
\end{itemize}\end{multicols}}

\begin{sommaire}
\begin{multicols}{2}
\begin{enumerate}
\item Texte support : Jonas gründet eine Farm
\item Compréhension du texte
\item Lexique thématique : la ferme et l'élevage
\item Grammaire 1 : le futur I
\item Grammaire 2 : la finalité (um ... zu)
\item Commentaire modèle et production guidée
\item Activité d'intégration
\item Méthodes et exercices résolus
\item Exercices
\item Corrigés des exercices
\item Fiche bilan
\end{enumerate}
\end{multicols}
\end{sommaire}

\begin{objectifs}
\begin{itemize}
\item À la fin de cette leçon, je sais comprendre globalement puis en détail un texte allemand sur la création d'une activité agropastorale.
\item À la fin de cette leçon, je sais employer correctement le vocabulaire thématique (die Farm, das Vieh, der Mais, das Feld, anbauen, züchten) avec articles et pluriels.
\item À la fin de cette leçon, je sais former et conjuguer le futur I (werden + infinitif) pour parler de projets.
\item À la fin de cette leçon, je sais exprimer la finalité avec « um ... zu + infinitif ».
\item À la fin de cette leçon, je sais distinguer futur I et présent à valeur de futur.
\item À la fin de cette leçon, je sais répondre à des questions de compréhension en phrases complètes.
\item À la fin de cette leçon, je sais résumer un texte en quelques phrases avec mes propres mots.
\item À la fin de cette leçon, je sais rédiger un court paragraphe présentant un projet agricole ou d'élevage.
\end{itemize}
\end{objectifs}

\begin{prerequis}
\begin{itemize}
\item Le présent des verbes réguliers et irréguliers (Seconde)
\item Les articles définis et indéfinis au nominatif et à l'accusatif (Seconde)
\item Les verbes à particule séparable (anbauen, aufstehen...) vus en Seconde
\item Le verbe werden au présent (auxiliaire du passif, déjà rencontré)
\item La place de l'infinitif en fin de phrase (avec les modaux, Seconde)
\end{itemize}
\end{prerequis}

\newpage

\coursec{Texte support : Jonas gründet eine Farm}

Au Cameroun, beaucoup de jeunes diplômés retournent au village pour créer une ferme : du maïs à Bafoussam, un élevage de bovins à Ngaoundéré, des volailles à Dschang. En Allemagne aussi, des jeunes créent des exploitations agricoles modernes, souvent biologiques. Mais comment présenter son projet, parler de ses champs et de son bétail en allemand ? Comment dire ce qu'on \emph{fera} demain et \emph{pourquoi} on le fera ?

\textbf{Problématique :} comment comprendre et exploiter un texte allemand qui présente la création d'une activité agropastorale, en repérant les informations essentielles, le vocabulaire de la ferme et les deux outils grammaticaux du projet : le futur I et la construction de la finalité \all{um ... zu} ?

\courssub{Lecture globale du texte}

\begin{methode}[Lire un texte allemand en cinq étapes]
\begin{enumerate}
\item Lire le \cle{titre} et anticiper : de quoi va parler le texte ? (\all{gründen} = fonder, créer ; \all{die Farm} = la ferme.)
\item Faire une \cle{première lecture silencieuse}, sans dictionnaire, pour saisir le sens général.
\item Repérer les \cle{mots transparents}, proches du français : \all{die Farm}, \all{der Mais}, \all{der Markt}, \all{der Kredit}, \all{die Bank}, \all{die Familie}.
\item Faire une \cle{deuxième lecture} avec le glossaire, en soulignant les mots nouveaux.
\item \cle{Reformuler} le texte en français, avec ses propres mots, sans traduire mot à mot.
\end{enumerate}
\end{methode}

\begin{textebox}[Jonas gründet eine Farm]
Jonas ist 22 Jahre alt und kommt aus Bayern, im Süden von Deutschland. Er hat eine landwirtschaftliche Schule besucht und hat einen großen Traum: er will seine eigene Farm gründen.

Nächstes Jahr wird er diesen Traum verwirklichen. Zuerst wird er mit der Bank sprechen, um einen Kredit zu bekommen. Mit diesem Geld wird er ein großes Feld kaufen und einen Stall bauen. Dann wird er Mais auf seinen Feldern anbauen, denn Mais ist ein wichtiges Produkt für die Region. Danach wird er Rinder züchten: zwanzig Kühe werden auf seiner Farm leben. Jonas wird jeden Tag früh aufstehen, um das Vieh zu füttern und die Felder zu pflegen.

Schließlich wird er seine Produkte — Mais, Milch und Fleisch — auf dem Markt in der Stadt verkaufen. Er will hart arbeiten, um seine Familie zu ernähren und den Kredit zurückzuzahlen. „Die erste Ernte wird bestimmt gut“, sagt Jonas optimistisch. „Ich werde meine Farm modern und biologisch machen.“

Seine Eltern sind stolz auf ihn, und die Jugendlichen des Dorfes beobachten sein Projekt mit großem Interesse.
\end{textebox}

\textbf{Glossaire (mots nouveaux) :}

\begin{center}
\begin{tabularx}{\linewidth}{|X|X|}\hline
\rowcolor{popblueL} \textbf{Deutsch} & \textbf{Français} \\ \hline
die Farm, -en & la ferme \\ \hline
das Vieh (nur Singular) & le bétail (collectif) \\ \hline
der Mais (nur Singular) & le maïs \\ \hline
das Feld, -er & le champ \\ \hline
anbauen (baut an, baute an, hat angebaut) & cultiver \\ \hline
züchten (züchtet, züchtete, hat gezüchtet) & élever (des animaux) \\ \hline
der Kredit, -e & le crédit, le prêt \\ \hline
die Ernte, -n & la récolte \\ \hline
der Stall, Ställe & l'étable, l'écurie \\ \hline
gründen (gründet, gründete, hat gegründet) & créer, fonder \\ \hline
ernähren (ernährt, ernährte, hat ernährt) & nourrir \\ \hline
das Rind, -er / die Kuh, Kühe & le bovin / la vache \\ \hline
füttern (füttert, fütterte, hat gefüttert) & nourrir (les animaux) \\ \hline
pflegen (pflegt, pflegte, hat gepflegt) & entretenir, soigner \\ \hline
zurückzahlen (zahlt zurück, zahlte zurück, hat zurückgezahlt) & rembourser \\ \hline
stolz (auf + Akkusativ) & fier (de) \\ \hline
\end{tabularx}
\end{center}

\begin{attention}[Pièges de traduction]
Ne traduis jamais mot à mot ! Deux pièges classiques :
\begin{itemize}
\item \all{das Vieh} est \textbf{singulier} en allemand mais désigne un \textbf{collectif} : « le bétail ». On dit \all{Das Vieh ist auf der Weide.} (« Le bétail est au pâturage »), jamais de pluriel \all{die Viehe}.
\item \all{der Mais} et \all{das Geld} n'ont pas de pluriel courant, comme « le maïs » et « l'argent » en français.
\end{itemize}
\end{attention}

\courssub{Repérage des informations générales}

Après la première lecture, réponds aux cinq questions fondamentales du journaliste et du lecteur : \all{Wer ? Wo ? Was ? Wann ? Warum ?}

\begin{center}
\begin{tabularx}{\linewidth}{|X|X|X|}\hline
\rowcolor{popblueL} \textbf{Frage} & \textbf{Français} & \textbf{Réponse dans le texte} \\ \hline
\all{Wer ?} & Qui ? & Jonas, 22 ans, jeune Bavarois \\ \hline
\all{Wo ?} & Où ? & En Bavière, au sud de l'Allemagne \\ \hline
\all{Was ?} & Quoi ? & Il fonde une ferme : champs de maïs et élevage de bovins \\ \hline
\all{Wann ?} & Quand ? & L'année prochaine (\all{nächstes Jahr}) \\ \hline
\all{Warum ?} & Pourquoi ? & Pour nourrir sa famille et vendre ses produits au marché \\ \hline
\end{tabularx}
\end{center}

\begin{exoresolu}[Questions de compréhension globale]
Réponds en allemand, avec une phrase complète : 1) \all{Woher kommt Jonas ?} 2) \all{Was will Jonas gründen ?} 3) \all{Warum spricht er mit der Bank ?}
\tcblower
1) \all{Jonas kommt aus Bayern, im Süden von Deutschland.} (« Jonas vient de Bavière, au sud de l'Allemagne. »)\par
2) \all{Er will seine eigene Farm gründen.} (« Il veut créer sa propre ferme. »)\par
3) \all{Er spricht mit der Bank, um einen Kredit zu bekommen.} (« Il parle avec la banque pour obtenir un crédit. ») — Remarque la construction de finalité \all{um ... zu bekommen} : l'infinitif avec \all{zu} se place à la fin.
\end{exoresolu}

\courssub{Deuxième lecture détaillée : le futur I et la finalité}

Relis le texte et \cle{souligne} tous les verbes conjugués avec \all{werden} + infinitif : c'est le \textbf{futur I}, le temps du projet de Jonas.

\begin{exemplebox}[Le futur I dans le texte]
\all{Nächstes Jahr wird er diesen Traum verwirklichen.} « L'année prochaine, il réalisera ce rêve. »\par
\all{Zuerst wird er mit der Bank sprechen.} « D'abord, il parlera avec la banque. »\par
\all{Dann wird er Mais auf seinen Feldern anbauen.} « Ensuite, il cultivera du maïs dans ses champs. »\par
\all{Zwanzig Kühe werden auf seiner Farm leben.} « Vingt vaches vivront dans sa ferme. »\par
\all{Schließlich wird er seine Produkte auf dem Markt verkaufen.} « Finalement, il vendra ses produits au marché. »
\end{exemplebox}

\begin{propriete}[Le futur I : forme et emploi]
\textbf{Forme :} auxiliaire \all{werden} conjugué au présent (en 2\up{e} position) + \textbf{infinitif à la fin} de la phrase.\par
\begin{center}
\begin{tabular}{|l|l|}\hline
\rowcolor{popblueL} \textbf{Personne} & \textbf{Conjugaison de \all{werden}} \\ \hline
\all{ich} & \all{werde} \\ \hline
\all{du} & \all{wirst} \\ \hline
\all{er / sie / es} & \all{wird} \\ \hline
\all{wir} & \all{werden} \\ \hline
\all{ihr} & \all{werdet} \\ \hline
\all{sie / Sie} & \all{werden} \\ \hline
\end{tabular}
\end{center}
\textbf{Emploi :} exprimer un projet, une prévision, une promesse situés dans l'avenir : \all{Ich werde hart arbeiten.} « Je travaillerai dur. »\par
\textbf{Comparaison :} le français utilise un verbe conjugué au futur (« il cultivera ») ; l'allemand utilise une construction à deux éléments encadrant la phrase : \all{Er \textbf{wird} Mais \textbf{anbauen}.}\par
\textbf{Attention verbes à particule :} au futur I, le verbe à particule \textbf{ne se sépare pas}, l'infinitif reste groupé à la fin : \all{Er wird früh \textbf{aufstehen}} (et non \all{auf ... stehen}).
\end{propriete}

Repère maintenant les constructions \all{um ... zu} + infinitif : elles expriment le \textbf{but}, la \textbf{finalité} (« pour, afin de »).

\begin{exemplebox}[La finalité dans le texte]
\all{Er spricht mit der Bank, um einen Kredit zu bekommen.} « Il parle avec la banque pour obtenir un crédit. »\par
\all{Jonas wird früh aufstehen, um das Vieh zu füttern.} « Jonas se lèvera tôt pour nourrir le bétail. »\par
\all{Er will hart arbeiten, um seine Familie zu ernähren.} « Il veut travailler dur pour nourrir sa famille. »
\end{exemplebox}

\begin{propriete}[La construction de finalité \all{um ... zu}]
\textbf{Forme :} \all{um} + (compléments) + \all{zu} + infinitif, placé \textbf{en fin de proposition}. La virgule précède toujours \all{um}.\par
\textbf{Emploi :} exprimer le but d'une action, quand le sujet des deux verbes est le même.\par
\textbf{Attention aux verbes à particule :} le \all{zu} s'insère entre la particule et le verbe radical :\par
\all{aufstehen} $\rightarrow$ \all{um früh auf\textbf{zu}stehen}\par
\all{zurückzahlen} $\rightarrow$ \all{um den Kredit zurück\textbf{zu}zahlen}\par
\all{anbauen} $\rightarrow$ \all{um Mais an\textbf{zu}bauen}
\end{propriete}

\begin{center}
\begin{tabularx}{\linewidth}{|X|X|X|}\hline
\rowcolor{popblueL} \textbf{Français} & \textbf{Deutsch} & \textbf{Remarque} \\ \hline
Il travaille pour nourrir sa famille. & \all{Er arbeitet, um seine Familie zu ernähren.} & \all{zu} + infinitif à la fin \\ \hline
Il se lève tôt pour nourrir le bétail. & \all{Er steht früh auf, um das Vieh zu füttern.} & virgule avant \all{um} \\ \hline
Il parle à la banque pour rembourser le crédit. & \all{Er spricht mit der Bank, um den Kredit zurückzuzahlen.} & verbe à particule : \all{zurück\textbf{zu}zahlen} \\ \hline
\end{tabularx}
\end{center}

\courssub{Les articulations logiques du texte}

Le récit de Jonas suit une chronologie claire grâce à des \cle{connecteurs temporels}. Releve-les : ils seront indispensables pour ton résumé et tes propres productions.

\begin{center}
\begin{tabularx}{\linewidth}{|X|X|X|}\hline
\rowcolor{popblueL} \textbf{Connecteur} & \textbf{Français} & \textbf{Exemple du texte} \\ \hline
\all{zuerst} & d'abord & \all{Zuerst wird er mit der Bank sprechen.} \\ \hline
\all{dann} & puis, ensuite & \all{Dann wird er Mais anbauen.} \\ \hline
\all{danach} & après cela & \all{Danach wird er Rinder züchten.} \\ \hline
\all{schließlich} & finalement, enfin & \all{Schließlich wird er seine Produkte verkaufen.} \\ \hline
\end{tabularx}
\end{center}

\begin{attention}[Place du verbe après un connecteur]
Quand la phrase commence par \all{zuerst}, \all{dann}, \all{danach} ou \all{schließlich}, le verbe conjugué reste en \textbf{2\up{e} position} et le sujet passe après : \all{Dann \textbf{wird er} Mais anbauen} (et non \all{Dann er wird...}). C'est l'inversion, la règle d'or de la phrase allemande.
\end{attention}

\courssub{Déclinaison des noms principaux (Nominatif / Accusatif)}

Pour construire tes propres phrases, maîtrise le cas des compléments d'objet direct (accusatif) rencontrés dans le texte.

\begin{center}
\begin{tabular}{|l|l|l|l|l|}\hline
\rowcolor{popblueL} \textbf{Nom} & \textbf{Nom. Sg.} & \textbf{Akk. Sg.} & \textbf{Nom. Pl.} & \textbf{Akk. Pl.} \\ \hline
\all{der Kredit} & \all{der Kredit} & \all{den Kredit} & \all{die Kredite} & \all{die Kredite} \\ \hline
\all{das Feld} & \all{das Feld} & \all{das Feld} & \all{die Felder} & \all{die Felder} \\ \hline
\all{die Farm} & \all{die Farm} & \all{die Farm} & \all{die Farmen} & \all{die Farmen} \\ \hline
\all{das Vieh} & \all{das Vieh} & \all{das Vieh} & \textendash & \textendash \\ \hline
\all{der Mais} & \all{der Mais} & \all{den Mais} & \textendash & \textendash \\ \hline
\all{der Stall} & \all{der Stall} & \all{den Stall} & \all{die Ställe} & \all{die Ställe} \\ \hline
\end{tabular}
\end{center}

\courssub{Lecture orale : la prononciation}

Lis le texte à voix haute après le professeur. Deux sons demandent une attention particulière :

\begin{itemize}
\item Le \all{z} se prononce « \textbf{ts} » : \all{züchten} « \emph{tsu-cht'n} », \all{zwanzig} « \emph{tsvan-tsik} », \all{zuerst} « \emph{tsou-erst} ».
\item Le \all{au} se prononce « \textbf{ao} » (comme dans « \emph{caoutchouc} ») : \all{anbauen} « \emph{an-bao-n} », \all{aufstehen} « \emph{aof-chté-n} », \all{auch} « \emph{aokh} ».
\item Le \all{w} se prononce « \textbf{v} » : \all{werden} « \emph{vèr-d'n} », \all{die Landwirtschaft} « \emph{land-virt-chaft} ».
\end{itemize}

\begin{experience}[Lecture et jeu de rôle]
Par groupes de trois : un élève lit le texte à voix haute, un deuxième pose les questions \all{Wer ? Wo ? Was ? Wann ? Warum ?}, le troisième répond sans regarder le texte. Puis échangez les rôles. Terminez par un mini-jeu de rôle : l'un est Jonas, l'autre le directeur de la banque (\all{der Bankdirektor}) qui demande : \all{Was werden Sie anbauen ? Warum brauchen Sie einen Kredit ?}
\end{experience}

\courssub{Carte mentale du texte}

\begin{popfigure}
\begin{tikzpicture}[mindmap, grow cyclic, every node/.style={concept, align=center},
root concept/.append style={concept color=popblue, font=\large\bfseries},
level 1/.append style={level distance=3.2cm, sibling angle=90},
level 1 concept/.append style={concept color=poporange}]
\node[root concept] {Jonas' Farm}
  child { node[align=center]{Felder\\ Mais anbauen\\ Ernte} }
  child { node[align=center]{Vieh\\ Rinder züchten\\ Stall} }
  child { node[align=center]{Geld\\ Kredit\\ Bank} }
  child { node[align=center]{Zukunft\\ Pläne\\ Markt} };
\end{tikzpicture}
\legende{Carte mentale du texte : les quatre idées principales autour du projet de Jonas.}
\end{popfigure}

\begin{exoresolu}[Résumé guidé du texte]
Complète ce résumé avec les mots suivants : \all{Kredit — Mais — Rinder — Markt — Familie}.\par
\all{Jonas kommt aus Bayern. Nächstes Jahr wird er eine Farm gründen. Er spricht mit der Bank, um einen \dots\ zu bekommen. Dann wird er \dots\ anbauen und \dots\ züchten. Schließlich wird er seine Produkte auf dem \dots\ verkaufen, um seine \dots\ zu ernähren.}
\tcblower
\all{Jonas kommt aus Bayern. Nächstes Jahr wird er eine Farm gründen. Er spricht mit der Bank, um einen \textbf{Kredit} zu bekommen. Dann wird er \textbf{Mais} anbauen und \textbf{Rinder} züchten. Schließlich wird er seine Produkte auf dem \textbf{Markt} verkaufen, um seine \textbf{Familie} zu ernähren.}\par
Traduction : « Jonas vient de Bavière. L'année prochaine, il créera une ferme. Il parle avec la banque pour obtenir un crédit. Ensuite, il cultivera du maïs et élèvera des bovins. Finalement, il vendra ses produits au marché pour nourrir sa famille. »
\end{exoresolu}

\begin{savaistu}[L'agriculture en Allemagne et au Cameroun]
En Allemagne, l'agriculture biologique (\all{die biologische Landwirtschaft}) occupe plus de 10 \% des surfaces agricoles, et les jeunes agriculteurs comme Jonas reçoivent souvent des aides de l'État pour s'installer. Au Cameroun, le secteur agropastoral emploie la majorité de la population rurale : maïs dans l'Ouest, élevage dans l'Adamaoua, cacao dans le Centre et le Sud. Savoir présenter un projet agricole en allemand est un vrai atout : l'Allemagne coopère avec le Cameroun dans de nombreux programmes de développement rural !
\end{savaistu}

\begin{aretenir}[Ce qu'il faut retenir de cette section]
\begin{itemize}
\item Lire un texte en cinq étapes : titre, lecture globale, mots transparents, lecture détaillée, reformulation.
\item Les cinq questions de repérage : \all{Wer ? Wo ? Was ? Wann ? Warum ?}
\item Le vocabulaire de la ferme : \all{die Farm, das Vieh, der Mais, das Feld, der Stall, die Ernte, der Kredit, anbauen, züchten}.
\item Le futur I : \all{werden} conjugué + infinitif en fin de phrase (verbe à particule non séparé).
\item La finalité : \all{um ... zu} + infinitif, avec virgule avant \all{um} ; \all{zu} s'insère dans le verbe à particule (\all{aufzustehen}).
\item Les connecteurs chronologiques : \all{zuerst, dann, danach, schließlich} + inversion du sujet (verbe en 2\up{e} position).
\item Déclinaison accusatif des noms masculins : \all{der} $\rightarrow$ \all{den} (ex. \all{den Kredit, den Stall, den Mais}).
\end{itemize}
\end{aretenir}

\coursec{Compréhension du texte}

Après avoir découvert le texte \emph{Jonas gründet eine Farm} dans la section précédente, nous allons maintenant l'exploiter méthodiquement. La compréhension de l'écrit ne se limite pas à « comprendre le sens global » : elle exige une lecture active, une localisation précise des informations et une reformulation rigoureuse en phrases complètes. C'est ce que nous allons entraîner ici, en mobilisant le vocabulaire agropastoral et les structures grammaticales (futur I, \emph{um ... zu}) vues dans cette leçon.

\courssub{Le texte support : « Jonas gründet eine Farm »}

Pour travailler en autonomie, relis attentivement le texte ci-dessous. Les mots difficiles sont glossés sous le texte.

\begin{textebox}[Texte support : Jonas gründet eine Farm]
Jonas Müller kommt aus einer kleinen Stadt in der Nähe von Köln. Er ist 28 Jahre alt und hat Agrarwissenschaften studiert. Sein Traum ist es, eine eigene Farm in Kamerun zu gründen. Nächstes Jahr wird er nach Yaoundé reisen, um ein Grundstück zu kaufen. Er will Mais und Maniok anbauen und Hühner sowie Ziegen züchten.

Für sein Projekt braucht Jonas viel Geld. Er hat zwar etwas Erspartes, aber das reicht nicht für den Kauf des Landes, den Bau der Ställe und den Kauf der Maschinen. Deshalb spricht er nächste Woche mit seinem Bankberater, um einen Kredit zu beantragen. Die Bank wird einen Businessplan sehen wollen. Jonas hat ihn schon geschrieben: Er plant, die Ernte auf dem Markt in Mfoundi und in Supermärkten in Yaoundé zu verkaufen.

Jonas kennt die Risiken: Die Trockenzeit kann lang sein, und die Preise für Futtermittel steigen oft. Aber er ist optimistisch. „Die Nachfrage nach lokalen Produkten wächst“, sagt er. „Mit harter Arbeit und guter Planung wird meine Farm rentabel sein.“
\legende{Glossaire : \textbf{die Agrarwissenschaften} (pl.) : les sciences agronomiques ; \textbf{gründen} : créer, fonder ; \textbf{das Grundstück, die Grundstücke} : le terrain ; \textbf{anbauen} (séparable : er baut an) : cultiver ; \textbf{der Maniok} : le manioc ; \textbf{züchten} : élever (des animaux) ; \textbf{das Ersparte} : les économies ; \textbf{der Stall, die Ställe} : l'étable, le poulailler ; \textbf{der Businessplan, die Businesspläne} : le plan d'affaires ; \textbf{beantragen} : demander (officiellement) ; \textbf{die Trockenzeit, die Trockenzeiten} : la saison sèche ; \textbf{das Futtermittel, die Futtermittel} : l'aliment pour bétail ; \textbf{die Nachfrage} : la demande ; \textbf{rentabel} : rentable.}
\end{textebox}

\courssub{Stratégie de compréhension : Lire – Repérer – Citer – Rédiger}

Avant de répondre aux questions, il faut adopter une démarche en quatre étapes. Cette méthode vaut pour tout texte de compréhension (examen, concours, vie professionnelle).

\begin{popfigure}
\begin{tikzpicture}[
    node distance=1.2cm and 1.5cm,
    box/.style={rectangle, rounded corners, draw=popblue, thick, fill=popblueL, text width=2.8cm, align=center, font=\small},
    arrow/.style={-Stealth, thick, popblue}
]
\node[box, align=center] (lire){\textbf{1. LIRE}\\Globalement\\Quel sujet ? Quel ton ?};
\node[box, right=of lire, align=center] (reperer){\textbf{2. REPÉRER}\\Mots-clés\\Question $\rightarrow$ paragraphe};
\node[box, right=of reperer, align=center] (zitieren){\textbf{3. CITER}\\Preuve textuelle\\Copier le segment exact};
\node[box, right=of zitieren, align=center] (redigieren){\textbf{4. RÉDIGER}\\Phrase complète\\Reprise + réponse};

\draw[arrow] (lire) -- (reperer);
\draw[arrow] (reperer) -- (zitieren);
\draw[arrow] (zitieren) -- (redigieren);
\end{tikzpicture}
\legende{Organigramme de la méthode de compréhension écrite}
\end{popfigure}

\begin{methode}[Répondre à une question de compréhension en allemand]
\emph{Consigne :} la réponse doit être une phrase complète en allemand, qui reprend les mots de la question et cite l'information du texte.
\begin{enumerate}
\item \textbf{Analyser la question} : identifier le mot interrogatif (\all{Wer, Was, Wo, Wann, Warum, Wie}) et le sujet et le verbe attendus.
\item \textbf{Localiser l'information} : souligner dans le texte la phrase qui contient la réponse.
\item \textbf{Construire la réponse} : reprendre le sujet de la question (souvent \all{er/sie/es} ou le nom) + verbe conjugué (2\up{e} position) + information.
\item \textbf{Vérifier} : la réponse répond-elle exactement à la question ? L'orthographe (majuscules des noms, Umlaute) est-elle correcte ?
\end{enumerate}
\end{methode}

\begin{exemplebox}[Modèles de réponse]
\textbf{Frage :} Was wird Jonas anbauen? \\
\textbf{Repérage dans le texte :} „Er will Mais und Maniok anbauen …“ \\
\textbf{Antwort :} \all{Jonas wird Mais und Maniok anbauen.} (« Jonas cultivera du maïs et du manioc. »)

\textbf{Frage :} Warum spricht er mit der Bank? \\
\textbf{Repérage :} „Deshalb spricht er nächste Woche mit seinem Bankberater, um einen Kredit zu beantragen.“ \\
\textbf{Antwort :} \all{Er spricht mit der Bank, um einen Kredit zu beantragen.} (« Il parle à la banque pour demander un crédit. »)
\end{exemplebox}

\courssub{Compréhension globale}

Réponds en allemand, en phrases complètes, en te servant de la méthode ci-dessus.

\begin{enumerate}
\item Worum geht es in dem Text? (De quoi parle le texte ?)
\item Was ist das Ziel von Jonas? (Quel est l'objectif de Jonas ?)
\item Wo will Jonas seine Farm gründen? (Où veut-il créer sa ferme ?)
\end{enumerate}

\begin{corrige}[Exercices de compréhension globale]
\begin{enumerate}
\item \all{Der Text handelt von Jonas Müller, der eine Farm in Kamerun gründen will.} (Le texte parle de Jonas Müller, qui veut fonder une ferme au Cameroun.)
\item \all{Sein Ziel ist es, eine eigene Farm in Kamerun zu gründen, Mais und Maniok anzubauen und Tiere zu züchten.} (Son but est de fonder sa propre ferme au Cameroun, de cultiver du maïs et du manioc et d'élever des animaux.)
\item \all{Er will seine Farm in Kamerun, in der Nähe von Yaoundé, gründen.} (Il veut fonder sa ferme au Cameroun, près de Yaoundé.)
\end{enumerate}
\end{corrige}

\courssub{Compréhension détaillée}

Réponds précisément aux questions suivantes en t'appuyant \textbf{obligatoirement} sur une citation du texte (copie le segment exact entre guillemets).

\begin{enumerate}
\item Woher kommt Jonas?
\item Was wird er nächstes Jahr machen?
\item Was will er anbauen und züchten?
\item Wofür braucht er Geld? (Nenne drei Posten.)
\item Warum spricht er mit der Bank?
\item Was will er auf dem Markt verkaufen?
\end{enumerate}

\begin{exoresolu}[Question détaillée n° 5]
\textbf{Énoncé :} Warum spricht er mit der Bank? Begründe mit dem Text.
\tcblower
\textbf{Solution détaillée :}
\begin{enumerate}
\item \textbf{Mot interrogatif :} \all{Warum} → on attend une cause ou un but (« parce que… », « pour… »).
\item \textbf{Localisation :} paragraphe 2, phrase 3 : „Deshalb spricht er nächste Woche mit seinem Bankberater, \textbf{um einen Kredit zu beantragen}.“
\item \textbf{Construction de la réponse :} sujet \all{er} + verbe \all{spricht} (2\up{e} position) + complément de but \all{um ... zu} + infinitif à la fin.
\item \textbf{Réponse finale :} \all{Er spricht mit der Bank, um einen Kredit zu beantragen.}
\end{enumerate}
\end{exoresolu}

\begin{corrige}[Questions détaillées 1–4 et 6]
\begin{enumerate}
\item \all{Jonas kommt aus einer kleinen Stadt in der Nähe von Köln.} (Citation : „Jonas Müller kommt aus einer kleinen Stadt in der Nähe von Köln.“)
\item \all{Er wird nach Yaoundé reisen, um ein Grundstück zu kaufen.} (Citation : „Nächstes Jahr wird er nach Yaoundé reisen, um ein Grundstück zu kaufen.“)
\item \all{Er will Mais und Maniok anbauen und Hühner sowie Ziegen züchten.} (Citation : „Er will Mais und Maniok anbauen und Hühner sowie Ziegen züchten.“)
\item \all{Er braucht Geld für den Kauf des Landes, den Bau der Ställe und den Kauf der Maschinen.} (Citation : „… nicht für den Kauf des Landes, den Bau der Ställe und den Kauf der Maschinen.“)
\item \all{Er will die Ernte auf dem Markt in Mfoundi und in Supermärkten in Yaoundé verkaufen.} (Citation : „… die Ernte auf dem Markt in Mfoundi und in Supermärkten in Yaoundé zu verkaufen.“)
\end{enumerate}
\end{corrige}

\courssub{Vrai ou faux ? (Richtig / Falsch)}

Détermine si les affirmations sont justes ou fausses. \textbf{Justifie chaque réponse par une citation exacte du texte.}

\begin{enumerate}
\item Jonas hat bereits ein Grundstück in Yaoundé gekauft.
\item Er hat Agrarwirtschaft in Köln studiert.
\item Die Bank verlangt einen Businessplan.
\item Jonas will nur Mais anbauen.
\item Er ist pessimistisch wegen der Trockenzeit.
\end{enumerate}

\begin{corrige}[Vrai ou faux]
\begin{enumerate}
\item \textbf{Falsch.} Citation : „Nächstes Jahr \textbf{wird} er nach Yaoundé reisen, \textbf{um ein Grundstück zu kaufen}.“ (Futur : il n'a pas encore acheté.)
\item \textbf{Falsch.} Citation : „Er \textbf{hat Agrarwissenschaften studiert}.“ (Il a étudié les sciences agronomiques, pas « l'économie agricole » ; le texte ne dit pas non plus qu'il a étudié à Cologne : il vient seulement d'une ville \emph{près de} Cologne.)
\item \textbf{Richtig.} Citation : „Die Bank \textbf{wird einen Businessplan sehen wollen}.“
\item \textbf{Falsch.} Citation : „Er will \textbf{Mais und Maniok} anbauen und \textbf{Hühner sowie Ziegen} züchten.“
\item \textbf{Falsch.} Citations : „\textbf{Aber er ist optimistisch}.“ et „\textbf{Mit harter Arbeit und guter Planung wird meine Farm rentabel sein.}“
\end{enumerate}
\end{corrige}

\courssub{Questions personnelles et réaction (Produktion)}

Ces questions n'ont pas de réponse unique dans le texte. Elles demandent ton avis, en allemand, avec le vocabulaire de la leçon.

\begin{enumerate}
\item Möchtest du auch eine Farm in Kamerun gründen? Warum (nicht)?
\item Welche Tiere würdest du züchten: Hühner, Ziegen, Rinder oder Fische? Begründe.
\item Was sind deiner Meinung nach die größten Risiken für einen Landwirt in Kamerun?
\end{enumerate}

\begin{exemplebox}[Pistes de réponse (vocabulaire utile)]
\begin{enumerate}
\item \all{Ja, ich möchte eine Farm gründen, weil die Landwirtschaft in Kamerun wichtig ist.} / \all{Nein, ich möchte keine Farm gründen, weil die Arbeit sehr hart und das Risiko groß ist.}
\item \all{Ich würde Hühner züchten, weil sie wenig Platz brauchen und Eier und Fleisch schnell Geld bringen.}
\item \all{Die größten Risiken sind die lange Trockenzeit, teure Futtermittel und schlechte Straßen für den Transport.}
\end{enumerate}
\end{exemplebox}

\courssub{Technique du résumé (Zusammenfassung)}

Un résumé en allemand (niveau A2) doit reprendre les \textbf{trois idées principales} en \textbf{une phrase chacune}, sans détails superflus, avec ses propres mots (reformulation) ou des citations courtes.

\begin{methode}[Écrire un résumé en trois phrases]
\begin{enumerate}
\item \textbf{Identifier les trois macro-informations} : (1) Qui ? Quel projet ? (2) Quels moyens, quelles étapes ? (3) Quel résultat, quelle opinion ?
\item \textbf{Rédiger la phrase 1 (sujet + projet)} : \all{Jonas Müller will eine Farm in Kamerun gründen.}
\item \textbf{Rédiger la phrase 2 (moyens / problème)} : \all{Er braucht einen Kredit von der Bank und kauft Land und Tiere.}
\item \textbf{Rédiger la phrase 3 (perspective)} : \all{Trotz Risiken wie der Trockenzeit ist er optimistisch.}
\item \textbf{Relier si possible} avec des connecteurs (\all{deshalb}, \all{aber}, \all{außerdem}, \all{trotzdem}).
\end{enumerate}
\end{methode}

\begin{exoresolu}[Résumé modèle]
\textbf{Énoncé :} Fasse den Text in drei Sätzen zusammen.
\tcblower
\textbf{Solution :} \all{Jonas Müller, ein Agrarwissenschaftler, will nächstes Jahr eine Farm in Kamerun gründen. Dafür braucht er einen Bankkredit, um Land, Ställe und Tiere zu kaufen. Trotz der Risiken wie der langen Trockenzeit glaubt er an den Erfolg seiner Farm.} \\
(Jonas Müller, un agronome, veut fonder une ferme au Cameroun l'année prochaine. Pour cela, il a besoin d'un crédit bancaire pour acheter un terrain, des étables et des animaux. Malgré les risques comme la longue saison sèche, il croit au succès de sa ferme.)
\end{exoresolu}

\courssub{Entraînement : questions types d'examen}

Voici les questions classiques qui tombent souvent sur ce type de texte. Entraîne-toi à y répondre sans regarder le corrigé.

\begin{textebox}[Questions types — Entraînement]
1. Wer ist Jonas Müller? (Alter, Herkunft, Studium, Plan) \\
2. Was wird Jonas nächstes Jahr konkret tun? (Nenne drei Handlungen.) \\
3. Wofür braucht Jonas einen Kredit? (Nenne die drei Ausgaben.) \\
4. Was will er auf dem Markt in Mfoundi verkaufen? \\
5. Nenne zwei Risiken, die Jonas im Text nennt. \\
6. Welcher Satz im Text zeigt, dass Jonas gut vorbereitet ist? \\
7. Bilde einen Satz mit „um ... zu“: Jonas spricht mit der Bank. Er will einen Kredit bekommen.
\end{textebox}

\begin{corrige}[Questions types 1–7]
\begin{enumerate}
\item \all{Jonas Müller ist 28 Jahre alt, kommt aus der Nähe von Köln und hat Agrarwissenschaften studiert. Er plant, eine Farm in Kamerun zu gründen.}
\item \all{Er wird nach Yaoundé reisen, ein Grundstück kaufen und Mais und Maniok anbauen sowie Hühner und Ziegen züchten.}
\item \all{Er braucht den Kredit für den Kauf des Landes, den Bau der Ställe und den Kauf der Maschinen.}
\item \all{Er will die Ernte (Mais und Maniok) verkaufen.}
\item \all{Die lange Trockenzeit und die steigenden Preise für Futtermittel.}
\item \all{„Jonas hat ihn schon geschrieben“ — er hat schon einen Businessplan geschrieben.}
\item \all{Jonas spricht mit der Bank, um einen Kredit zu bekommen.}
\end{enumerate}
\end{corrige}

\courssub{Tableau récapitulatif : mots interrogatifs et structure de la réponse}

\begin{center}
\begin{tabularx}{\linewidth}{|l|X|X|}\hline
\rowcolor{popblueL}
\textbf{Fragewort} & \textbf{Information cherchée} & \textbf{Exemple de réponse} \\ \hline
\all{Wer} & Personne (sujet) & \all{Jonas gründet die Farm.} \\ \hline
\all{Was} & Chose / action & \all{Er baut Mais an.} \\ \hline
\all{Wo / Wohin} & Lieu / direction & \all{Er fährt nach Yaoundé.} \\ \hline
\all{Wann} & Moment & \all{Er reist nächstes Jahr.} \\ \hline
\all{Warum} & Cause / raison & \all{Er braucht Geld, weil die Farm teuer ist.} \\ \hline
\all{Wie} & Manière / moyen & \all{Er arbeitet hart und plant gut.} \\ \hline
\all{Wie viel / Wie viele} & Quantité & \all{Er braucht viel Geld und viele Tiere.} \\ \hline
\end{tabularx}
\end{center}

\begin{attention}[Pièges fréquents des francophones en compréhension écrite]
\begin{enumerate}
\item \textbf{Ne pas citer le texte} : une réponse sans citation (« Zitat: … ») est incomplète en examen. → Copie toujours le segment exact entre guillemets.
\item \textbf{Ordre des mots dans la réponse} : en allemand, le verbe conjugué est en \textbf{2\up{e} position} dans la proposition principale.
\begin{itemize}
    \item Si la question commence par \all{Warum ...?}, réponds par une principale : \all{Er spricht mit der Bank, weil er Geld braucht.} ou avec un but : \all{..., um Geld zu bekommen.}
    \item \textbf{Attention :} ne commence jamais une réponse par \all{Weil ...} seul. En français « Parce qu'il a besoin d'argent » peut tenir seul ; en allemand, \all{Weil er Geld braucht} est une subordonnée \textbf{incomplète} sans proposition principale.
\end{itemize}
\item \textbf{Confondre \all{Wo} et \all{Wohin}} : \all{Wo} (où, état, sans mouvement) → \all{in Kamerun} ; \all{Wohin} (vers où, mouvement) → \all{nach Yaoundé}.
\item \textbf{Oublier la majuscule aux noms} : \all{die farm} → \all{die Farm} ; \all{der mais} → \all{der Mais}.
\item \textbf{Verbe à particule séparable} : dans la réponse, la particule retourne à la fin : \all{Er \textbf{baut} Mais \textbf{an}.} (et non \all{Er anbaut Mais}).
\end{enumerate}
\end{attention}

\begin{savaistu}[Le Cameroun et l'agropastoral : un lien fort]
L'Allemagne coopère avec le Cameroun dans le secteur agricole, notamment à travers la \textbf{GIZ} (Deutsche Gesellschaft für Internationale Zusammenarbeit). Des projets d'appui à l'agriculture et à l'élevage forment de jeunes agriculteurs — un peu comme Jonas — aux techniques modernes, à la gestion d'une exploitation et à l'accès au crédit. Le vocabulaire de cette leçon (\all{Businessplan}, \all{Kredit}, \all{anbauen}, \all{züchten}) est donc directement utilisable dans un stage ou un projet réel au Cameroun.
\end{savaistu}

\coursec{Lexique thématique : la ferme et l'élevage}

Dans la section précédente, nous avons lu et compris le texte \all{Jonas gründet eine Farm}. Nous y avons rencontré de nombreux mots nouveaux : \all{die Farm}, \all{das Vieh}, \all{der Mais}, \all{anbauen}, \all{züchten}... Pour bien parler de la création d'une activité agropastorale — un thème très concret au Cameroun, où beaucoup de jeunes se lancent dans l'agriculture et l'élevage —, il faut d'abord maîtriser ce vocabulaire. Cette section organise ces mots en quatre champs lexicaux clairs, avec pour chaque nom son \cle{article} et son \cle{pluriel}.

\courssub{Observer d'abord : le vocabulaire en contexte}

\begin{textebox}[Ein Tag auf der Farm]
Jonas steht jeden Morgen um fünf Uhr auf. Zuerst geht er in den Stall und füttert das Vieh: die Kühe, die Schweine, die Hühner und die Ziegen. Dann melkt er die Kühe. Die Milch bringt er später auf den Markt in der Stadt.\\[4pt]
Um sieben Uhr fährt Jonas mit dem Traktor auf das Feld. Dort baut er Mais und Kartoffeln an. Sein Bruder pflanzt Bananen, und seine Schwester sät Reis. Im September ist die Ernte: Die Familie erntet das Getreide und verkauft es auf dem Markt.\\[4pt]
Jonas hat große Pläne. „Ich will mehr Rinder züchten“, sagt er. „Dafür brauche ich Geld. Die Bank gibt mir einen Kredit, und ich investiere das Geld in meine Farm. Nächstes Jahr kaufe ich einen neuen Traktor und ein größeres Feld.“\\[4pt]
Am Abend ist Jonas müde, aber glücklich. Seine Farm wächst, und die Familie hat genug zu essen.
\end{textebox}

\textbf{Glossaire}

\begin{center}
\begin{tabularx}{\linewidth}{|X|X|}
\hline
\rowcolor{popblueL} \textbf{Deutsch} & \textbf{Français} \\
\hline
aufstehen (steht auf, ist aufgestanden) & se lever \\
\hline
füttern & nourrir (les animaux) \\
\hline
melken & traire \\
\hline
die Milch (sing.) & le lait \\
\hline
säen & semer \\
\hline
wachsen (wächst, ist gewachsen) & pousser, grandir \\
\hline
genug zu essen haben & avoir assez à manger \\
\hline
\end{tabularx}
\end{center}

\textbf{Questions de compréhension rapide}
\begin{enumerate}
\item \all{Was füttert Jonas am Morgen?} (Que nourrit Jonas le matin ?)
\item \all{Was baut er auf dem Feld an?} (Que cultive-t-il au champ ?)
\item \all{Wer gibt ihm einen Kredit?} (Qui lui donne un crédit ?)
\end{enumerate}

\courssub{Champ lexical 1 : l'exploitation agricole}

\begin{definition}[Les bâtiments et les lieux]
\begin{itemize}
\item \all{die Farm, -en} : la ferme (mot moderne, d'origine anglaise)
\item \all{der Bauernhof, ¨-e} : la ferme, l'exploitation agricole (mot traditionnel)
\item \all{das Feld, -er} : le champ
\item \all{der Stall, ¨-e} : l'étable, l'écurie
\item \all{die Ernte, -n} : la récolte
\item \all{der Traktor, -en} : le tracteur
\end{itemize}
\end{definition}

\begin{exemplebox}[Exemples traduits]
\all{Jonas gründet eine Farm in der Nähe von Yaoundé.} « Jonas crée une ferme près de Yaoundé. »\\
\all{Der Traktor steht vor dem Stall.} « Le tracteur est devant l'étable. »\\
\all{Auf den Feldern wächst Mais.} « Dans les champs pousse du maïs. »
\end{exemplebox}

\courssub{Champ lexical 2 : les cultures}

\begin{definition}[Les plantes cultivées et les verbes de la culture]
Noms : \all{der Mais} (le maïs, sans pluriel), \all{der Reis} (le riz, sans pluriel), \all{die Banane, -n} (la banane), \all{die Kartoffel, -n} (la pomme de terre), \all{das Getreide} (les céréales, sans pluriel).\\
Verbes : \all{anbauen} (baut an, hat angebaut) : cultiver ; \all{ernten} : récolter ; \all{säen} : semer ; \all{pflanzen} : planter.
\end{definition}

\begin{attention}[Verbe à particule séparable : anbauen]
\all{anbauen} est un verbe à \cle{particule séparable}. À l'indicatif présent, la particule \all{an-} va à la fin de la phrase : \all{Er baut Mais an.} « Il cultive du maïs. » Ne dites jamais \all{Er anbaut Mais}. Au \all{Perfekt}, le \all{ge-} se place entre la particule et le verbe : \all{Er hat Mais angebaut.} « Il a cultivé du maïs. »
\end{attention}

\courssub{Champ lexical 3 : l'élevage}

\begin{definition}[Les animaux et les verbes de l'élevage]
Noms : \all{das Vieh} (le bétail, toujours au singulier), \all{das Rind, -er} (le bovin), \all{die Kuh, ¨-e} (la vache), \all{das Schwein, -e} (le cochon), \all{das Huhn, ¨-er} (la poule), \all{die Ziege, -n} (la chèvre).\\
Verbes : \all{züchten} : élever ; \all{füttern} : nourrir ; \all{melken} : traire.
\end{definition}

\begin{attention}[« das Vieh » n'a pas de pluriel]
\all{das Vieh} désigne le bétail en général et reste toujours au \cle{singulier}, comme « le bétail » en français. Pour parler de plusieurs bovins, on dit \all{viele Rinder} « beaucoup de bovins » — jamais \all{viele Viehe}. De même, \all{der Mais}, \all{der Reis} et \all{das Getreide} n'ont pas de pluriel : ce sont des noms de matière.
\end{attention}

\begin{exemplebox}[Exemples traduits]
\all{Sie züchtet Rinder.} « Elle élève des bovins. »\\
\all{Der Bauer füttert die Schweine und die Hühner.} « Le paysan nourrit les cochons et les poules. »\\
\all{Jeden Morgen melkt er die Kühe.} « Chaque matin, il trait les vaches. »
\end{exemplebox}

\courssub{Champ lexical 4 : l'argent et le projet}

\begin{definition}[L'aspect économique de la ferme]
Noms : \all{der Kredit, -e} (le crédit, le prêt), \all{das Geld, -er} (l'argent), \all{der Markt, ¨-e} (le marché).\\
Verbes : \all{verkaufen} (vendre), \all{kaufen} (acheter), \all{gründen} (fonder, créer), \all{investieren} (investir).
\end{definition}

\begin{exemplebox}[Exemples traduits]
\all{Wir verkaufen die Ernte auf dem Markt.} « Nous vendons la récolte au marché. »\\
\all{Die Bank gibt Jonas einen Kredit.} « La banque accorde un crédit à Jonas. »\\
\all{Er investiert das Geld in seine Farm.} « Il investit l'argent dans sa ferme. »
\end{exemplebox}

\courssub{Tableau récapitulatif du vocabulaire}

\begin{center}
\begin{tabularx}{\linewidth}{|l|l|X|}
\hline
\rowcolor{popblueL} \textbf{Deutsch} & \textbf{Pluriel} & \textbf{Français} \\
\hline
die Farm & -en & la ferme \\
\hline
der Bauernhof & ¨-e & la ferme (traditionnelle) \\
\hline
das Feld & -er & le champ \\
\hline
die Ernte & -n & la récolte \\
\hline
der Stall & ¨-e & l'étable \\
\hline
der Traktor & -en & le tracteur \\
\hline
der Mais & — & le maïs \\
\hline
der Reis & — & le riz \\
\hline
die Banane & -n & la banane \\
\hline
die Kartoffel & -n & la pomme de terre \\
\hline
das Getreide & — & les céréales \\
\hline
das Vieh & — & le bétail \\
\hline
das Rind & -er & le bovin \\
\hline
die Kuh & ¨-e & la vache \\
\hline
das Schwein & -e & le cochon \\
\hline
das Huhn & ¨-er & la poule \\
\hline
die Ziege & -n & la chèvre \\
\hline
der Kredit & -e & le crédit \\
\hline
das Geld & -er & l'argent \\
\hline
der Markt & ¨-e & le marché \\
\hline
\end{tabularx}
\end{center}

\begin{propriete}[La majuscule des noms]
En allemand, \cle{tous les noms} prennent une majuscule, où qu'ils soient dans la phrase : \all{das Feld}, \all{der Mais}, \all{die Kuh}. C'est une différence majeure avec le français, où seuls les noms propres prennent une majuscule. Un francophone qui écrit \all{das feld} ou \all{die kuh} fait une faute d'orthographe. La majuscule est d'ailleurs un excellent indice pour repérer les noms dans un texte.
\end{propriete}

\courssub{Les familles de mots}

En allemand, les mots d'une même famille se construisent de façon très régulière. Observer ces familles aide à deviner le sens des mots nouveaux.

\begin{center}
\begin{tabularx}{\linewidth}{|X|X|X|}
\hline
\rowcolor{popblueL} \textbf{Verbe} & \textbf{Nom dérivé} & \textbf{Français} \\
\hline
anbauen (cultiver) & der Anbau & la culture (des plantes) \\
\hline
züchten (élever) & die Zucht ; der Züchter & l'élevage ; l'éleveur \\
\hline
ernten (récolter) & die Ernte & la récolte \\
\hline
verkaufen (vendre) & der Verkauf & la vente \\
\hline
\end{tabularx}
\end{center}

\begin{exemplebox}[La famille en action]
\all{Der Anbau von Mais ist hier wichtig.} « La culture du maïs est importante ici. »\\
\all{Die Zucht von Rindern braucht viel Geld.} « L'élevage de bovins demande beaucoup d'argent. »\\
\all{Der Züchter verkauft seine Tiere.} « L'éleveur vend ses animaux. »
\end{exemplebox}

\courssub{Carte mentale : Auf der Farm}

\begin{popfigure}
\begin{center}
\begin{tikzpicture}[
  mindmap,
  grow cyclic,
  every node/.style={concept},
  root concept/.append style={concept color=popgreen, font=\large\bfseries},
  level 1/.append style={level distance=3.2cm, sibling angle=90},
  level 2/.append style={level distance=2.4cm, sibling angle=45},
  concept connection/.append style={color=popmuted}
]
\node[root concept]{Auf der Farm}
  child[concept color=popblue] { node {Gebäude}
    child { node {der Stall} }
    child { node {der Bauernhof} }
  }
  child[concept color=popgold] { node {Pflanzen}
    child { node {der Mais} }
    child { node {die Kartoffel} }
    child { node {das Feld} }
  }
  child[concept color=poporange] { node {Tiere}
    child { node {die Kuh} }
    child { node {das Schwein} }
    child { node {die Ziege} }
  }
  child[concept color=poppurple] { node {Geld}
    child { node {der Kredit} }
    child { node {der Markt} }
  };
\end{tikzpicture}
\end{center}
\legende{Carte mentale du vocabulaire : quatre branches pour organiser et mémoriser les mots de la ferme.}
\end{popfigure}

\courssub{Expressions utiles à retenir par cœur}

\begin{aretenir}[Les cinq expressions clés du thème]
\begin{itemize}
\item \all{eine Farm gründen} : créer une ferme
\item \all{Mais anbauen} : cultiver du maïs
\item \all{Vieh züchten} : élever du bétail
\item \all{auf dem Markt verkaufen} : vendre au marché
\item \all{einen Kredit bekommen} : obtenir un crédit
\end{itemize}
Ces expressions reviendront dans toutes vos productions écrites et orales sur ce chapitre. Notez la construction \all{auf dem Markt} : en allemand, on est « sur » le marché (\all{auf} + datif), pas « à » le marché.
\end{aretenir}

\begin{methode}[Mémoriser le vocabulaire : la méthode du triplet]
\begin{enumerate}
\item N'apprenez jamais un nom allemand seul. Apprenez toujours le \cle{triplet} : article + nom + pluriel.
\item Écrivez-le sous cette forme abrégée : \all{der Mais, —} ; \all{das Feld, -er} ; \all{die Kuh, ¨-e}. Le tiret signale la terminaison du pluriel, le tréma (¨) signale que la voyelle du radical se modifie.
\item Interrogez-vous à voix haute : « le champ ? » → \all{das Feld, -er}. Répétez chaque triplet trois fois.
\item Réutilisez le mot dans une phrase complète : \all{Das Feld ist groß.} Un mot employé dans une phrase est un mot mémorisé.
\item Regroupez les mots par famille (carte mentale ci-dessus) : le cerveau retient mieux les réseaux que les listes isolées.
\end{enumerate}
\end{methode}

\begin{savaistu}[D'où viennent ces mots ?]
\all{der Bauer} signifie « le paysan ». En ajoutant le suffixe \all{-n-} et le mot \all{der Hof} (la cour, la cour de ferme), on obtient \all{der Bauernhof} : littéralement « la cour du paysan », c'est-à-dire la ferme. Le mot \all{die Farm}, lui, est un \cle{anglicisme} très courant dans l'allemand moderne, surtout chez les jeunes et dans les médias économiques. Au Cameroun, on observe le même phénomène : beaucoup de jeunes entrepreneurs parlent de leur « farm » pour désigner leur exploitation de cacao, de maïs ou leur élevage de poules !
\end{savaistu}

\begin{exoresolu}[Compléter avec le bon mot et le bon article]
Énoncé — Complétez les phrases avec : \all{das Feld}, \all{die Ernte}, \all{das Vieh}, \all{der Kredit}, \all{der Markt}.
\begin{enumerate}
\item Jonas baut Mais auf \ldots\ an.
\item Er füttert \ldots\ im Stall.
\item Im September verkauft die Familie \ldots\ .
\item Die Bank gibt ihm \ldots\ .
\item Er verkauft die Kartoffeln auf \ldots\ .
\end{enumerate}
\tcblower
Solution détaillée :
\begin{enumerate}
\item \all{Jonas baut Mais auf dem Feld an.} — \all{auf} + datif (lieu) : \all{das Feld} devient \all{dem Feld}.
\item \all{Er füttert das Vieh im Stall.} — accusatif, article neutre inchangé ; \all{das Vieh} reste au singulier.
\item \all{Im September verkauft die Familie die Ernte.} — accusatif féminin : \all{die} ne change pas.
\item \all{Die Bank gibt ihm einen Kredit.} — accusatif masculin : \all{der Kredit} devient \all{einen Kredit}.
\item \all{Er verkauft die Kartoffeln auf dem Markt.} — \all{auf} + datif : \all{der Markt} devient \all{dem Markt}.
\end{enumerate}
Remarque : observez comme le choix du cas (accusatif ou datif) transforme l'article. C'est pourquoi il faut toujours connaître le genre du nom.
\end{exoresolu}

\begin{exercice}[À vous de jouer]
Traduisez en allemand :
\begin{enumerate}
\item Le paysan cultive du maïs et du riz.
\item Elle élève des chèvres et des poules.
\item Nous vendons la récolte au marché de Douala.
\item La banque donne un crédit à Jonas.
\item Le matin, il nourrit le bétail dans l'étable.
\end{enumerate}
\end{exercice}

\begin{corrige}[Exercice « À vous de jouer »]
\begin{enumerate}
\item \all{Der Bauer baut Mais und Reis an.} (attention : particule \all{an} en fin de phrase)
\item \all{Sie züchtet Ziegen und Hühner.} (pluriels : \all{die Ziege, -n} ; \all{das Huhn, ¨-er})
\item \all{Wir verkaufen die Ernte auf dem Markt in Douala.}
\item \all{Die Bank gibt Jonas einen Kredit.} (\all{der Kredit} → accusatif \all{einen Kredit})
\item \all{Am Morgen füttert er das Vieh im Stall.} (\all{das Vieh} reste au singulier ; \all{im} = \all{in dem})
\end{enumerate}
\end{corrige}

\begin{attention}[Erreurs fréquentes des francophones]
\begin{itemize}
\item Oublier la majuscule : écrire \all{das feld} au lieu de \all{das Feld}.
\item Traduire « élever » par \all{erhöhen} (qui signifie « augmenter ») au lieu de \all{züchten} pour les animaux.
\item Dire \all{auf der Markt} au lieu de \all{auf dem Markt} (datif masculin).
\item Mettre un pluriel à \all{das Vieh}, \all{der Mais} ou \all{der Reis}.
\item Oublier de séparer la particule : \all{Er anbaut Mais} au lieu de \all{Er baut Mais an.}
\end{itemize}
\end{attention}

Ce lexique est désormais votre boîte à outils. Dans la section suivante, nous verrons comment projeter la ferme de Jonas dans l'avenir grâce au \cle{futur I}.

\coursec{Grammaire 1 : le futur I}

Dans la section précédente, nous avons découvert le vocabulaire de la ferme et de l'élevage : \all{die Farm}, \all{das Vieh}, \all{der Mais}, \all{anbauen}, \all{züchten}. Ce lexique nous a permis de décrire ce que Jonas possède et ce qu'il fait aujourd'hui. Mais le texte \all{„Jonas gründet eine Farm“} ne parle pas seulement du présent : Jonas a des \cle{projets}. Il annonce ce qu'il fera demain, le mois prochain, l'année prochaine. Pour exprimer l'avenir, l'allemand dispose d'un temps spécifique : le \cle{futur I} (\all{das Futur I}). C'est ce temps que nous allons observer, comprendre et pratiquer dans cette section.

\courssub{Observation : comment Jonas parle de l'avenir}

Relisons quelques phrases du texte support. Observons-les attentivement avant d'énoncer la règle.

\begin{exemplebox}[Phrases tirées du texte]
\begin{itemize}
\item \all{Er wird eine Farm gründen.} « Il créera une ferme. »
\item \all{Er wird Mais anbauen.} « Il cultivera du maïs. »
\item \all{Wir werden Rinder züchten.} « Nous élèverons des bovins. »
\item \all{Nächstes Jahr wird er die erste Ernte verkaufen.} « L'année prochaine, il vendra la première récolte. »
\item \all{Er wird nicht aufgeben.} « Il n'abandonnera pas. »
\end{itemize}
\end{exemplebox}

\textbf{Questions d'observation :}
\begin{enumerate}
\item Quel verbe revient dans toutes ces phrases ? (\all{wird}, \all{werden} : c'est le verbe \all{werden}.)
\item Où se trouve ce verbe dans la phrase ? (En deuxième position, comme tout verbe conjugué allemand.)
\item Où se trouve l'autre verbe (\all{gründen}, \all{anbauen}, \all{züchten}, \all{verkaufen}, \all{aufgeben}) ? (Tout à la fin de la phrase, à l'infinitif.)
\end{enumerate}

Nous constatons donc que le futur allemand se construit avec \textbf{deux éléments} : une forme conjuguée de \all{werden} et un infinitif rejeté en fin de phrase. C'est exactement le même mécanisme que le \all{Perfekt} avec \all{haben}/\all{sein} et le participe passé final : l'allemand « encadre » la phrase entre le verbe conjugué et la forme verbale finale. On parle de \cle{cadre verbal} (\all{die Satzklammer}).

\courssub{Formation du futur I}

\begin{propriete}[Règle de formation du futur I]
Le \cle{futur I} se forme avec :
\begin{itemize}
\item l'auxiliaire \all{werden} \textbf{conjugué au présent}, placé en \textbf{deuxième position} dans la phrase déclarative (en première position dans la question sans mot interrogatif) ;
\item le verbe principal à l'\textbf{infinitif}, placé en \textbf{dernière position} de la phrase.
\end{itemize}
\all{Ich werde nächstes Jahr eine Farm gründen.} « Je créerai une ferme l'année prochaine. »
\end{propriete}

Voici la conjugaison complète de l'auxiliaire \all{werden} au présent, à connaître absolument par cœur :

\begin{center}
\begin{tabular}{|l|l|l|}
\hline
\rowcolor{popblueL} \textbf{Personne} & \textbf{werden (Präsens)} & \textbf{Exemple au futur I} \\
\hline
ich & werde & \all{Ich werde Mais anbauen.} \\
\hline
du & wirst & \all{Du wirst hart arbeiten.} \\
\hline
er / sie / es & wird & \all{Er wird eine Farm gründen.} \\
\hline
wir & werden & \all{Wir werden Rinder züchten.} \\
\hline
ihr & werdet & \all{Ihr werdet die Ernte verkaufen.} \\
\hline
sie / Sie & werden & \all{Sie werden erfolgreich sein.} \\
\hline
\end{tabular}
\end{center}

\begin{attention}[Attention à la 2e et à la 3e personne du singulier]
La forme \all{du wirst} est irrégulière : le \emph{e} du radical \all{werd-} devient \emph{i}. On écrit \all{du wirst} et non \all{du werdest}. De même, à la 3e personne : \all{er wird} (et non \all{er werdet}). Seule la forme \all{ihr werdet} garde le \emph{e}. Retiens aussi que \all{werden} est ici un \textbf{auxiliaire} : il ne porte pas le sens, c'est l'infinitif final qui donne le sens de l'action.
\end{attention}

Le schéma suivant visualise la structure de la phrase au futur I : le verbe conjugué et l'infinitif forment une sorte de « cadre » (\all{Satzklammer}) qui entoure tous les autres éléments.

\begin{popfigure}
\begin{tikzpicture}[
box/.style={draw=popblue, fill=popblueL, rounded corners=2pt, align=center, minimum height=0.9cm, font=\small},
arr/.style={-{Stealth[length=2.5mm]}, thick, popdark}]
\node[box, align=center] (s) at (0,0){\textbf{Sujet}\\ \all{Er}};
\node[box, fill=poporangeL, draw=poporange, align=center] (w) at (2.6,0){\textbf{werden}\\ \all{wird}};
\node[box, align=center] (m) at (5.6,0){\textbf{Compléments}\\ \all{nächstes Jahr}\\ \all{Mais}};
\node[box, fill=popgreenL, draw=popgreen, align=center] (i) at (8.8,0){\textbf{Infinitif}\\ \all{anbauen}};
\draw[arr] (s) -- (w);
\draw[arr] (w) -- (m);
\draw[arr] (m) -- (i);
\draw[arr, poppurple, dashed] (w.south) .. controls (4.2,-1.6) and (7,-1.6) .. (i.south);
\node[font=\small\itshape, poppurple] at (5.6,-1.9) {le cadre verbal : werden \ldots\ Infinitif};
\end{tikzpicture}
\legende{Structure de la phrase au futur I : \all{werden} en 2e position, infinitif en dernière position.}
\end{popfigure}

\courssub{Emplois du futur I}

Le futur I s'emploie dans quatre situations principales :

\begin{enumerate}
\item \textbf{Un projet ou une action à venir}, souvent annoncé par un complément de temps (\all{morgen} « demain », \all{nächstes Jahr} « l'année prochaine », \all{bald} « bientôt », \all{später} « plus tard ») : \all{Wir werden nächsten Monat Hühner kaufen.} « Nous achèterons des poules le mois prochain. »
\item \textbf{Une prédiction}, ce que l'on suppose pour l'avenir : \all{Die Ernte wird gut sein.} « La récolte sera bonne. »
\item \textbf{Une promesse} : \all{Ich werde dir helfen, das verspreche ich!} « Je t'aiderai, je te le promets ! »
\item \textbf{Une intention ferme} : \all{Ich werde nie aufgeben!} « Je n'abandonnerai jamais ! »
\end{enumerate}

\begin{exemplebox}[Exemples traduits]
\begin{itemize}
\item \all{Wir werden Rinder züchten.} « Nous élèverons des bovins. »
\item \all{Wirst du hart arbeiten?} « Travailleras-tu dur ? »
\item \all{Sie werden die Ernte verkaufen.} « Ils vendront la récolte. »
\item \all{Morgen wird es regnen.} « Demain, il pleuvra. » (prédiction)
\item \all{Ihr werdet bald ein großes Feld haben.} « Vous aurez bientôt un grand champ. »
\end{itemize}
\end{exemplebox}

\begin{popfigure}
\begin{tikzpicture}[
lbl/.style={font=\small, align=center},
pt/.style={circle, fill=popblue, inner sep=2.5pt}]
\draw[-{Stealth[length=3mm]}, very thick, popdark] (-0.5,0) -- (10.5,0);
\node[pt, fill=popmuted] (p) at (1.5,0) {};
\node[pt, fill=popgreen] (pr) at (5,0) {};
\node[pt, fill=poporange] (f) at (8.5,0) {};
\node[lbl, below=0.15cm of p, align=center]{Passé\\ \all{Er hat gearbeitet.}};
\node[lbl, below=0.15cm of pr, align=center]{Présent\\ \all{Er arbeitet.}};
\node[lbl, below=0.15cm of f, poporange, align=center]{\textbf{Futur I}\\ \all{Er wird arbeiten.}};
\node[font=\small\itshape, popmuted] at (10,0.5) {temps};
\end{tikzpicture}
\legende{Frise des temps : le futur I situe l'action après le moment présent.}
\end{popfigure}

\courssub{Futur I ou présent à valeur de futur ?}

Voici une particularité importante de l'allemand. Quand le moment de l'action est déjà indiqué par un complément de temps, les Allemands utilisent très souvent le \textbf{présent} au lieu du futur I, surtout à l'oral :

\begin{center}
\begin{tabularx}{\linewidth}{|X|X|X|}
\hline
\rowcolor{popblueL} \textbf{Français} & \textbf{Deutsch (présent + temps)} & \textbf{Deutsch (Futur I)} \\
\hline
Demain, je vais à Bamenda. & \all{Morgen fahre ich nach Bamenda.} & \all{Morgen werde ich nach Bamenda fahren.} \\
\hline
L'année prochaine, je crée une ferme. & \all{Nächstes Jahr gründe ich eine Farm.} & \all{Nächstes Jahr werde ich eine Farm gründen.} \\
\hline
Ce soir, nous mangeons du maïs grillé. & \all{Heute Abend essen wir Mais.} & \all{Heute Abend werden wir Mais essen.} \\
\hline
\end{tabularx}
\end{center}

\begin{savaistu}[Le futur... au présent !]
À l'oral, les Allemands préfèrent souvent le présent accompagné d'un complément de temps : \all{„Nächstes Jahr gründe ich eine Farm“} est plus naturel et plus fréquent que \all{„Nächstes Jahr werde ich eine Farm gründen“}. Le futur I sert alors surtout à insister sur une promesse, une prédiction ou une intention forte. En français, nous faisons pareil avec « demain, je pars » au lieu de « demain, je partirai ».
\end{savaistu}

\courssub{Cas particuliers : verbes à particule et négation}

\textbf{1. Les verbes à particule séparable.} Au futur I, l'infinitif reste \textbf{entier} en fin de phrase : la particule ne se sépare pas. C'est une différence capitale avec le présent.

\begin{center}
\begin{tabularx}{\linewidth}{|X|X|}
\hline
\rowcolor{popblueL} \textbf{Présent (particule séparée)} & \textbf{Futur I (infinitif entier en fin)} \\
\hline
\all{Er baut Mais an.} & \all{Er wird Mais anbauen.} \\
\hline
\all{Er gibt nicht auf.} & \all{Er wird nicht aufgeben.} \\
\hline
\all{Der Markt fängt um 7 Uhr an.} & \all{Der Markt wird um 7 Uhr anfangen.} \\
\hline
\end{tabularx}
\end{center}

\textbf{2. La négation.} On place \all{nicht} avant l'élément nié (souvent juste avant l'infinitif) et \all{kein-} devant un nom :

\begin{itemize}
\item \all{Er wird nicht aufgeben.} « Il n'abandonnera pas. »
\item \all{Er wird kein Geld verlieren.} « Il ne perdra pas d'argent. »
\item \all{Wir werden das Feld nicht verkaufen.} « Nous ne vendrons pas le champ. »
\end{itemize}

\begin{attention}[Erreur fréquente n° 1 : la place de l'infinitif]
En français, on dit « je vais \emph{créer} une ferme » : l'infinitif suit immédiatement l'auxiliaire. En allemand, c'est impossible ! Tout le reste de la phrase passe \textbf{avant} l'infinitif final. On dit \all{Ich werde nächstes Jahr eine Farm gründen} et jamais \all{Ich werde gründen eine Farm}. Même dans une phrase longue, l'infinitif attend patiemment la toute dernière place.
\end{attention}

\begin{attention}[Erreur fréquente n° 2 : les faux amis werden et bekommen]
Ne confonds pas \all{werden} avec \all{bekommen}. Employé seul (sans infinitif), \all{werden} signifie « devenir » : \all{Er wird Bauer.} « Il devient paysan. » Le verbe \all{bekommen} signifie « obtenir, recevoir » : \all{Er bekommt Geld.} « Il reçoit de l'argent. » Et surtout, \all{bekommen} ne veut jamais dire « devenir » ! Pour le futur, retiens : \all{werden} + infinitif = futur ; \all{werden} + nom/adjectif = devenir.
\end{attention}

\courssub{Texte de consolidation : les projets de la famille Ngo Bell}

\begin{textebox}[„Die Pläne der Familie Ngo Bell“]
Die Familie Ngo Bell wohnt in der Nähe von Bafoussam. Sie hat ein kleines Feld und ein paar Hühner, aber sie hat große Pläne. „Nächstes Jahr werden wir mehr Land kaufen“, sagt der Vater. „Wir werden Mais und Bohnen anbauen, und wir werden auch Ziegen züchten.“ Die Mutter ergänzt: „Ich werde ein kleines Geschäft am Markt eröffnen. Dort werde ich Eier, Gemüse und Mais verkaufen.“ Die Tochter Amina ist siebzehn Jahre alt. Sie sagt: „Ich werde die Schule beenden, und dann werde ich Landwirtschaft studieren. Später werde ich meinen Eltern helfen.“ Der Sohn Junior lacht: „Und ich werde der beste Fußballspieler von Bafoussam werden! Aber am Wochenende werde ich auf dem Feld arbeiten, das verspreche ich.“ Die Großmutter ist skeptisch: „Werdet ihr wirklich so hart arbeiten?“ „Ja, Großmutter“, antwortet die Familie, „wir werden nicht aufgeben! Wir werden kein Geld verlieren, wir werden erfolgreich sein.“ Am Abend sitzt die Familie zusammen und plant die Zukunft. Bald wird hier eine große Farm stehen.
\end{textebox}

\textbf{Glossaire :} \all{der Plan, die Pläne} « le projet » ; \all{in der Nähe von} « près de » ; \all{das Land} « ici : la terre, le terrain » ; \all{die Bohne, -n} « le haricot » ; \all{die Ziege, -n} « la chèvre » ; \all{ergänzen} « ajouter » ; \all{das Geschäft, -e} « le commerce, la boutique » ; \all{eröffnen} « ouvrir » ; \all{beenden} « terminer » ; \all{die Landwirtschaft} « l'agriculture » ; \all{skeptisch} « sceptique » ; \all{erfolgreich} « qui réussit, prospère » ; \all{die Zukunft} « l'avenir ».

\textbf{Questions de compréhension :}
\begin{enumerate}
\item \all{Was wird die Familie nächstes Jahr kaufen?}
\item \all{Was wird die Mutter am Markt verkaufen?}
\item \all{Was wird Amina nach der Schule studieren?}
\item \all{Wann wird Junior auf dem Feld arbeiten?}
\item Relevez dans le texte trois phrases au futur I et traduisez-les.
\end{enumerate}

\begin{corrige}[Questions de compréhension]
\begin{enumerate}
\item \all{Nächstes Jahr wird die Familie mehr Land kaufen.} « L'année prochaine, la famille achètera plus de terrain. »
\item \all{Sie wird Eier, Gemüse und Mais verkaufen.} « Elle vendra des œufs, des légumes et du maïs. »
\item \all{Nach der Schule wird sie Landwirtschaft studieren.} « Après l'école, elle étudiera l'agriculture. »
\item \all{Am Wochenende wird er auf dem Feld arbeiten.} « Le week-end, il travaillera au champ. »
\item Exemples : \all{„Wir werden Mais und Bohnen anbauen“} « Nous cultiverons du maïs et des haricots » ; \all{„Wir werden nicht aufgeben“} « Nous n'abandonnerons pas » ; \all{„Bald wird hier eine große Farm stehen“} « Bientôt, il y aura ici une grande ferme ».
\end{enumerate}
\end{corrige}

\courssub{Entraînement}

\begin{exoresolu}[Mets les phrases au futur I]
Énoncé : transforme les phrases suivantes au futur I. a) \all{Er gründet eine Farm.} b) \all{Wir bauen Mais an.} c) \all{Ich verkaufe die Ernte.} d) \all{Gibst du auf?}
\tcblower
Solution détaillée : on conjugue \all{werden} à la personne du sujet (2e position) et on renvoie le verbe principal à l'infinitif en fin de phrase. Pour les verbes à particule, l'infinitif reste entier.
\begin{itemize}
\item a) \all{Er wird eine Farm gründen.} « Il créera une ferme. »
\item b) \all{Wir werden Mais anbauen.} « Nous cultiverons du maïs. » (particule \all{an-} rattachée : \all{anbauen})
\item c) \all{Ich werde die Ernte verkaufen.} « Je vendrai la récolte. »
\item d) \all{Wirst du aufgeben?} « Abandonneras-tu ? » (à l'interrogative, \all{werden} passe en première position, l'infinitif reste en dernier)
\end{itemize}
\end{exoresolu}

\begin{exercice}[À ton tour]
\begin{enumerate}
\item Conjugue \all{werden} aux personnes demandées : du \ldots , er \ldots , wir \ldots , ihr \ldots , sie (pluriel) \ldots
\item Traduis en allemand : a) « Demain, j'achèterai des poules. » b) « Nous n'abandonnerons pas. » c) « Elle ne perdra pas d'argent. » d) « Travaillerez-vous (\all{Sie}) sur le champ ? »
\item Réponds à la question par une phrase complète au futur I : \all{Was wirst du nach dem Abitur machen?}
\end{enumerate}
\end{exercice}

\begin{corrige}[Exercice « À ton tour »]
\begin{enumerate}
\item \all{du wirst}, \all{er wird}, \all{wir werden}, \all{ihr werdet}, \all{sie werden}.
\item a) \all{Morgen werde ich Hühner kaufen.} b) \all{Wir werden nicht aufgeben.} c) \all{Sie wird kein Geld verlieren.} d) \all{Werden Sie auf dem Feld arbeiten?}
\item Exemple de réponse : \all{Nach dem Abitur werde ich Landwirtschaft studieren.} « Après le baccalauréat, j'étudierai l'agriculture. »
\end{enumerate}
\end{corrige}

\begin{aretenir}[L'essentiel du futur I]
\begin{itemize}
\item Formation : \all{werden} conjugué (2e position) + infinitif (dernière position) : \all{Ich werde arbeiten.}
\item Conjugaison : \all{ich werde, du wirst, er wird, wir werden, ihr werdet, sie/Sie werden}.
\item Emplois : projet, prédiction, promesse, intention.
\item Verbes à particule : l'infinitif reste entier en fin : \all{Er wird Mais anbauen.}
\item Négation : \all{nicht} avant l'élément nié, \all{kein-} devant un nom.
\item À l'oral, le présent + complément de temps remplace souvent le futur I : \all{Morgen fahre ich nach Bamenda.}
\end{itemize}
\end{aretenir}

\coursec{Grammaire 2 : la finalité (um ... zu)}

Dans la section précédente, nous avons appris à parler des projets de Jonas grâce au \cle{futur I} : \all{Er wird eine Farm gründen.} Mais un projet n'a de sens que s'il a un \textbf{but}. Pourquoi Jonas parle-t-il à la banque ? Pourquoi travaille-t-il si dur ? Pour répondre à ces questions, l'allemand dispose d'une construction très élégante et très fréquente : la proposition finale avec \cle{um ... zu}. C'est l'objet de cette section.

\courssub{Observation : repérons la construction dans le texte}

Relisons attentivement deux phrases du texte support \emph{Jonas gründet eine Farm} :

\begin{exemplebox}[Phrases du texte]
\all{Um einen Kredit zu bekommen, spricht er mit der Bank.}\\
« Pour obtenir un crédit, il parle à la banque. »

\all{Er arbeitet jeden Tag hart, um seine Familie zu ernähren.}\\
« Il travaille dur chaque jour pour nourrir sa famille. »
\end{exemplebox}

Observons et comparons avec le français :

\begin{itemize}
\item En français, le but s'exprime simplement avec « \textbf{pour} + infinitif » : \emph{pour obtenir}, \emph{pour nourrir}. Un seul petit mot suffit.
\item En allemand, la proposition de but est \textbf{encadrée} : elle s'ouvre avec \all{um} et se ferme avec \all{zu} + infinitif, placé \textbf{à la toute fin} de la proposition.
\item Entre \all{um} et \all{zu} + infinitif, on place tous les compléments : \all{um \textbf{einen Kredit} zu bekommen}, \all{um \textbf{seine Familie} zu ernähren}.
\item Le sujet du verbe principal et le sujet sous-entendu de l'infinitif sont \textbf{les mêmes} : c'est \emph{er} (Jonas) qui parle à la banque et c'est \emph{er} qui obtiendra le crédit.
\end{itemize}

\begin{definition}[La proposition finale avec \all{um ... zu}]
La construction \cle{um ... zu} + infinitif exprime le \textbf{but}, la \textbf{finalité} d'une action : elle répond à la question \all{Wozu?} / \all{Warum?} (« pour quoi faire ? », « dans quel but ? »). Elle correspond au français « \textbf{pour} + infinitif » ou « \textbf{afin de} + infinitif ».
\end{definition}

\courssub{La règle : forme et construction}

\begin{propriete}[\all{um ... zu} + infinitif = pour + infinitif]
La proposition finale est un \textbf{cadre} :
\begin{itemize}
\item \all{um} \textbf{ouvre} la proposition (première place) ;
\item les \textbf{compléments} (objets, lieux, temps) se placent au milieu ;
\item \all{zu} + \textbf{infinitif} \textbf{ferme} la proposition (dernière place). Le \all{zu} est collé devant l'infinitif.
\end{itemize}
Exemple : \all{Er arbeitet hart, \textbf{um} seine Familie \textbf{zu ernähren}.} = « Il travaille dur pour nourrir sa famille. »
\end{propriete}

Le schéma suivant visualise cette structure en cadre, qui rappelle la place du verbe conjugué dans la phrase allemande : l'allemand aime « encadrer » ses énoncés.

\begin{popfigure}
\begin{center}
\begin{tikzpicture}[
  node distance=0.3cm and 0.2cm,
  box/.style={draw=popblue, fill=popblueL, rounded corners=2pt, align=center, font=\small, inner sep=4pt},
  mainbox/.style={draw=poporange, fill=poporangeL, rounded corners=2pt, align=center, font=\small, inner sep=4pt},
  lbl/.style={font=\footnotesize\itshape, text=popmuted}
]
\node[box] (um) {\all{um}};
\node[box, right=of um, align=center] (compl){compléments\\\all{einen Kredit}};
\node[box, right=of compl, align=center] (zu){\all{zu} + Infinitif\\\all{bekommen}};
\node[mainbox, right=of zu, align=center] (verb){Verbe conjugué\\\all{spricht}};
\node[mainbox, right=of verb, align=center] (rest){Sujet + suite\\\all{er mit der Bank}};

\draw[-{Stealth}, thick, popblue] (um) -- (compl);
\draw[-{Stealth}, thick, popblue] (compl) -- (zu);
\draw[thick, popdark] (zu) -- (verb);
\draw[-{Stealth}, thick, poporange] (verb) -- (rest);

\node[lbl, below=0.5cm of compl.south] {proposition finale (le cadre \all{um ... zu})};
\node[lbl, below=0.5cm of verb.south] {proposition principale};

\draw[decorate, decoration={brace, amplitude=4pt, mirror}, popblue] (um.south west) -- (zu.south east);
\draw[decorate, decoration={brace, amplitude=4pt, mirror}, poporange] (verb.south west) -- (rest.south east);
\end{tikzpicture}
\end{center}
\legende{Structure en cadre : \all{um} + compléments + \all{zu} + infinitif, puis le verbe conjugué de la proposition principale.}
\end{popfigure}

\courssub{La place de la proposition finale}

La proposition avec \all{um ... zu} peut occuper deux positions dans la phrase :

\begin{enumerate}
\item \textbf{Après la proposition principale} (position la plus fréquente) : on la sépare par une virgule.\\
\all{Er spricht mit der Bank, um einen Kredit zu bekommen.}\\
« Il parle à la banque pour obtenir un crédit. »
\item \textbf{Au début de la phrase} (pour insister sur le but) : la proposition finale occupe alors la \textbf{première place}, et le verbe conjugué de la proposition principale vient immédiatement après, en \textbf{deuxième position}.\\
\all{Um einen Kredit zu bekommen, spricht er mit der Bank.}\\
« Pour obtenir un crédit, il parle à la banque. »
\end{enumerate}

\begin{attention}[Le verbe conjugué reste en deuxième position !]
Quand la proposition finale est placée au début, elle compte comme \textbf{une seule place} (la première). Le verbe conjugué de la proposition principale vient donc juste après la virgule, et le sujet se place après le verbe :\\
\all{Um Geld zu verdienen, \textbf{verkauft} \textbf{er} Mais.} — et jamais \all{Um Geld zu verdienen, er verkauft Mais.}
\end{attention}

\courssub{Le cas particulier des verbes à particule séparable}

Que se passe-t-il quand l'infinitif est un verbe à particule séparable comme \all{anbauen} (cultiver), \all{aufstehen} (se lever), \all{einkaufen} (faire les courses) ? Le \all{zu} ne se place pas devant le verbe entier : il \textbf{s'insère entre la particule et le verbe}, et le tout s'écrit en \textbf{un seul mot}.

\begin{propriete}[Verbes séparables avec \all{zu}]
Particule + \all{zu} + verbe, en un seul mot :\\
\all{anbauen} $\rightarrow$ \all{an\textbf{zu}bauen} ; \all{aufstehen} $\rightarrow$ \all{auf\textbf{zu}stehen} ; \all{einkaufen} $\rightarrow$ \all{ein\textbf{zu}kaufen}.\\
Exemple : \all{Jonas braucht Geld, um Mais \textbf{anzubauen}.} = « Jonas a besoin d'argent pour cultiver du maïs. »
\end{propriete}

\begin{attention}[Jamais \all{zu anbauen} !]
L'erreur classique du francophone est de coller le \all{zu} devant le verbe séparable. Retenez bien : le \all{zu} se place \textbf{entre} la particule et le verbe : \all{anzubauen}, \all{aufzustehen}, \all{einzukaufen} — jamais \all{zu anbauen}, jamais \all{zu aufstehen}.\\
\all{Er steht früh auf, um die Tiere zu füttern.} devient avec un verbe séparable à l'infinitif : \all{Er geht früh ins Bett, um früh \textbf{aufzustehen}.} (« Il se couche tôt pour se lever tôt. »)
\end{attention}

\courssub{Exemples nombreux et variés}

\begin{exemplebox}[La finalité dans la vie de la ferme]
\all{Sie lernt Deutsch, um in Deutschland zu studieren.}\\
« Elle apprend l'allemand pour étudier en Allemagne. »

\all{Um Geld zu verdienen, verkauft er Mais auf dem Markt.}\\
« Pour gagner de l'argent, il vend du maïs au marché. »

\all{Wir fahren nach Douala, um neue Maschinen zu kaufen.}\\
« Nous allons à Douala pour acheter de nouvelles machines. »

\all{Die Bauern düngen das Feld, um mehr Gemüse anzubauen.}\\
« Les paysans fertilisent le champ pour cultiver plus de légumes. »

\all{Um das Vieh zu füttern, steht Jonas jeden Morgen um fünf Uhr auf.}\\
« Pour nourrir le bétail, Jonas se lève chaque matin à cinq heures. »

\all{Ich spare Geld, um nächstes Jahr ein Stück Land zu kaufen.}\\
« J'économise de l'argent pour acheter un terrain l'année prochaine. »
\end{exemplebox}

\courssub{Comparaison français / allemand}

\begin{center}
\begin{tabularx}{\linewidth}{|X|X|X|}
\hline
\rowcolor{popblueL} Français & Deutsch & Remarque \\
\hline
pour obtenir un crédit & \all{um einen Kredit zu bekommen} & un mot en français, un cadre en allemand \\
\hline
pour cultiver du maïs & \all{um Mais anzubauen} & verbe séparable : \all{zu} au milieu \\
\hline
pour se lever tôt & \all{um früh aufzustehen} & idem : \all{aufzu}stehen \\
\hline
pour la famille (nom) & \all{für die Familie} & devant un nom : \all{für}, pas \all{um ... zu} \\
\hline
afin que son fils apprenne & \all{damit sein Sohn lernt} & sujets différents : \all{damit} \\
\hline
\end{tabularx}
\end{center}

\begin{attention}[Ne traduis jamais « pour » par \all{für} devant un infinitif !]
\all{für} ne s'emploie \textbf{que devant un nom} : \all{Er arbeitet für die Familie.} (« Il travaille pour la famille. ») Devant un verbe à l'infinitif, « pour » se traduit obligatoirement par \all{um ... zu} :\\
\all{Er arbeitet, um die Familie zu ernähren.} (« Il travaille pour nourrir la famille. »)\\
Faux : \all{Er arbeitet für die Familie zu ernähren.}
\end{attention}

\courssub{La condition d'emploi : même sujet, sinon \all{damit}}

\begin{propriete}[\all{um ... zu} ou \all{damit} ?]
On n'emploie \all{um ... zu} que si le \textbf{sujet des deux verbes est identique}. Si les sujets sont différents, on utilise \all{damit} + une proposition subordonnée complète (verbe conjugué à la fin) :
\begin{itemize}
\item Même sujet : \all{Jonas lernt viel, um die Prüfung zu bestehen.} (« Jonas apprend beaucoup pour réussir l'examen » — c'est Jonas qui apprend et qui réussira.)
\item Sujets différents : \all{Jonas lernt viel, damit sein Sohn später die Farm leiten kann.} (« Jonas apprend beaucoup pour que son fils puisse plus tard diriger la ferme » — Jonas apprend, mais c'est le fils qui dirigera.)
\end{itemize}
\end{propriete}

\begin{savaistu}[Un exercice classique du Bac]
Au baccalauréat, la transformation \all{damit} $\leftrightarrow$ \all{um ... zu} est un exercice de grammaire très fréquent. La clé est toujours la même : regarde les sujets ! S'ils sont identiques, \all{damit} devient \all{um ... zu} : \all{Er spart Geld, damit er Land kauft.} $\rightarrow$ \all{Er spart Geld, um Land zu kaufen.} Si les sujets diffèrent, la transformation est impossible et \all{damit} doit rester.
\end{savaistu}

\courssub{Méthode : construire une phrase de but}

\begin{methode}[Transformer deux informations en une phrase avec \all{um ... zu}]
Énoncé de départ : « Il parle à la banque. But : obtenir un crédit. »
\begin{enumerate}
\item Identifier l'\textbf{action} : il parle à la banque $\rightarrow$ \all{Er spricht mit der Bank.}
\item Identifier le \textbf{but} et le mettre à l'infinitif : obtenir un crédit $\rightarrow$ \all{einen Kredit bekommen}.
\item Vérifier que le \textbf{sujet est le même} (ici : \emph{er} dans les deux cas) $\rightarrow$ \all{um ... zu} possible.
\item Construire le cadre : \all{um} + compléments + \all{zu} + infinitif $\rightarrow$ \all{um einen Kredit zu bekommen}.
\item Assembler, avec une virgule : \all{Er spricht mit der Bank, um einen Kredit zu bekommen.}
\item Variante avec insistance : placer le but en tête, verbe conjugué en deuxième position : \all{Um einen Kredit zu bekommen, spricht er mit der Bank.}
\end{enumerate}
\end{methode}

\courssub{Texte d'application : les projets de Jonas}

\begin{textebox}[Jonas plant die Zukunft seiner Farm]
Jonas steht jeden Morgen um fünf Uhr auf, um sein Vieh zu füttern. Er hat große Pläne: Er wird hart arbeiten, um seine Farm zu vergrößern. Zuerst spricht er mit der Bank, um einen Kredit zu bekommen. Mit diesem Geld wird er Maschinen kaufen, um mehr Mais anzubauen. Er wird auch Hühner und Ziegen züchten, um mehr Geld zu verdienen. Um frisches Wasser für die Tiere zu haben, wird er einen Brunnen bauen. Seine Frau verkauft Gemüse auf dem Markt in Yaoundé, um die Familie zu unterstützen. Jonas spart jeden Monat einen Teil des Geldes, um nächstes Jahr ein größeres Feld zu kaufen. Er schickt seine Kinder zur Schule, damit sie später die Farm leiten können. Um alles gut zu organisieren, schreibt er jeden Abend seine Ausgaben in ein Heft. Die Nachbarn kommen oft, um von Jonas zu lernen. Er erklärt ihnen alles, damit auch sie erfolgreich sein können. Jonas ist stolz: Er hat die Farm gegründet, um seiner Familie ein besseres Leben zu geben.

\textbf{Glossaire} : das Vieh (le bétail) ; füttern (nourrir, donner à manger) ; vergrößern (agrandir) ; der Kredit, -e (le crédit, le prêt) ; die Ziege, -n (la chèvre) ; der Brunnen, - (le puits) ; unterstützen (soutenir) ; die Ausgabe, -n (la dépense) ; das Heft, -e (le cahier) ; leiten (diriger) ; erfolgreich (qui réussit) ; stolz (fier).
\end{textebox}

\textbf{Questions de compréhension} (réponds en allemand, avec une phrase complète) :
\begin{enumerate}
\item \all{Warum steht Jonas jeden Morgen um fünf Uhr auf?}
\item \all{Warum spricht er mit der Bank?}
\item \all{Was wird er kaufen, um mehr Mais anzubauen?}
\item \all{Warum schickt Jonas seine Kinder zur Schule?} (attention : sujets différents !)
\end{enumerate}

\courssub{Exercice résolu}

\begin{exoresolu}[Construis des phrases de but]
Relie chaque action à son but avec \all{um ... zu} (ou \all{damit} si nécessaire).\\
a) Jonas travaille hart. But : nourrir sa famille.\\
b) Il va au marché. But : vendre du maïs.\\
c) Elle se couche tôt. But : se lever tôt.\\
d) Jonas économise de l'argent. But : son fils pourra étudier à l'université.
\tcblower
\textbf{Solution détaillée :}\\
a) Même sujet (Jonas) $\rightarrow$ \all{um ... zu}. \all{Jonas arbeitet hart, um seine Familie zu ernähren.} « Jonas travaille dur pour nourrir sa famille. »\\
b) Même sujet $\rightarrow$ \all{Er geht auf den Markt, um Mais zu verkaufen.} « Il va au marché pour vendre du maïs. »\\
c) Verbe séparable \all{aufstehen} : le \all{zu} s'insère au milieu $\rightarrow$ \all{aufzustehen}. \all{Sie geht früh ins Bett, um früh aufzustehen.} « Elle se couche tôt pour se lever tôt. »\\
d) Sujets différents (Jonas économise, mais c'est le fils qui étudiera) $\rightarrow$ \all{damit} obligatoire, verbe conjugué à la fin : \all{Jonas spart Geld, damit sein Sohn an der Universität studieren kann.} « Jonas économise de l'argent pour que son fils puisse étudier à l'université. »
\end{exoresolu}

\begin{exercice}[À toi de jouer]
Traduis en allemand : 1. Il apprend l'allemand pour travailler en Allemagne. 2. Pour acheter un champ, elle économise de l'argent. 3. Nous nous levons tôt pour nourrir le bétail (attention au verbe séparable !). 4. Jonas explique tout à ses voisins pour qu'ils réussissent aussi.
\end{exercice}

\begin{corrige}[Exercice : la finalité]
1. \all{Er lernt Deutsch, um in Deutschland zu arbeiten.}\\
2. \all{Um ein Feld zu kaufen, spart sie Geld.} (proposition finale en tête, verbe conjugué \all{spart} en deuxième position)\\
3. \all{Wir stehen früh auf, um das Vieh zu füttern.} — ici le verbe séparable \all{aufstehen} est le verbe conjugué de la principale ; l'infinitif \all{füttern} n'est pas séparable, donc \all{zu füttern}.\\
4. Sujets différents $\rightarrow$ \all{damit} : \all{Jonas erklärt seinen Nachbarn alles, damit sie auch erfolgreich sein können.}
\end{corrige}

\begin{aretenir}[L'essentiel de la finalité]
\begin{itemize}
\item \all{um ... zu} + infinitif = « pour + infinitif » : \all{um} ouvre, \all{zu} + infinitif ferme la proposition.
\item Verbes séparables : le \all{zu} s'insère \textbf{entre} particule et verbe : \all{anzubauen}, \all{aufzustehen}.
\item Proposition finale en tête $\rightarrow$ verbe conjugué en deuxième position : \all{Um Geld zu verdienen, verkauft er Mais.}
\item Jamais \all{für} devant un infinitif ; \all{für} seulement devant un nom.
\item Même sujet $\rightarrow$ \all{um ... zu} ; sujets différents $\rightarrow$ \all{damit} + verbe conjugué final.
\end{itemize}
\end{aretenir}

\coursec{Commentaire modèle et production guidée}

Dans la section précédente, tu as appris à exprimer la finalité avec la tournure \all{um \ldots\ zu} : \all{Jonas spricht mit der Bank, um einen Kredit zu bekommen.} « Jonas parle à la banque pour obtenir un crédit. » Tu maîtrises maintenant tous les outils de la leçon : le vocabulaire de la ferme et de l'élevage, le futur I et la construction de la finalité. Il est temps de les réunir dans deux tâches complètes, typiques de l'examen : le \cle{commentaire de texte} et la \cle{production guidée}. C'est exactement ce que l'on te demandera au baccalauréat : lire, comprendre, commenter, puis produire à ton tour.

\courssub{La méthode du commentaire de texte en quatre étapes}

Commenter un texte, ce n'est ni le traduire mot à mot ni le recopier. C'est montrer que tu l'as compris et que tu sais en parler avec tes propres mots. En allemand, on appelle cela \all{einen Text kommentieren}. La démarche se fait toujours en quatre étapes, dans un ordre précis.

\begin{methode}[Rédiger un commentaire de texte en quatre étapes]
\begin{enumerate}
\item \cle{Présenter} le texte : qui en est le sujet ? De quoi parle-t-on ? Commence toujours par la formule : \all{Der Text handelt von\ldots} « Le texte traite de\ldots ». Tu peux ajouter l'origine du texte si elle est connue : \all{Der Text stammt aus einer Zeitung.} « Le texte provient d'un journal. »
\item \cle{Résumer} le contenu avec tes propres phrases, en suivant l'ordre du texte. Utilise les connecteurs chronologiques : \all{Zuerst\ldots} « d'abord\ldots », \all{dann\ldots} « ensuite\ldots », \all{danach\ldots} « après cela\ldots », \all{schließlich\ldots} « finalement\ldots ».
\item \cle{Analyser} : relève le vocabulaire du thème (ici : la ferme, l'élevage), les temps et constructions grammaticales remarquables (ici : le futur I, \all{um \ldots\ zu}), et l'intention de l'auteur (\all{Die Absicht des Autors ist\ldots} « L'intention de l'auteur est\ldots »).
\item \cle{Juger} : donne ton avis personnel, toujours en dernière position, avec la formule : \all{Meiner Meinung nach\ldots} « À mon avis\ldots ». Justifie ton avis avec \all{weil} « parce que ».
\end{enumerate}
\end{methode}

\begin{popfigure}
\begin{center}
\begin{tikzpicture}[
  node distance=0.55cm,
  etape/.style={rectangle, rounded corners=4pt, draw=popblue, fill=popblueL, thick, align=center, text width=4.6cm, minimum height=1.1cm, font=\small},
  fleche/.style={-{Stealth[length=3mm]}, thick, popdark}
]
\node[etape, align=center] (e1){\textbf{1. Présentation}\\ \all{Der Text handelt von\ldots}};
\node[etape, below=of e1, align=center] (e2){\textbf{2. Résumé}\\ \all{Zuerst\ldots, dann\ldots, schließlich\ldots}};
\node[etape, below=of e2, align=center] (e3){\textbf{3. Analyse}\\ Lexique, grammaire, intention};
\node[etape, below=of e3, draw=poporange, fill=poporangeL, align=center] (e4){\textbf{4. Avis personnel}\\ \all{Meiner Meinung nach\ldots}};
\draw[fleche] (e1) -- (e2);
\draw[fleche] (e2) -- (e3);
\draw[fleche] (e3) -- (e4);
\end{tikzpicture}
\end{center}
\legende{Organigramme du commentaire de texte : quatre étapes dans un ordre obligatoire.}
\end{popfigure}

\begin{aretenir}[Les formules-clés du commentaire]
\begin{center}
\begin{tabularx}{\linewidth}{|X|X|}
\hline
\rowcolor{popblueL} \textbf{Deutsch} & \textbf{Français} \\
\hline
\all{Der Text handelt von\ldots} & Le texte traite de\ldots \\
\hline
\all{Zuerst\ldots, dann\ldots, schließlich\ldots} & D'abord\ldots, ensuite\ldots, finalement\ldots \\
\hline
\all{Der Autor benutzt das Futur I.} & L'auteur utilise le futur I. \\
\hline
\all{Die Absicht des Textes ist\ldots} & L'intention du texte est\ldots \\
\hline
\all{Meiner Meinung nach\ldots} & À mon avis\ldots \\
\hline
\all{Ich finde das Projekt interessant, weil\ldots} & Je trouve le projet intéressant, parce que\ldots \\
\hline
\end{tabularx}
\end{center}
\end{aretenir}

\begin{propriete}[Rappel : le futur I (Futur I)]
Le futur I s forme avec l'auxiliaire \all{werden} au présent + l'infinitif du verbe principal \textbf{rejeté en fin de phrase}.
\begin{center}
\begin{tabular}{|l|l|l|}
\hline
\rowcolor{popblueL} \textbf{Personne} & \textbf{werden (Präsens)} & \textbf{Exemple : gründen} \\
\hline
ich & werde & \all{ich werde gründen} \\
\hline
du & wirst & \all{du wirst gründen} \\
\hline
er / sie / es & wird & \all{er wird gründen} \\
\hline
wir & werden & \all{wir werden gründen} \\
\hline
ihr & werdet & \all{ihr werdet gründen} \\
\hline
sie / Sie & werden & \all{sie werden gründen} \\
\hline
\end{tabular}
\end{center}
\emph{Emploi :} projets, intentions, prédictions pour l'avenir. \emph{Marqueurs fréquents :} \all{nächstes Jahr}, \all{in Zukunft}, \all{bald}, \all{später}.
\end{propriete}

\courssub{Commentaire modèle rédigé}

Voici un commentaire complet du texte de la leçon, \all{Jonas gründet eine Farm}. Lis-le attentivement et repère les quatre étapes de la méthode.

\begin{textebox}[Commentaire modèle : Jonas gründet eine Farm]
Der Text handelt von Jonas, einem jungen Bauern aus Bayern. Er wird nächstes Jahr eine Farm gründen. Zuerst wird er Mais anbauen, dann wird er Rinder züchten. Um einen Kredit zu bekommen, spricht er mit der Bank. Der Text benutzt oft das Futur I, weil Jonas über seine Pläne spricht. Man findet auch viele Wörter zum Thema Landwirtschaft: das Feld, das Vieh, der Mais. Meiner Meinung nach ist sein Projekt sehr interessant, weil die Landwirtschaft wichtig für die Zukunft ist.
\end{textebox}

\textbf{Glossaire :} \all{der Bauer, -n} « le paysan, l'agriculteur » ; \all{gründen} « fonder, créer » ; \all{der Kredit, -e} « le crédit » ; \all{der Plan, „-e} « le projet, le plan » ; \all{die Landwirtschaft} « l'agriculture » ; \all{die Zukunft} « l'avenir » ; \all{das Vieh} « le bétail » (collectif, pas de pluriel) ; \all{das Feld, „-er} « le champ » ; \all{züchten} « élever ».

Analysons ensemble la structure de ce modèle : la première phrase présente le texte (étape 1). Les trois phrases suivantes résument le contenu avec \all{zuerst}, \all{dann} et la construction de finalité \all{um \ldots\ zu} (étape 2). Les deux phrases suivantes analysent la grammaire (futur I) et le lexique (étape 3). La dernière phrase donne l'avis personnel, justifié par \all{weil} (étape 4). Remarque que le commentaire est rédigé au présent pour parler du texte, et au futur I seulement pour rapporter les projets de Jonas.

\begin{attention}[L'inversion après « Meiner Meinung nach »]
Quand une phrase allemande commence par un élément autre que le sujet, le verbe conjugué reste en deuxième position et le sujet vient après : c'est l'\cle{inversion}. On dit donc : \all{Meiner Meinung nach \textbf{ist} sein Projekt interessant.} « À mon avis, son projet est intéressant. » Ne dis jamais \all{Meiner Meinung nach sein Projekt ist\ldots}, calque fautif du français « à mon avis son projet est\ldots ». Même règle après \all{zuerst} et \all{dann} : \all{Zuerst \textbf{wird} er Mais anbauen.} « D'abord, il cultivera du maïs. »
\end{attention}

\courssub{Production guidée : « Mein Projekt »}

À ton tour maintenant. Tu vas présenter ton propre projet agricole ou d'élevage au Cameroun, en suivant une trame en cinq points. Cette trame reprend exactement les acquis de la leçon : le vocabulaire de la ferme, le futur I et la finalité.

\begin{methode}[La trame de production « Mein Projekt »]
\begin{enumerate}
\item \cle{Qui je suis} : \all{Ich heiße\ldots\ und ich komme aus\ldots} « Je m'appelle\ldots\ et je viens de\ldots ».
\item \cle{Ce que je ferai} (futur I) : \all{Nächstes Jahr werde ich eine Farm gründen.} « L'année prochaine, je créerai une ferme. »
\item \cle{Ce que je cultiverai ou élèverai} (vocabulaire du thème) : \all{Ich werde Mais anbauen und Hühner züchten.} « Je cultiverai du maïs et j'élèverai des poules. »
\item \cle{Pourquoi} (finalité avec \all{um \ldots\ zu}) : \all{Ich werde Mais anbauen, um Geld zu verdienen.} « Je cultiverai du maïs pour gagner de l'argent. »
\item \cle{Conclusion} : \all{Meiner Meinung nach ist die Landwirtschaft die Zukunft Kameruns.} « À mon avis, l'agriculture est l'avenir du Cameroun. »
\end{enumerate}
\end{methode}

\begin{exemplebox}[Phrases-modèles pour ta production]
\all{Ich werde Hühner züchten.} « J'élèverai des poules. »

\all{Ich werde Mais anbauen, um Geld zu verdienen.} « Je cultiverai du maïs pour gagner de l'argent. »

\all{Ich werde Schweine züchten, um meine Familie zu ernähren.} « J'élèverai des porcs pour nourrir ma famille. »

\all{Ich werde Kakao anbauen, um ihn nach Europa zu verkaufen.} « Je cultiverai du cacao pour le vendre en Europe. »

\all{Um einen Kredit zu bekommen, werde ich mit der Bank sprechen.} « Pour obtenir un crédit, je parlerai avec la banque. »
\end{exemplebox}

\begin{attention}[Verbes à particule séparable dans « um … zu »]
Avec un verbe à particule séparable (ex. \all{anbauen}, \all{aufbauen}, \all{mitmachen}), la particule se colle devant l'infinitif après \all{zu} : \all{zu} + \all{an} + \all{bauen} = \all{anzubauen}.
Exemple : \all{Ich will Mais anbauen.} $\rightarrow$ \all{Um Mais \textbf{anzubauen}, brauche ich Land.} « Pour cultiver du maïs, j'ai besoin de terre. »
Ne sépare jamais la particule dans une proposition \all{um \ldots\ zu} !
\end{attention}

\begin{center}
\begin{tabularx}{\linewidth}{|X|X|X|}
\hline
\rowcolor{popblueL} \textbf{Français} & \textbf{Deutsch} & \textbf{Remarque} \\
\hline
Je créerai une ferme. & \all{Ich werde eine Farm gründen.} & futur I : \all{werden} + infinitif en fin \\
\hline
J'élèverai des poules. & \all{Ich werde Hühner züchten.} & \all{das Huhn, „-er} au pluriel \\
\hline
Je cultiverai du maïs. & \all{Ich werde Mais anbauen.} & particule \all{an-} en fin de phrase \\
\hline
pour gagner de l'argent & \all{um Geld zu verdienen} & finalité : infinitif avec \all{zu} en fin \\
\hline
À mon avis, c'est important. & \all{Meiner Meinung nach ist das wichtig.} & inversion après la formule d'avis \\
\hline
\end{tabularx}
\end{center}

\begin{exoresolu}[Rédiger le début d'une production « Mein Projekt »]
Énoncé : tu es Aïcha, élève à Bafoussam. Rédige quatre phrases de production en suivant la trame : présentation, projet au futur I, culture et élevage, finalité avec \all{um \ldots\ zu}.
\tcblower
Solution détaillée :

\all{Ich heiße Aïcha und ich komme aus Bafoussam.} « Je m'appelle Aïcha et je viens de Bafoussam. » (étape 1 : présentation au présent)

\all{Nächstes Jahr werde ich eine kleine Farm gründen.} « L'année prochaine, je créerai une petite ferme. » (étape 2 : futur I, \all{werde} en deuxième position, infinitif \all{gründen} à la fin)

\all{Ich werde Mais und Bohnen anbauen und ich werde auch Hühner züchten.} « Je cultiverai du maïs et des haricots et j'élèverai aussi des poules. » (étape 3 : vocabulaire du thème ; attention, la particule \all{an-} du verbe \all{anbauen} va en fin de proposition)

\all{Ich werde Hühner züchten, um Eier auf dem Markt zu verkaufen.} « J'élèverai des poules pour vendre des œufs au marché. » (étape 4 : finalité ; \all{um} en tête de la subordonnée, infinitif \all{zu verkaufen} à la toute fin ; \all{das Ei, -er} « l'œuf », \all{der Markt, „-e} « le marché »)
\end{exoresolu}

\courssub{Production orale et correction collective}

La dernière étape est orale : tu présentes ton projet en une minute devant la classe. Parle lentement, regarde ton public, et n'écris pas un texte complet à lire : note seulement cinq mots-clés, un par étape de la trame.

\begin{experience}[Présentation orale : « Mein Projekt » en une minute]
Par groupes de quatre, chaque élève présente son projet en une minute. Les autres écoutent et remplissent la grille d'évaluation ci-dessous. Ensuite, deux ou trois volontaires présentent devant toute la classe, et le groupe classe corrige collectivement : on relève d'abord les réussites, puis une ou deux erreurs à corriger ensemble au tableau. Chronomètre en main, le professeur arrête la présentation à soixante secondes : apprends à aller à l'essentiel.
\end{experience}

\begin{corrige}[Grille de critères pour la correction collective]
\begin{center}
\begin{tabularx}{\linewidth}{|X|c|X|}
\hline
\rowcolor{popgreenL} \textbf{Critère} & \textbf{Points} & \textbf{Ce qu'on attend} \\
\hline
Vocabulaire du thème & 4 & au moins quatre mots de la leçon (\all{die Farm, das Vieh, der Mais, anbauen, züchten}) \\
\hline
Futur I correct & 4 & \all{werden} conjugué en deuxième position, infinitif à la fin \\
\hline
Finalité \all{um \ldots\ zu} & 4 & au moins une construction correcte, infinitif avec \all{zu} en fin \\
\hline
Clarté et prononciation & 4 & débit calme, phrases complètes, durée respectée \\
\hline
Structure en cinq étapes & 4 & trame respectée, conclusion avec \all{Meiner Meinung nach} \\
\hline
\end{tabularx}
\end{center}
Total : vingt points. Une erreur de place du verbe coûte un point ; un mot français glissé dans la phrase coûte un point.
\end{corrige}

\begin{exercice}[À toi de jouer]
Rédige ta production complète « Mein Projekt » en huit à dix lignes, en suivant la trame en cinq étapes. Utilise au moins trois verbes au futur I, deux constructions avec \all{um \ldots\ zu} et cinq mots du vocabulaire de la ferme. Puis entraîne-toi à la présenter oralement en une minute, sans lire.
\end{exercice}

\begin{savaistu}[L'agriculture, un pont entre le Cameroun et l'Allemagne]
Au Cameroun, le ministère de l'Agriculture et du Développement rural (MINADER) encourage les jeunes à créer des exploitations agricoles et pastorales. De son côté, la coopération allemande, notamment l'agence GIZ, très présente à Yaoundé, finance des projets de formation agricole et des partenariats avec de jeunes entrepreneurs camerounais. Savoir présenter son projet en allemand — exactement l'exercice que tu viens de faire — peut donc un jour t'ouvrir les portes d'une bourse de formation ou d'un partenariat international. La leçon de grammaire d'aujourd'hui est peut-être le début d'un vrai dossier de candidature.
\end{savaistu}

\newpage

\coursec{Activité d'intégration}

\coursec{Leçon ALA08 : La création d'une activité agropastorale}

\courssub{Objectifs de l'activité}

Cette activité d'intégration clôt la leçon ALA08 en faisant travailler les élèves comme de jeunes entrepreneurs. À l'issue de la séance, chaque élève sera capable de :

\begin{itemize}
\item réinvestir le \cle{vocabulaire} de la ferme et de l'élevage (\all{die Farm, das Vieh, der Mais, das Feld, anbauen, züchten, die Ernte, der Kredit}) dans un contexte de communication authentique ;
\item employer correctement le \cle{futur I} (\all{werden + Infinitiv}) pour présenter un projet et des prévisions ;
\item exprimer la \cle{finalité} avec la construction \all{um ... zu + Infinitiv} ;
\item rédiger en groupe un court texte cohérent de 8 à 10 lignes (\all{„Unser Projekt“}) ;
\item présenter oralement ce projet devant la classe, comme devant le jury d'une banque, avec une prononciation soignée ;
\item écouter les projets des autres groupes, poser des questions simples et voter de manière argumentée.
\end{itemize}

\courssub{Situation}

La banque régionale \all{„AgroBank Kamerun“} organise chaque année un concours intitulé \all{„Junge Unternehmer“} (« Jeunes Entrepreneurs »). Le concours s'adresse aux jeunes de 16 à 22 ans qui veulent créer une activité agropastorale dans leur région : élevage de poules, culture de maïs, élevage de bovins, maraîchage, etc. Le projet gagnant reçoit un \all{Kredit} (un crédit, un prêt) symbolique pour démarrer l'activité.

Votre classe de Première A du lycée de Yaoundé participe au concours. Par groupes de trois ou quatre, vous allez imaginer un projet agropastoral pour votre région, le rédiger par écrit, puis le présenter oralement devant la classe. Le professeur joue le rôle du directeur de la banque, et la classe entière forme le jury.

Voici l'affiche officielle du concours, affichée au tableau :

\begin{textebox}[Affiche du concours]
\textbf{AgroBank Kamerun — Wettbewerb „Junge Unternehmer“}\par
Du bist jung? Du hast eine Idee?\par
Gründe eine Farm in deiner Region!\par
\medskip
Hühner züchten in Yaoundé? Mais anbauen in Bafoussam? Rinder züchten in Ngaoundéré?\par
\medskip
Schreibe dein Projekt (8–10 Zeilen) und präsentiere es vor der Jury!\par
Das beste Projekt bekommt einen Kredit von der AgroBank!\par
\medskip
\emph{Viel Glück!}
\end{textebox}

\begin{definition}[Glossaire de l'affiche]
\begin{itemize}
\item \all{der Wettbewerb, die Wettbewerbe} : le concours
\item \all{die Idee, die Ideen} : l'idée
\item \all{gründen} : fonder, créer
\item \all{das Huhn, die Hühner} : la poule, les poules
\item \all{das Rind, die Rinder} : le bovin
\item \all{die Zeile, die Zeilen} : la ligne
\item \all{die Jury, die Jurys} : le jury
\item \all{bekommen} : recevoir, obtenir
\item \all{Viel Glück!} : Bonne chance !
\end{itemize}
\end{definition}

\courssub{Consignes}

\begin{enumerate}
\item \textbf{Comprendre l'affiche (5 min).} Lisez l'affiche du concours en silence, puis à voix haute. Répondez oralement aux questions : \all{Wer organisiert den Wettbewerb? Wer kann mitmachen? Was bekommt das beste Projekt?} (Qui organise le concours ? Qui peut participer ? Que reçoit le meilleur projet ?)

\item \textbf{Choisir le projet (5 min).} En groupe de trois ou quatre, choisissez une activité agropastorale et une région : par exemple l'élevage de poules à Yaoundé, la culture de maïs à Bafoussam, l'élevage de bovins à Ngaoundéré, la culture de tomates à Douala. Donnez un nom à votre projet : \all{„Unser Projekt heißt ...“}.

\item \textbf{Préparer le lexique (10 min).} Dans votre cahier, listez au moins \textbf{six mots} du vocabulaire de la leçon que vous utiliserez, avec leur article et leur pluriel : par exemple \all{die Farm, -en ; das Feld, -er ; das Vieh (kein Plural) ; der Mais (kein Plural) ; die Ernte, -n ; der Kredit, -e}. Vérifiez les articles dans le tableau de lexique ci-dessous.

\item \textbf{Rédiger le texte (20 min).} Rédigez ensemble un texte de 8 à 10 lignes intitulé \all{„Unser Projekt“}. Le texte doit contenir obligatoirement :
\begin{itemize}
\item au moins \textbf{6 mots} du lexique thématique (noms ou verbes) ;
\item au moins \textbf{3 verbes au futur I} (\all{werden + Infinitiv} à la fin de la phrase) ;
\item au moins \textbf{2 constructions de finalité} \all{um ... zu + Infinitiv} ;
\item une phrase de conclusion qui convainc le jury : \all{„Unser Projekt ist gut für unsere Region.“}
\end{itemize}

\item \textbf{Corriger en groupe (10 min).} Relisez votre texte : verbe conjugué en deuxième position, infinitif à la fin, majuscule à tous les noms, article correct devant chaque nom. Le professeur circule et aide.

\item \textbf{Présenter oralement (15 min).} Chaque groupe présente son projet devant la classe en 2 minutes maximum. Chaque membre du groupe lit une partie du texte. Parlez lentement, regardez le jury, prononcez bien le \all{w} allemand (comme « v » français) et le \all{v} allemand (comme « f » français).

\item \textbf{Écouter et voter (10 min).} Pendant les présentations, les autres groupes écoutent et notent : un mot de vocabulaire entendu, un verbe au futur I, une construction \all{um ... zu}. Ensuite, posez une question simple au groupe : \all{„Wie viele Hühner werdet ihr kaufen?“} (Combien de poules allez-vous acheter ?) À la fin, la classe vote pour le meilleur projet, et le directeur de la banque (le professeur) accorde le \all{Kredit} symbolique au groupe gagnant.
\end{enumerate}

\courssub{Ressources à mobiliser}

\begin{aretenir}[Le futur I (Futur I)]
\textbf{Formation :} \all{werden} (conjugué, en deuxième position) + \textbf{infinitif à la fin de la phrase}.\par
\all{Ich werde eine Farm gründen.} « Je vais créer une ferme / Je créerai une ferme. »\par
\all{Wir werden Mais anbauen.} « Nous allons cultiver du maïs / Nous cultiverons du maïs. »\par
\textbf{Conjugaison de \all{werden} (verbe fort) :}
\begin{center}
\begin{tabular}{|l|l|l|l|}\hline
\rowcolor{popblueL} Person & Singular & & Plural \\ \hline
1. & \all{ich werde} & & \all{wir werden} \\ \hline
2. & \all{du wirst} & & \all{ihr werdet} \\ \hline
3. & \all{er/sie/es wird} & & \all{sie/Sie werden} \\ \hline
\end{tabular}
\end{center}
\emph{Emploi :} projet, intention, prévision pour le futur proche ou lointain (équivalent français : futur simple ou « aller + infinitif »).
\end{aretenir}

\begin{aretenir}[La finalité : \all{um ... zu + Infinitiv}]
\textbf{Structure :} \all{um} (en début de proposition subordonnée) + \textbf{zu + infinitif} à la toute fin de la phrase.\par
\all{Wir brauchen einen Kredit, um Hühner zu kaufen.} « Nous avons besoin d'un crédit pour acheter des poules. »\par
\emph{Règle de l'infinitif avec \all{zu} :} l'infinitif (éventuellement précédé de ses compléments) se place entre \all{um} et \all{zu}.\par
\emph{Verbes à particule séparable :} \all{zu} s'insère entre la particule et le radical : \all{um Mais an|zu|bauen} « pour cultiver du maïs », \all{um Hühner an|zu|schaffen} « pour se procurer des poules ».
\end{aretenir}

\begin{center}
\begin{tabularx}{\linewidth}{|X|X|X|}\hline
\rowcolor{popblueL} Français & Deutsch & Remarque \\ \hline
la ferme (exploitation) & die Farm, -en & anglicisme, synonyme : \all{der Bauernhof, -höfe} \\ \hline
le champ & das Feld, -er & neutre \\ \hline
le bétail & das Vieh & collectif, pas de pluriel \\ \hline
le maïs & der Mais & pas de pluriel \\ \hline
cultiver & anbauen & particule séparable (\all{er baut an}) \\ \hline
élever & züchten & régulier \\ \hline
la récolte & die Ernte, -n & \\ \hline
le crédit, le prêt & der Kredit, -e & \\ \hline
vendre & verkaufen & régulier \\ \hline
avoir besoin de & brauchen & + accusatif \\ \hline
nourrir (animaux) & füttern & régulier \\ \hline
le marché & der Markt, -\"e & \\ \hline
l'œuf & das Ei, -er & neutre \\ \hline
la viande & das Fleisch & pas de pluriel \\ \hline
\end{tabularx}
\end{center}

\begin{attention}[Pièges à éviter pour les francophones]
\begin{itemize}
\item \textbf{Ordre des mots (Futur I) :} Ne dites pas \all{*ich werde gründen eine Farm}. L'infinitif va \textbf{à la fin de la phrase} : \all{Ich werde eine Farm gründen.}
\item \textbf{Ordre des mots (Finalité) :} Ne dites pas \all{*um zu kaufen Hühner}. Tout le groupe nominal se place entre \all{um} et \all{zu + Infinitiv} : \all{um Hühner zu kaufen}.
\item \textbf{Majuscules :} N'oubliez jamais la majuscule aux noms : \all{die Farm, das Feld, der Mais, das Vieh, der Kredit}.
\item \textbf{Faux ami / Anglicisme :} \all{die Farm} vient de l'anglais \emph{farm}. Il désigne une exploitation agricole moderne, souvent grande. Le terme traditionnel est \all{der Bauernhof}. Ne le confondez pas avec « ferme » au sens industriel (usine).
\item \textbf{Verbe \all{brauchen} :} Il régit l'accusatif (\all{brauchen einen Kredit}), pas de préposition.
\item \textbf{Prononciation :} \all{w} = [v] (comme \emph{v}oyage) ; \all{v} = [f] (comme \emph{f}erme) ; \all{z} = [ts] (comme \emph{pizza}).
\end{itemize}
\end{attention}

\courssub{Production guidée : modèle et corrigé commenté}

\begin{exoresolu}[Unser Projekt : „Hühnerfarm Mfoundi“]
Énoncé : rédigez un texte modèle de 8 à 10 lignes pour le concours, sur le thème de l'élevage de poules à Yaoundé, avec 6 mots du lexique, 3 verbes au futur I et 2 constructions \all{um ... zu}.
\tcblower
\textbf{Production modèle :}\par
\medskip
\begin{textebox}[Unser Projekt : „Hühnerfarm Mfoundi“]
Unser Projekt heißt „Hühnerfarm Mfoundi“. Wir sind vier Schüler aus Yaoundé und wir werden eine kleine Farm gründen. Zuerst werden wir ein Feld hinter der Schule kaufen. Dann brauchen wir einen Kredit, um zweihundert Hühner zu kaufen. Wir werden die Hühner jeden Tag füttern, um gesunde Tiere zu haben. Wir werden die Eier auf dem Markt in Mokolo verkaufen. Mit dem Geld werden wir mehr Vieh kaufen. Unser Projekt ist gut für unsere Region, denn die Familien in Yaoundé brauchen Eier und Fleisch. Bitte, geben Sie uns den Kredit!
\end{textebox}
\medskip
\textbf{Traduction :}\par
« Notre projet s'appelle “Ferme de poules Mfoundi”. Nous sommes quatre élèves de Yaoundé et nous allons créer une petite ferme. D'abord, nous allons acheter un champ derrière l'école. Ensuite, nous avons besoin d'un crédit pour acheter deux cents poules. Nous nourrirons les poules chaque jour pour avoir des animaux en bonne santé. Nous vendrons les œufs au marché de Mokolo. Avec l'argent, nous achèterons plus de bétail. Notre projet est bon pour notre région, car les familles de Yaoundé ont besoin d'œufs et de viande. S'il vous plaît, donnez-nous le crédit ! »\par
\medskip
\textbf{Commentaire des choix de langue :}\par
\begin{itemize}
\item \textbf{Lexique (7 occurrences)} : \all{die Farm, das Feld, der Kredit, die Hühner (Pl. de das Huhn), füttern, der Markt, das Vieh} — tous employés avec l'article correct ou la forme conjuguée adéquate.
\item \textbf{Futur I (5 occurrences)} : \all{wir werden ... gründen / kaufen / füttern / verkaufen / kaufen} — \all{werden} en deuxième position, infinitif toujours en position finale.
\item \textbf{Finalité (2 constructions)} : \all{um zweihundert Hühner zu kaufen} et \all{um gesunde Tiere zu haben} — les compléments (\all{zweihundert Hühner}, \all{gesunde Tiere}) sont correctement placés entre \all{um} et \all{zu + Infinitiv}.
\item \textbf{Connecteurs logiques} : \all{zuerst} (d'abord), \all{dann} (ensuite), \all{denn} (car) assurent la cohérence textuelle et l'argumentation.
\item \textbf{Conclusion persuasive} : la dernière phrase \all{„Bitte, geben Sie uns den Kredit!“} s'adresse directement au jury avec le \all{Sie} de politesse (impératif formel) — exemple authentique d'interaction orale.
\end{itemize}
\end{exoresolu}

\courssub{Bilan de l'activité}

À l'issue de cette activité, faites le point collectivement avec la classe :

\begin{itemize}
\item \textbf{Sur le lexique} : chaque groupe a manipulé au moins six mots du thème de la ferme avec leurs articles ; le vote final a permis de réentendre ces mots dans plusieurs contextes différents (élevage, culture, vente).
\item \textbf{Sur la grammaire} : le futur I a été utilisé pour exprimer des projets et des prévisions réels ; la construction \all{um ... zu} a permis de justifier chaque étape du projet (pourquoi acheter, pourquoi nourrir, pourquoi vendre).
\item \textbf{Sur la communication} : les élèves ont rédigé un texte court mais complet, l'ont présenté oralement devant un « jury », ont écouté les autres et ont argumenté leur vote — quatre compétences du programme réunies en une seule tâche (compréhension, expression écrite, expression orale, interaction).
\item \textbf{Sur la culture économique} : l'activité a montré qu'on peut parler d'agriculture et d'entrepreneuriat en allemand, dans un contexte camerounais concret (Yaoundé, Bafoussam, Ngaoundéré, marché Mokolo), ce qui donne du sens à l'apprentissage de la langue.
\end{itemize}

\begin{experience}[Pour aller plus loin : Devoir individuel]
En devoir à la maison, chaque élève rédige individuellement une version améliorée du projet de son groupe (10 à 12 lignes), en ajoutant \textbf{au moins une phrase supplémentaire au futur I} et \textbf{une phrase supplémentaire avec \all{um ... zu}}. Les meilleurs textes seront affichés au mur de la salle d'allemand sous le titre \all{„Junge Unternehmer der Klasse 1A“}.
\end{experience}

\newpage

\coursec{Méthodes et exercices résolus}

\coursec{La création d'une activité agropastorale}

\courssub{Texte support : Jonas gründet eine Farm}

\begin{textebox}[Texte original — Contexte camerounais]
Jonas ist 24 Jahre alt und kommt aus Yaoundé. Er hat einen Traum: er will eine eigene Farm gründen. Nächstes Jahr wird er ein großes Feld in der Nähe von Bafoussam kaufen. Auf diesem Feld wird er Mais und Gemüse anbauen. Er wird auch Hühner und Ziegen züchten. Zurzeit lernt Jonas viel über Landwirtschaft. Er liest Bücher und schaut Videos im Internet, um moderne Methoden kennenzulernen. Seine Familie unterstützt ihn. Sein Vater wird ihm Geld für das Feld geben. Später will Jonas einen Stall bauen und einen Traktor kaufen. Er arbeitet hart, um seinen Traum zu verwirklichen. Die Farm soll seine Familie ernähren und Einkommen bringen.
\end{textebox}

\begin{definition}[Lexique du texte — Mots nouveaux avec article et pluriel]
\begin{center}
\begin{tabularx}{\linewidth}{|X|X|X|}
\hline
\rowcolor{popblueL} \textbf{Deutsch (Artikel, Plural)} & \textbf{Français} & \textbf{Exemple / Famille} \\
\hline
die Farm, -en & la ferme (agricole) & eine Farm gründen \\
das Feld, die Felder & le champ & auf dem Feld, ein großes Feld kaufen \\
das Vieh (kein Plural) & le bétail, le cheptel & Vieh züchten, das Vieh füttern \\
der Mais (kein Plural) & le maïs & Mais anbauen, der Mais wächst \\
das Gemüse (kein Plural) & les légumes & Gemüse anbauen, frisches Gemüse \\
das Huhn, die Hühner & la poule / le poulet & Hühner züchten, ein Huhn legen \\
die Ziege, -n & la chèvre & Ziegen züchten, Ziegenmilch \\
die Landwirtschaft & l'agriculture & über Landwirtschaft lernen \\
der Landwirt, -e / die Landwirtin, -nen & l'agriculteur / l'agricultrice & Landwirt werden \\
der Stall, die Ställe & l'étable, le hangar & einen Stall bauen \\
der Traktor, -en & le tracteur & einen Traktor kaufen \\
die Ernte, -n & la récolte & die Ernte verbessern, die Ernte einfahren \\
gründen & créer, fonder & eine Farm gründen, ein Unternehmen gründen \\
anbauen (trennbar) & cultiver (plantes) & Mais anbauen, der Anbau \\
züchten & élever (animaux) & Tiere züchten, die Zucht \\
unterstützen & soutenir & die Familie unterstützen \\
verwirklichen & réaliser (rêve, projet) & einen Traum verwirklichen \\
ernähren & nourrir & die Familie ernähren \\
\hline
\end{tabularx}
\end{center}
\end{definition}

\courssub{Savoir-faire 1 : Comprendre un texte (global puis détaillé)}

\begin{methode}[Lire et comprendre un texte sur la création d'une ferme]
\begin{enumerate}
\item \cle{Lire le titre et le premier paragraphe} : ils annoncent le thème (ici : \all{Jonas gründet eine Farm} — qui ? quoi ? où ?).
\item \cle{Première lecture globale} : repérer les personnages, le lieu, le moment, le projet. Ne pas s'arrêter sur chaque mot inconnu.
\item \cle{Deuxième lecture détaillée} : souligner les mots du champ lexical de la ferme (\all{das Feld, das Vieh, der Mais, anbauen, züchten}) et les marqueurs de temps et de but (\all{nächstes Jahr, später, um ... zu, will}).
\item \cle{Deviner les mots inconnus} par le contexte et les familles de mots : \all{bauen} → \all{anbauen} → \all{der Anbau} ; \all{wirklich} → \all{verwirklichen}.
\item \cle{Répondre aux questions} en reprenant les mots du texte, avec des phrases complètes : sujet + verbe en 2\textsuperscript{e} position.
\item \cle{Vérifier} : chaque réponse contient un verbe conjugué, une majuscule aux noms, la bonne place du verbe (infinitif à la fin après modal ou \all{werden}).
\end{enumerate}
\end{methode}

\begin{exoresolu}[Compréhension de texte — Type Bac]
\textbf{Énoncé.} Lisez le texte \all{Jonas gründet eine Farm} et répondez en allemand par des phrases complètes.
\begin{enumerate}
\item Wo will Jonas seine Farm gründen?
\item Was will er anbauen und was will er züchten?
\item Warum lernt er viel über Landwirtschaft?
\item Was wird sein Vater tun?
\item Was ist das Ziel der Farm? (Zwei Aspekte)
\end{enumerate}
\tcblower
\textbf{Solution détaillée.}
\begin{enumerate}
\item La question porte sur le \cle{lieu} (\all{wo}). Réponse : \all{Jonas will seine Farm in der Nähe von Bafoussam gründen.} « Jonas veut créer sa ferme près de Bafoussam. » Le verbe modal \all{will} est en 2\textsuperscript{e} position, l'infinitif \all{gründen} à la fin.
\item La question porte sur \cle{deux activités} : culture (\all{anbauen}) et élevage (\all{züchten}). Réponse : \all{Er will Mais und Gemüse anbauen und Hühner und Ziegen züchten.} « Il veut cultiver du maïs et des légumes et élever des poules et des chèvres. » Les deux infinitifs (\all{anbauen, züchten}) sont en fin de phrase.
\item La question \all{warum} appelle une \cle{finalité} : on répond avec \all{um ... zu} + infinitif. Réponse : \all{Er lernt viel über Landwirtschaft, um moderne Methoden kennenzulernen.} « Il apprend beaucoup sur l'agriculture pour connaître des méthodes modernes. » Virgule obligatoire avant \all{um}, infinitif \all{kennenzulernen} (verbe à particule : \all{zu} entre \all{kennen} et \all{lernen}) à la fin.
\item Question sur le futur (\all{was wird ... tun}). Réponse : \all{Sein Vater wird ihm Geld für das Feld geben.} « Son père lui donnera de l'argent pour le champ. » Futur I : \all{wird ... geben}. \all{ihm} = datif (à qui ?).
\item Question sur le but final (\all{Ziel}). Réponse : \all{Die Farm soll seine Familie ernähren und Einkommen bringen.} « La ferme doit nourrir sa famille et apporter des revenus. » \all{soll} (verbe modal \all{sollen} : obligation/but) + deux infinitifs à la fin.
\end{enumerate}
\textbf{Conclusion méthode} : une réponse complète reprend les termes de la question, conjugue le verbe et respecte l'ordre des mots allemand (verbe en 2, infinitifs à la fin).
\end{exoresolu}

\courssub{Savoir-faire 2 : Employer le vocabulaire de la ferme et de l'élevage}

\begin{methode}[Mémoriser et utiliser le lexique thématique]
\begin{enumerate}
\item \cle{Apprendre chaque nom avec son article et son pluriel} : \all{die Farm, -en} ; \all{das Feld, -er} ; \all{das Vieh} (singulier seulement, collectif) ; \all{der Mais} (singulier seulement) ; \all{das Gemüse} (singulier seulement) ; \all{das Huhn, die Hühner} ; \all{die Ziege, -n}.
\item \cle{Regrouper par familles de mots} (Wortfamilien) : \all{bauen → anbauen → der Anbau} ; \all{züchten → die Zucht → der Züchter} ; \all{die Landwirtschaft → der Landwirt} ; \all{ernähren → die Ernährung}.
\item \cle{Associer verbe et complément} (Kollokationen) : \all{Mais anbauen}, \all{Gemüse anbauen}, \all{Vieh züchten}, \all{Hühner züchten}, \all{eine Farm gründen}, \all{ein Feld kaufen}, \all{einen Stall bauen}.
\item \cle{Réutiliser le vocabulaire dans ses propres phrases} : écrire trois phrases minimum par mot nouveau.
\item \cle{Éviter les faux amis et calques} : \all{die Farm} = la ferme agricole (pas \all{der Hof} = la cour/ferme bâtie) ; « un champ » = \all{das Feld} (pas \all{der Champ} qui n'existe pas) ; « élever des animaux » = \all{Tiere züchten} (pas \all{erziehen} = éduquer des enfants).
\end{enumerate}
\end{methode}

\begin{exoresolu}[Vocabulaire — Phrases à compléter et transformer]
\textbf{Énoncé A.} Complétez avec le mot qui convient (forme correcte) : \all{das Vieh — der Mais — das Feld — anbauen — züchten — der Stall}.
\begin{enumerate}
\item Jonas kauft ein großes \dots\ in der Nähe von Bafoussam.
\item Auf diesem \dots\ will er \dots\ und Gemüse \dots.
\item Er will auch Hühner und Ziegen \dots.
\item \dots\ braucht viel Platz, Futter und Wasser.
\item Für die Tiere baut er einen \dots.
\end{enumerate}
\textbf{Énoncé B.} Mettez au pluriel (là où c'est possible) : \all{die Farm, das Feld, das Vieh, der Mais, die Ziege, der Traktor, die Ernte}.
\tcblower
\textbf{Solution détaillée A.}
\begin{enumerate}
\item \all{Jonas kauft ein großes Feld in der Nähe von Bafoussam.} (Accusatif neutre : \all{ein großes Feld}).
\item \all{Auf diesem Feld will er Mais und Gemüse anbauen.} (\all{diesem} = datif neutre après \all{auf} + lieu ; \all{anbauen} infinitif soudé à la fin).
\item \all{Er will auch Hühner und Ziegen züchten.} (\all{züchten} infinitif à la fin).
\item \all{Das Vieh braucht viel Platz, Futter und Wasser.} (Sujet singulier neutre \all{das Vieh} → verbe 3\textsuperscript{e} pers. sg. \all{braucht}).
\item \all{Für die Tiere baut er einen Stall.} (Accusatif masculin : \all{einen Stall} ; verbe \all{baut} en 2\textsuperscript{e} position).
\end{enumerate}
\textbf{Solution détaillée B.}
\all{die Farmen, die Felder, (das Vieh — pas de pluriel), (der Mais — pas de pluriel), die Ziegen, die Traktoren, die Ernten}.
\textbf{Conclusion méthode} : vérifier toujours l'article (genre), le nombre (singulier/pluriel), le cas (ici surtout accusatif/datif) et la place de la particule séparable (\all{anbauen} reste un mot à l'infinitif).
\end{exoresolu}

\courssub{Grammaire 1 : Le futur I (Futur I)}

\begin{propriete}[Formation et emploi du Futur I]
\begin{itemize}
\item \cle{Forme} : \all{werden} (conjugué au présent, 2\textsuperscript{e} position) + \cle{infinitif du verbe principal à la fin de la phrase}.
\item \cle{Conjugaison de \all{werden} (verbe fort, voyelle changeante)} :
\begin{center}
\begin{tabular}{|l|l|l|l|}
\hline
\rowcolor{popblueL} \textbf{Personne} & \textbf{werden} & \textbf{Exemple : gründen} & \textbf{Traduction} \\
\hline
ich & werde & ich werde gründen & je créerai / je vais créer \\
du & wirst & du wirst gründen & tu créeras / tu vas créer \\
er / sie / es & \textbf{wird} & er wird gründen & il/elle créera / il/elle va créer \\
wir & werden & wir werden gründen & nous créerons / nous allons créer \\
ihr & werdet & ihr werdet gründen & vous créerez / vous allez créer \\
sie / Sie & werden & sie werden gründen & ils/elles créeront / vous créerez \\
\hline
\end{tabular}
\end{center}
\item \cle{Emplois} : Projet futur (\all{nächstes Jahr, bald, später, in Zukunft}), prédiction, promesse, hypothèse.
\item \cle{Ordre des mots (Satzbau)} :
\begin{itemize}
\item Phrase simple : \all{Subjekt + werden + Compléments + Infinitiv.} (\all{Jonas wird eine Farm gründen.})
\item Phrase commençant par un complément (Topicalisation) : \all{Complément + werden + Subjekt + Compléments + Infinitiv.} (\all{Nächstes Jahr wird Jonas eine Farm gründen.})
\end{itemize}
\item \cle{Verbes à particule séparable} : l'infinitif final se ressoude : \all{anbauen, züchten, aufbauen}. Jamais \all{an bauen} ou \all{zu anbauen} (sauf avec \all{um ... zu} : \all{um anzubauen}).
\end{itemize}
\end{propriete}

\begin{experience}[Schéma visuel — Structure du Futur I]
\begin{popfigure}
\begin{tikzpicture}[node distance=0.8cm, every node/.style={font=\small}]
\node[draw, rounded corners, fill=popblueL, text=popdark] (S) {Sujet};
\node[draw, rounded corners, fill=poporangeL, text=popdark, right=of S] (W) {werden (conjugué)};
\node[draw, rounded corners, fill=popgreenL, text=popdark, right=of W] (K) {Compléments (Acc, Dat, Temps...)};
\node[draw, rounded corners, fill=poppinkL, text=popdark, right=of K] (I) {Infinitif (verbe principal)};
\draw[-{Stealth[length=3mm]}, thick, popblue] (S) -- (W);
\draw[-{Stealth[length=3mm]}, thick, poporange] (W) -- (K);
\draw[-{Stealth[length=3mm]}, thick, popgreen] (K) -- (I);
\node[below=0.5cm of W, align=center, text=popmuted] {\textbf{Position 2} \textbf{Position finale}};
\end{tikzpicture}
\legende{Ordre des mots : Le verbe conjugué \textit{werden} est ancré en 2\textsuperscript{e} position ; l'infinitif de sens va à la fin.}
\end{popfigure}
\end{experience}

\begin{methode}[Mettre une phrase au Futur I — Étapes]
\begin{enumerate}
\item Identifier le sujet (qui ?) → conjuguer \all{werden} pour ce sujet.
\item Placer \all{werden} en \textbf{2\textsuperscript{e} position}.
\item Garder les compléments (objet, lieu, temps) dans l'ordre habituel (Temps - Lieu - Objet / Datif avant Accusatif).
\item Mettre le verbe principal à l'\textbf{infinitif à la toute fin}.
\item \cle{Vérifier} : particule séparable ressoudée (\all{anbauen}), majuscules aux noms.
\end{enumerate}
\end{methode}

\begin{exoresolu}[Grammaire — Mettre au Futur I]
\textbf{Énoncé.} Transformez les phrases du présent au Futur I.
\begin{enumerate}
\item Jonas gründet eine Farm.
\item Er baut Mais und Gemüse an.
\item Seine Familie hilft ihm.
\item Ihr besucht die Farm bald.
\item Wir lernen viel über Landwirtschaft.
\end{enumerate}
\tcblower
\textbf{Solution détaillée.}
\begin{enumerate}
\item \all{Jonas wird eine Farm gründen.} (Sujet \all{Jonas} → \all{wird} ; infinitif \all{gründen} à la fin).
\item \all{Er wird Mais und Gemüse anbauen.} (Piège : présent \all{baut ... an} → futur \all{anbauen} soudé à la fin).
\item \all{Seine Familie wird ihm helfen.} (\all{die Familie} = singulier → \all{wird} ; \all{helfen} + datif \all{ihm}).
\item \all{Ihr werdet die Farm bald besuchen.} (Sujet \all{ihr} → \all{werdet} ; pas \all{werden} !).
\item \all{Wir werden viel über Landwirtschaft lernen.} (Sujet \all{wir} → \all{werden} ; infinitif \all{lernen}).
\end{enumerate}
\textbf{Conclusion méthode} : Futur I = \all{werden} (conjugué, pos. 2) + Infinitif (pos. finale). Terminaison de \all{werden} = terminaison du présent (sauf \all{er/sie/es wird}, \all{ihr werdet}).
\end{exoresolu}

\courssub{Grammaire 2 : La finalité (um ... zu)}

\begin{propriete}[La proposition de but avec \all{um ... zu}]
\begin{itemize}
\item \cle{Condition sine qua non} : Le sujet de la proposition principale et celui de l'infinitif sont \textbf{identiques} (coréférence). \all{Jonas arbeitet hart, um Erfolg zu haben.} (Jonas = sujet des deux actions).
\item \cle{Structure} : \all{Hauptsatz, um + Compléments + zu + Infinitiv (toute fin).}
\item \cle{Ponctuation} : \textbf{Virgule obligatoire} avant \all{um}.
\item \cle{Verbe à particule séparable} : \all{zu} s'insère \textbf{entre} la particule et le radical : \all{um Mais anzubauen}, \all{um Hühner zuzüchten}, \all{um den Stall aufzubauen}.
\item \cle{Verbe composé / préfixe inséparable} : \all{zu} devant le verbe : \all{um zu lernen}, \all{um zu besuchen}, \all{um zu verwirklichen}.
\item \cle{Traduction} : « pour / afin de / dans le but de » + infinitif présent.
\item \cle{Sujets différents} : Utiliser \all{damit} + phrase complète conjuguée (ex. \all{Jonas arbeitet hart, damit seine Familie gut lebt.} — niveau B1, à reconnaître).
\end{itemize}
\end{propriete}

\begin{methode}[Construire une phrase avec \all{um ... zu} — Étapes]
\begin{enumerate}
\item Écrire la phrase principale (verbe conjugué en 2).
\item Ajouter une \textbf{virgule}.
\item Écrire \all{um}.
\item Placer les compléments (temps, lieu, objet, manière) \textbf{entre \all{um} et \all{zu}}.
\item Écrire \all{zu} + \textbf{infinitif} (verbe principal, particule ressoudée si séparable) \textbf{à la toute fin}.
\item Vérifier : même sujet ? virgule ? \all{zu} bien placé ?
\end{enumerate}
\end{methode}

\begin{exoresolu}[Grammaire — Relier deux phrases avec \all{um ... zu}]
\textbf{Énoncé.} Reliez les deux phrases en une seule avec \all{um ... zu}. Attention au même sujet !
\begin{enumerate}
\item Jonas spart Geld. Er kauft ein Feld.
\item Er lernt Deutsch. Er arbeitet später mit einer Firma aus Deutschland.
\item Er kauft Hühner. Er züchtet sie auf seiner Farm.
\item Er baut einen Stall. Er schützt das Vieh vor Regen.
\item Wir lernen Vokabeln. Wir bestehen die Prüfung.
\end{enumerate}
\tcblower
\textbf{Solution détaillée.}
\begin{enumerate}
\item \all{Jonas spart Geld, um ein Feld zu kaufen.} (Sujet \all{Jonas/Er} identique ; \all{zu kaufen} final).
\item \all{Er lernt Deutsch, um später mit einer Firma aus Deutschland zu arbeiten.} (Compléments \all{später, mit einer Firma...} entre \all{um} et \all{zu arbeiten}).
\item \all{Er kauft Hühner, um sie auf seiner Farm zu züchten.} (Pronom objet \all{sie} placé juste après \all{um} ; \all{züchten} final).
\item \all{Er baut einen Stall, um das Vieh vor Regen zu schützen.} (Verbe \all{schützen} inséparable → \all{zu schützen} ; \all{vor Regen} = protection contre la pluie).
\item \all{Wir lernen Vokabeln, um die Prüfung zu bestehen.} (Sujet \all{Wir} identique ; \all{bestehen} final).
\end{enumerate}
\textbf{Conclusion méthode} : Schéma universel — \all{Hauptsatz, um + (Objet/Compléments) + zu + Infinitiv}.
\end{exoresolu}

\courssub{Expression écrite et orale guidée : Traduction (Thème) et Production}

\begin{methode}[Traduire du français vers l'allemand (Thème) et rédiger un paragraphe]
\begin{enumerate}
\item \cle{Analyser le temps} : Projet futur certain → Futur I (\all{werden} + infinitif). Intention présente → \all{wollen} + infinitif. Action actuelle → Présent.
\item \cle{Construire le squelette phrase par phrase} : Sujet + Verbe conjugué (2) + Compléments + Infinitif/Participe (fin).
\item \cle{Traduire « pour » + infinitif} systématiquement par \all{um ... zu} (jamais \all{für ... zu}).
\item \cle{Gérer les verbes à particule} : \all{cultiver} = \all{anbauen} (infinitif \all{anbauen}) ; \all{élever} = \all{züchten}.
\item \cle{Pour un paragraphe guidé (4–6 phrases)} : 1. Présenter la personne/le projet (\all{Ich will / Ich werde ... gründen}). 2. Lieu (\all{in der Nähe von ...}). 3. Cultures (\all{anbauen}). 4. Élevage (\all{züchten}). 5. But / Préparation (\all{um ... zu} / \all{lernen}). 6. Soutien / Conclusion (\all{Familie hilft / wird helfen}).
\item \cle{Relire (Check-list)} : Majuscules \textbf{à tous les noms} ? Umlaute (ä, ö, ü) et ß ? Terminaisons verbes ? Cas (accusatif/datif) ? Place verbe (2 / fin) ? Virgule avant \all{um} ?
\end{enumerate}
\end{methode}

\begin{exoresolu}[Thème et Production écrite — Type Bac]
\textbf{Énoncé A (Thème).} Traduisez en allemand :
« L'année prochaine, mon frère créera une petite ferme près de Bafoussam. Il cultivera du maïs et élèvera des chèvres. Il travaille dur pour réussir. »

\textbf{Énoncé B (Production guidée).} Rédigez 5 à 6 phrases en allemand : votre projet de ferme (lieu, cultures, animaux, préparation, but, aide familiale).
\tcblower
\textbf{Solution détaillée A.}
\begin{itemize}
\item « L'année prochaine » = \all{Nächstes Jahr} (complément de temps en tête → inversion).
\item « mon frère créera » = Futur I → \all{wird mein Bruder ... gründen}.
\item « une petite ferme près de Bafoussam » = \all{eine kleine Farm in der Nähe von Bafoussam} (accusatif féminin : \all{eine kleine Farm}).
\item « il cultivera du maïs » = \all{er wird Mais anbauen} (particule ressoudée).
\item « et élèvera des chèvres » = \all{und Ziegen züchten} (deuxième infinitif coordonné, pas de \all{werden} répété).
\item « Il travaille dur » = Présent (action actuelle) → \all{Er arbeitet hart}.
\item « pour réussir » = \all{um erfolgreich zu sein} ou \all{um Erfolg zu haben}.
\end{itemize}
\textbf{Traduction complète :}
\all{Nächstes Jahr wird mein Bruder eine kleine Farm in der Nähe von Bafoussam gründen. Er wird Mais anbauen und Ziegen züchten. Er arbeitet hart, um erfolgreich zu sein.}

\textbf{Solution détaillée B (Modèle de production — Niveau A2).}
\all{Ich will später eine Farm in der Nähe von Yaoundé gründen. Nächstes Jahr werde ich ein Feld kaufen. Ich werde Mais, Maniok und Bohnen anbauen. Ich werde auch Hühner und Schweine züchten. Ich lerne jetzt viel über Landwirtschaft, um meine Farm gut zu führen. Meine Familie wird mir helfen.}

\textit{Traduction : « Je veux créer plus tard une ferme près de Yaoundé. L'année prochaine j'achèterai un champ. Je cultiverai du maïs, du manioc et des haricots. J'élèverai aussi des poules et des porcs. J'apprends maintenant beaucoup sur l'agriculture pour bien gérer ma ferme. Ma famille m'aidera. »}

\textbf{Conclusion méthode} : Le texte utilise \all{wollen} (intention), Futur I (projets datés), présent (action actuelle), \all{um ... zu} (but). Chaque phrase a son verbe conjugué en 2\textsuperscript{e} position.
\end{exoresolu}

\begin{attention}[Les 5 erreurs qui coûtent des points au Bac — Check-list finale]
\begin{enumerate}
\item \cle{Majuscules aux noms oubliées} : \all{die farm, das vieh, der mais} → \textbf{Fautif}. Toujours : \all{die Farm, das Vieh, der Mais}. Toute copie avec noms en minuscules perd des points d'orthographe systématiquement.
\item \cle{Verbe mal placé au Futur I} : \all{Jonas wird gründen eine Farm} → \textbf{Faux}. Règle d'or : \all{werden} en 2, \textbf{infinitif à la fin} : \all{Jonas wird eine Farm gründen.}
\item \cle{Particule séparable coupée à l'infinitif} : \all{Er wird Mais an bauen} ou \all{um zu anbauen} → \textbf{Faux}. Correct : \all{Er wird Mais anbauen} / \all{um Mais anzubauen} (un seul mot).
\item \cle{« Pour » traduit par \all{für}} : \all{Er arbeitet für erfolgreich zu sein} → \textbf{Faux}. \all{für} + nom (\all{für die Familie}). But + verbe = \all{um ... zu} : \all{um erfolgreich zu sein}.
\item \cle{Calques du français (Anglicismes/Gallicismes)} : « Créer une activité » ≠ \all{eine Aktivität kreieren}. Dire : \all{eine Farm gründen}, \all{ein Unternehmen aufbauen}. « Faire de l'élevage » ≠ \all{Elevage machen}. Dire : \all{Vieh züchten}, \all{Tierhaltung betreiben}. « Prendre une décision » ≠ \all{eine Entscheidung nehmen}. Dire : \all{sich entscheiden}, \all{eine Entscheidung treffen}.
\end{enumerate}
\end{attention}

\newpage

\coursec{Exercices}

\courssub{Je vérifie mes connaissances}

\begin{exercice}[QCM : le futur I et la finalité]
Entoure la bonne réponse.
\begin{enumerate}
\item Demain, nous \dots\ des légumes sur le champ.
\begin{itemize}
\item a) bauen \dots an
\item b) werden \dots anbauen
\item c) wollen \dots anbauen
\end{itemize}
\item Jonas spricht mit der Bank, \dots\ einen Kredit zu bekommen.
\begin{itemize}
\item a) damit
\item b) um
\item c) für
\end{itemize}
\item Quelle phrase est au futur I ?
\begin{itemize}
\item a) Sie züchtet Vieh.
\item b) Sie hat Vieh gezüchtet.
\item c) Sie wird Vieh züchten.
\end{itemize}
\item \all{um \dots\ zu} exprime :
\begin{itemize}
\item a) la cause
\item b) la conséquence
\item c) la finalité
\end{itemize}
\end{enumerate}
\end{exercice}

\begin{exercice}[Vrai ou faux ? Justifie ta réponse]
Réponds par \all{richtig} ou \all{falsch} et justifie en une phrase en français.
\begin{enumerate}
\item Au futur I, l'infinitif se place à la fin de la phrase.
\item Le verbe \all{anbauen} est un verbe à particule inséparable.
\item On peut utiliser \all{um \dots\ zu} seulement si les deux actions ont le même sujet.
\item \all{das Vieh} signifie « le champ ».
\end{enumerate}
\end{exercice}

\begin{exercice}[Vocabulaire : la ferme et l'élevage]
Complète chaque phrase avec le mot qui convient : \all{die Farm}, \all{das Vieh}, \all{der Mais}, \all{das Feld}, \all{anbauen}, \all{züchten}.
\begin{enumerate}
\item Jonas kauft ein großes \dots\ hinter dem Dorf.
\item Auf seiner \dots\ arbeiten fünf Personen.
\item Im März will er \dots\ und Gemüse \dots\ .
\item Er will auch Rinder und Ziegen \dots\ .
\item \dots\ braucht jeden Tag Wasser und Futter.
\end{enumerate}
\end{exercice}

\begin{exercice}[Conjugaison : mets les phrases au futur I]
Transforme les phrases comme dans l'exemple : \all{Er baut Mais an.} $\rightarrow$ \all{Er wird Mais anbauen.}
\begin{enumerate}
\item Sie züchtet Rinder.
\item Wir kaufen ein Feld.
\item Ich spreche mit der Bank.
\item Ihr verkauft das Gemüse auf dem Markt.
\item Die Familie arbeitet auf der Farm.
\end{enumerate}
\end{exercice}

\courssub{Je m'entraîne}

\begin{exercice}[Relie les phrases avec \all{um \dots\ zu}]
Exemple : \all{Er spricht mit der Bank. Er will einen Kredit bekommen.} $\rightarrow$ \all{Er spricht mit der Bank, um einen Kredit zu bekommen.}
\begin{enumerate}
\item Jonas kauft ein Feld. Er will Mais anbauen.
\item Sie lernt Deutsch. Sie will in Deutschland studieren.
\item Wir arbeiten jeden Tag. Wir wollen das Vieh füttern.
\item Er verkauft Gemüse. Er will Geld verdienen.
\item Ich gehe zur Schule. Ich will Landwirt werden.
\end{enumerate}
\end{exercice}

\begin{exercice}[Texte à trous : le projet de Jonas]
Complète le texte avec les mots de la liste : \all{wird}, \all{um}, \all{Mais}, \all{Vieh}, \all{Feld}, \all{anzubauen}, \all{zu}, \all{züchten}.

\begin{textebox}[Jonas' Plan]
Jonas hat ein großes (1) \dots\ gekauft. Nächstes Jahr (2) \dots\ er dort (3) \dots\ und Gemüse (4) \dots\ . Er will auch (5) \dots\ (6) \dots\ : Rinder, Ziegen und Hühner. Er spricht mit der Bank, (7) \dots\ einen Kredit (8) \dots\ bekommen. Sein Traum ist eine moderne Farm für das ganze Dorf.
\end{textebox}
\end{exercice}

\begin{exercice}[Appariements : mots et définitions]
Relie chaque mot allemand à sa définition française.
\begin{center}
\begin{tabularx}{\linewidth}{|X|X|}
\hline
\rowcolor{popblueL} Deutsch & Français \\
\hline
1. das Vieh & a) cultiver (plantes, céréales) \\
\hline
2. anbauen & b) le bétail, les animaux de la ferme \\
\hline
3. züchten & c) le prêt (argent de la banque) \\
\hline
4. der Kredit & d) élever (des animaux) \\
\hline
5. der Markt & e) le marché \\
\hline
\end{tabularx}
\end{center}
\end{exercice}

\begin{exercice}[Thème : traduis en allemand]
\begin{enumerate}
\item L'année prochaine, nous élèverons des bovins pour nourrir le village.
\item Jonas achètera un champ pour cultiver du maïs.
\item Je parle avec le fermier pour apprendre le métier.
\item Ils vendront leurs légumes au marché de Douala.
\end{enumerate}
\end{exercice}

\courssub{Je me prépare au Bac}

\begin{exercice}[Compréhension écrite type Bac]
Lis le texte suivant, puis réponds aux questions en allemand, par des phrases complètes.

\begin{textebox}[Ein neuer Anfang in der Nähe von Bafoussam]
Seit einem Jahr hat Amina, 24 Jahre alt, eine kleine Farm in der Nähe von Bafoussam. Nach der Schule hat sie zwei Jahre auf der Farm ihres Onkels gearbeitet, um Erfahrung zu sammeln. Dann hat sie mit ihrer Familie gesprochen und ein Feld gekauft.

„Im ersten Jahr werde ich Mais und Bohnen anbauen“, sagt sie. „Später werde ich auch Hühner und Ziegen züchten, um mehr Geld zu verdienen.“ Amina hat einen kleinen Kredit von der Bank bekommen. Mit diesem Geld wird sie Werkzeuge und Samen kaufen.

Ihre Familie hilft ihr jeden Tag. Ihr Bruder füttert die Tiere, ihre Mutter verkauft das Gemüse auf dem Markt. „Viele junge Leute gehen in die Stadt“, erklärt Amina. „Aber ich bleibe hier, um mein Dorf zu entwickeln. In fünf Jahren wird meine Farm zwanzig Personen Arbeit geben.“
\end{textebox}

\textbf{Glossaire :} die Erfahrung, -en : l'expérience ; das Werkzeug, -e : l'outil ; der Samen, - : la semence ; entwickeln : développer.

\begin{enumerate}
\item Wer ist Amina und wo arbeitet sie ?
\item Warum hat Amina auf der Farm ihres Onkels gearbeitet ?
\item Was wird Amina im ersten Jahr anbauen ?
\item Warum will sie später Hühner und Ziegen züchten ?
\item Wer hilft Amina auf der Farm ? Was machen diese Personen ?
\item Was wird Aminas Farm in fünf Jahren machen ?
\item Vrai ou faux ? Justifie avec une phrase du texte :
\begin{itemize}
\item a) Amina hat kein Geld von der Bank bekommen.
\item b) Amina will in die Stadt gehen.
\end{itemize}
\end{enumerate}

\textbf{Questions de langue :}
\begin{enumerate}
\item Relève deux phrases au futur I et traduis-les en français.
\item Relève une phrase avec \all{um \dots\ zu} et explique ce qu'elle exprime.
\item Mets au futur I : \all{Amina kauft Werkzeuge.} / \all{Ihr Bruder füttert die Tiere.}
\end{enumerate}
\end{exercice}

\begin{exercice}[Thème et version type Bac]
\textbf{A. Thème.} Traduis en allemand (attention au futur I et à la place du verbe) :
\begin{enumerate}
\item Mon frère construira une ferme près de Yaoundé.
\item Nous cultivons du maïs pour vendre la récolte au marché.
\item Elle travaillera beaucoup pour aider sa famille.
\item L'année prochaine, j'élèverai des chèvres et des poules.
\end{enumerate}

\textbf{B. Version.} Traduis en français le texte suivant :

\begin{textebox}[Die Farm der Familie Berger]
Familie Berger lebt in Bayern. Im Frühling wird Herr Berger Kartoffeln und Weizen anbauen. Seine Frau Anna wird Gemüse im Garten pflanzen, um frische Produkte für die Familie zu haben. Der Sohn Max lernt in der Landwirtschaftsschule, um später die Farm zu übernehmen. „In zehn Jahren werden wir mehr Felder kaufen und mehr Vieh züchten“, sagt Herr Berger. Die Familie verkauft ihre Produkte jeden Samstag auf dem Markt in der Stadt.
\end{textebox}

\textbf{Glossaire :} der Weizen : le blé ; pflanzen : planter ; übernehmen : reprendre ; das Produkt, -e : le produit.
\end{exercice}

\begin{exercice}[Production écrite guidée : Dein Projekt für die Zukunft]
Rédige un texte de 80 à 100 mots en allemand sur le sujet suivant : \all{Mein Projekt für die Zukunft : eine Farm in meinem Dorf}.

Tu dois :
\begin{itemize}
\item présenter ton projet (quoi ? où ? avec qui ?) ;
\item utiliser au moins quatre verbes au futur I (\all{werden} + infinitif) ;
\item utiliser au moins deux fois la construction \all{um \dots\ zu} ;
\item employer au moins cinq mots du vocabulaire de la leçon (\all{die Farm}, \all{das Vieh}, \all{der Mais}, \all{das Feld}, \all{anbauen}, \all{züchten}, \all{der Kredit}, \all{der Markt}\dots).
\end{itemize}

\textbf{Aide :} commence par exemple par \all{Nach der Schule werde ich \dots} et termine par \all{Mein Traum ist \dots}.
\end{exercice}

\newpage

\coursec{Corrigés des exercices}

\begin{corrige}[Exercice 1 — QCM : Le futur I et la finalité]
\begin{enumerate}
\item \textbf{b)} \all{Morgen werden wir Gemüse auf dem Feld anbauen.} (Demain, nous cultiverons des légumes dans le champ.) \\
Le \cle{Futur I} se forme avec l'auxiliaire \all{werden} conjugué en \textbf{2ᵉ position} et l'infinitif du verbe principal à la \textbf{fin de phrase}. La réponse a) correspond au présent ; la réponse c) (\all{wollen} + infinitif) exprime la volonté, pas le futur.
\item \textbf{b)} \all{Jonas spricht mit der Bank, um einen Kredit zu bekommen.} (Jonas parle à la banque pour obtenir un prêt.) \\
La construction \cle{um ... zu} + infinitif exprime la \cle{finalité} (« dans le but de / pour »). \all{damit} introduit une subordonnée avec verbe conjugué (sujet potentiellement différent) ; \all{für} est une préposition suivie d'un nom (ex. \all{für den Kredit}).
\item \textbf{c)} \all{Sie wird Vieh züchten.} (Elle élèvera du bétail.) \\
\all{wird} = 3ᵉ personne du singulier de \all{werden} (2ᵉ position) + \all{züchten} (infinitif à la fin) = Futur I. La réponse a) est au présent, la réponse b) au \cle{Perfekt} (\all{hat ... gezüchtet}).
\item \textbf{c)} La finalité. La structure \all{um ... zu} répond à la question « \all{Warum?} / \all{Wozu?} » (pourquoi / dans quel but).
\end{enumerate}
\end{corrige}

\begin{corrige}[Exercice 2 — Vrai ou faux ? Justifie ta réponse]
\begin{enumerate}
\item \all{richtig}. Au Futur I, le verbe conjugué \all{werden} occupe la 2ᵉ position et l'infinitif principal est rejeté en fin de phrase : \all{Er wird Mais anbauen.}
\item \all{falsch}. \all{anbauen} est un \cle{verbe à particule séparable} : au présent, la particule \all{an-} se détache et va en fin de phrase → \all{Er baut Mais an.}
\item \all{richtig}. Avec \cle{um ... zu}, le sujet de l'infinitif doit être \textbf{identique} à celui de la proposition principale ; si les sujets diffèrent, on utilise \all{damit} + subordonnée (verbe conjugué à la fin).
\item \all{falsch}. \all{das Vieh} (neutre, pas de pluriel usuel) signifie « le bétail / les animaux d'élevage » ; « le champ » se dit \all{das Feld} (pl. \all{die Felder}).
\end{enumerate}
\end{corrige}

\begin{corrige}[Exercice 3 — Vocabulaire : la ferme et l'élevage]
\begin{enumerate}
\item Jonas kauft ein großes \textbf{Feld} hinter dem Dorf. (accusatif neutre : \all{ein großes Feld})
\item Auf seiner \textbf{Farm} arbeiten fünf Personen. (datif : \all{auf seiner Farm})
\item Im März will er \textbf{Mais} und Gemüse \textbf{anbauen}. (infinitif après \all{will} ; particule \all{an-} ressoudée à l'infinitif en fin de phrase)
\item Er will auch Rinder und Ziegen \textbf{züchten}. (infinitif \all{züchten} en fin de phrase)
\item \textbf{Das Vieh} braucht jeden Tag Wasser und Futter. (sujet nominatif, majuscule en tête de phrase)
\end{enumerate}
\textbf{Traduction :} 1. Jonas achète un grand champ derrière le village. 2. Cinq personnes travaillent dans sa ferme. 3. En mars, il veut cultiver du maïs et des légumes. 4. Il veut aussi élever des bovins et des chèvres. 5. Le bétail a besoin d'eau et de nourriture chaque jour.
\end{corrige}

\begin{corrige}[Exercice 4 — Conjugaison : mets les phrases au futur I]
\begin{enumerate}
\item \all{Sie wird Rinder züchten.} (Elle élèvera des bovins.)
\item \all{Wir werden ein Feld kaufen.} (Nous achèterons un champ.)
\item \all{Ich werde mit der Bank sprechen.} (Je parlerai avec la banque.)
\item \all{Ihr werdet das Gemüse auf dem Markt verkaufen.} (Vous vendrez les légumes au marché.)
\item \all{Die Familie wird auf der Farm arbeiten.} (La famille travaillera dans la ferme.)
\end{enumerate}
\textbf{Points de vigilance :}
\begin{itemize}
\item Conjuguer correctement \all{werden} : \all{ich werde}, \all{du wirst}, \all{er/sie/es wird}, \all{wir werden}, \all{ihr werdet}, \all{sie/Sie werden}.
\item Placer l'infinitif principal à la \textbf{toute fin} de la phrase (ou de la proposition).
\item Verbe à particule séparable (\all{anbauen}, \all{aufbauen}, \all{mitarbeiten}...) : la particule se \textbf{ressoude} à l'infinitif → \all{wird ... anbauen} (jamais \all{wird ... an ... bauen}).
\end{itemize}
\end{corrige}

\begin{corrige}[Exercice 5 — Relie les phrases avec \cle{um ... zu}]
\begin{enumerate}
\item \all{Jonas kauft ein Feld, um Mais anzubauen.} (Jonas achète un champ pour cultiver du maïs.)
\item \all{Sie lernt Deutsch, um in Deutschland zu studieren.} (Elle apprend l'allemand pour étudier en Allemagne.)
\item \all{Wir arbeiten jeden Tag, um das Vieh zu füttern.} (Nous travaillons chaque jour pour nourrir le bétail.)
\item \all{Er verkauft Gemüse, um Geld zu verdienen.} (Il vend des légumes pour gagner de l'argent.)
\item \all{Ich gehe zur Schule, um Landwirt zu werden.} (Je vais à l'école pour devenir agriculteur.)
\end{enumerate}
\textbf{Points de vigilance :}
\begin{itemize}
\item \textbf{Virgule obligatoire} devant \all{um}.
\item L'infinitif précédé de \all{zu} va en \textbf{fin de phrase}.
\item Verbe à particule séparable : \all{zu} s'insère \textbf{entre} la particule et le radical → \all{anzubauen} (et non \all{zu anbauen}), \all{mitzuarbeiten}, \all{aufzugeben}.
\item Sujet de l'infinitif = sujet de la principale.
\end{itemize}
\end{corrige}

\begin{corrige}[Exercice 6 — Texte à trous : le projet de Jonas]
\begin{enumerate}
\item \textbf{Feld} : accusatif neutre après \all{ein großes} (Jonas a acheté un grand champ).
\item \textbf{wird} : auxiliaire \all{werden} au Futur I, 3ᵉ pers. sg., en 2ᵉ position après \all{Nächstes Jahr}.
\item \textbf{Mais} : complément d'objet direct (le maïs), masculin.
\item \textbf{anbauen} : infinitif du verbe à particule séparable, en fin de phrase (Futur I : \all{wird ... anbauen}). \emph{Attention : pas de « zu » ici, car pas de construction de finalité.}
\item \textbf{Vieh} : complément de \all{züchten} (élever du bétail), neutre.
\item \textbf{züchten} : infinitif après le modal \all{will} (veut).
\item \textbf{um} : introduit la proposition de finalité (\all{um ... zu}).
\item \textbf{zu} : particule préfixée à l'infinitif \all{bekommen} (\all{um einen Kredit zu bekommen}).
\end{enumerate}
\textbf{Texte complet :}
\begin{textebox}[Projet de Jonas]
Jonas hat ein großes \textbf{Feld} gekauft. Nächstes Jahr \textbf{wird} er dort \textbf{Mais} und Gemüse \textbf{anbauen}. Er will auch \textbf{Vieh} \textbf{züchten} : Rinder, Ziegen und Hühner. Er spricht mit der Bank, \textbf{um} einen Kredit \textbf{zu} bekommen. Sein Traum ist eine moderne Farm für das ganze Dorf.
\end{textebox}
\textbf{Traduction :} « Jonas a acheté un grand champ. L'année prochaine, il y cultivera du maïs et des légumes. Il veut aussi élever du bétail : des bovins, des chèvres et des poules. Il parle avec la banque pour obtenir un prêt. Son rêve est une ferme moderne pour tout le village. »
\end{corrige}

\begin{corrige}[Exercice 7 — Appariements : mots et définitions]
\begin{center}
\begin{tabularx}{\linewidth}{|X|X|X|}
\hline
\rowcolor{popblueL} \textbf{Deutsch} & \textbf{Français} & \textbf{Réponse} \\
\hline
das Vieh & le bétail, les animaux de la ferme (coll. sg.) & 1 -- b \\
\hline
anbauen & cultiver (plantes, céréales, légumes) & 2 -- a \\
\hline
züchten & élever (des animaux), faire de la sélection & 3 -- d \\
\hline
der Kredit & le prêt (bancaire), le crédit & 4 -- c \\
\hline
der Markt & le marché (lieu de vente) & 5 -- e \\
\hline
\end{tabularx}
\end{center}
\textbf{Astuce mémorielle :} \all{anbauen} = « construire sur / planter » (\all{bauen} = construire) $\rightarrow$ plantes. \all{züchten} = « élever, dresser » $\rightarrow$ animaux. \textbf{Faux ami :} \all{der Kredit} = prêt bancaire (pas « carte de crédit» = \all{die Kreditkarte}).
\end{corrige}

\begin{corrige}[Exercice 8 — Thème : traduis en allemand]
\begin{enumerate}
\item \all{Nächstes Jahr werden wir Rinder züchten, um das Dorf zu ernähren.} \\
\emph{Vigilance :} \all{nächstes Jahr} (sans article, 2ᵉ position) → \all{werden} (1ᵉ pers. pl.) → infinitif \all{züchten} (fin) ; finalité \all{um ... zu} (même sujet « wir »).
\item \all{Jonas wird ein Feld kaufen, um Mais anzubauen.} \\
\emph{Vigilance :} verbe à particule \all{anbauen} → \all{zu} s'insère : \all{anzubauen}.
\item \all{Ich spreche mit dem Bauern, um den Beruf zu lernen.} \\
\emph{Vigilance :} présent (phrase française au présent) ; \all{mit} + \cle{Datif} → \all{dem Bauern} (n-déclinaison : \all{der Bauer, des Bauers, dem Bauern, den Bauern}) ; \all{der Beruf} (masc.).
\item \all{Sie werden ihr Gemüse auf dem Markt in Douala verkaufen.} \\
\emph{Vigilance :} Futur I (\all{werden} 3ᵉ pl.) ; \all{auf dem Markt} (lieu → Datif) ; \all{ihr Gemüse} (possessif 3ᵉ pl., neutre accusatif).
\end{enumerate}
\textbf{Barème indicatif (8 pts) :} 2 pts par phrase (1 pt temps/ordre, 1 pt vocabulaire/cas).
\end{corrige}

\begin{corrige}[Exercice 9 — Compréhension écrite type Bac]
\textbf{Réponses aux questions de compréhension (phrases complètes, allemand) :}
\begin{enumerate}
\item \all{Amina ist eine junge Frau. Sie ist 24 Jahre alt und arbeitet auf ihrer kleinen Farm in der Nähe von Bafoussam.}
\item \all{Sie hat auf der Farm ihres Onkels gearbeitet, um Erfahrung zu sammeln.}
\item \all{Im ersten Jahr wird sie Mais und Bohnen anbauen.}
\item \all{Sie will später Hühner und Ziegen züchten, um mehr Geld zu verdienen.}
\item \all{Ihre Familie hilft ihr: ihr Bruder füttert die Tiere und ihre Mutter verkauft das Gemüse auf dem Markt.}
\item \all{In fünf Jahren wird ihre Farm zwanzig Personen Arbeit geben.}
\item a) \all{falsch.} Justification : \all{„Amina hat einen kleinen Kredit von der Bank bekommen.“} \\
   b) \all{falsch.} Justification : \all{„Aber ich bleibe hier, um mein Dorf zu entwickeln.“}
\end{enumerate}
\textbf{Questions de langue :}
\begin{enumerate}
\item Deux phrases au Futur I relevées dans le texte :
      \begin{itemize}
      \item \all{„Im ersten Jahr werde ich Mais und Bohnen anbauen.“}
      \item \all{„In fünf Jahren wird meine Farm zwanzig Personen Arbeit geben.“}
      \end{itemize}
      (Autre occurrence possible : \all{„Später werde ich auch Hühner und Ziegen züchten.“})
\item \all{„Aber ich bleibe hier, um mein Dorf zu entwickeln.“} : la construction \cle{um ... zu} exprime ici la \textbf{finalité} (le but poursuivi) : elle reste \emph{afin de} / \emph{dans le but de} développer son village.
\item Transformations au Futur I :
      \begin{itemize}
      \item \all{Amina wird Werkzeuge kaufen.} (Amina achètera des outils.)
      \item \all{Ihr Bruder wird die Tiere füttern.} (Son frère nourrira les animaux.)
      \end{itemize}
\end{enumerate}
\textbf{Barème indicatif (20 pts) :} Compréhension 12 pts (environ 1,5 pt/question, justification V/F incluse) ; Langue 8 pts (relevés/traductions 4 pts, explication finalité 2 pts, transformations 2 pts).
\textbf{Consignes :} Répondre par des phrases complètes (verbe en 2ᵉ position) ; adapter la personne (ne pas recopier bêtement le « ich » du texte si la question porte sur « elle ») ; citation exacte pour le V/F.
\end{corrige}

\begin{corrige}[Exercice 10 — Thème et Version type Bac]
\textbf{A. Thème (Français $\rightarrow$ Allemand)}
\begin{enumerate}
\item \all{Mein Bruder wird eine Farm in der Nähe von Yaoundé bauen.} \\
      (\all{in der Nähe von} + Génitif formel / Datif courant : \all{von Yaoundé} ; \all{bauen} = construire/édifier).
\item \all{Wir bauen Mais an, um die Ernte auf dem Markt zu verkaufen.} \\
      (Présent $\rightarrow$ Présent ; \all{anbauen} séparé : \all{bauen ... an} ; \all{die Ernte} = la récolte).
\item \all{Sie wird viel arbeiten, um ihrer Familie zu helfen.} \\
      \textbf{Attention :} \cle{helfen} + \cle{Datif} → \all{ihrer Familie} (féminin sg., pas \all{*ihre Familie} accusatif).
\item \all{Nächstes Jahr werde ich Ziegen und Hühner züchten.} \\
      (\all{die Ziege, -n} = la chèvre ; \all{das Huhn, die Hühner} = la poule / les poules).
\end{enumerate}

\textbf{B. Version (Allemand $\rightarrow$ Français) — Corrigé}
\begin{textebox}[Texte original allemand (support de la version)]
Die Familie Berger lebt in Bayern. Im Frühling wird Herr Berger Kartoffeln und Weizen anbauen. Seine Frau Anna wird Gemüse im Garten pflanzen, um frische Produkte für die Familie zu haben. Der Sohn Max lernt an der Landwirtschaftsschule, um den Hof später zu übernehmen. „In zehn Jahren werden wir mehr Felder kaufen und mehr Vieh züchten“, sagt Herr Berger. Die Familie verkauft ihre Produkte jeden Samstag auf dem Stadtmarkt.
\end{textebox}

\textbf{Traduction française attendue :}
« La famille Berger vit en Bavière. Au printemps, M. Berger cultivera des pommes de terre et du blé. Sa femme Anna plantera des légumes dans le jardin pour avoir des produits frais pour la famille. Le fils Max apprend à l'école d'agriculture pour reprendre la ferme plus tard. „Dans dix ans, nous achèterons plus de champs et nous élèverons plus de bétail“, dit M. Berger. La famille vend ses produits chaque samedi au marché de la ville. »

\textbf{Barème indicatif (20 pts) :} Thème 8 pts (2 pts/phrase) ; Version 12 pts (fidélité sens 8 pts, qualité français 4 pts).
\textbf{Points de vigilance :} Thème : place de \all{werden} (2ᵉ) et infinitif (fin) ; Version : rendre les Futur I allemands par des futurs français ; respecter les guillemets allemands „ “ dans la citation.
\end{corrige}

\begin{corrige}[Exercice 11 — Production écrite guidée : \all{Dein Projekt für die Zukunft}]
\textbf{Proposition de rédaction modèle (environ 95 mots) :}
\begin{textebox}[Mein Projekt für die Zukunft]
Nach der Schule werde ich in meinem Dorf eine kleine Farm gründen. Zuerst werde ich mit meinem Vater sprechen, um ein Feld zu kaufen. Dann werde ich Mais, Bohnen und Gemüse anbauen. Ich werde auch Hühner und Ziegen züchten, um mehr Geld zu verdienen. Mein Bruder wird mir helfen: er wird jeden Morgen das Vieh füttern. Ich werde mit der Bank sprechen, um einen kleinen Kredit zu bekommen. Mit diesem Geld werde ich Werkzeuge und Samen kaufen. Später werden wir unsere Produkte auf dem Markt verkaufen. Mein Traum ist eine moderne Farm, die dem ganzen Dorf Arbeit gibt.
\end{textebox}

\textbf{Traduction :} « Après l'école, je créerai une petite ferme dans mon village. D'abord, je parlerai avec mon père pour acheter un champ. Ensuite, je cultiverai du maïs, des haricots et des légumes. J'élèverai aussi des poules et des chèvres pour gagner plus d'argent. Mon frère m'aidera : il nourrira le bétail chaque matin. Je parlerai avec la banque pour obtenir un petit prêt. Avec cet argent, j'achèterai des outils et des semences. Plus tard, nous vendrons nos produits au marché. Mon rêve est une ferme moderne qui donne du travail à tout le village. »

\textbf{Barème indicatif (20 pts) :}
\begin{itemize}
\item Respect du sujet et longueur (80--120 mots) : 4 pts
\item Futur I correct (au moins 4 emplois distincts) : 6 pts
\item \cle{um ... zu} correct (au moins 2 emplois, virgule, \all{zu} inséré si verbe séparable) : 4 pts
\item Vocabulaire de la leçon (min. 5 mots : \all{Farm, Feld, Mais, anbauen, züchten, Vieh, Kredit, Markt, Samen, Werkzeuge}...) : 4 pts
\item Orthographe, majuscules (noms), ponctuation : 2 pts
\end{itemize}
\textbf{Points de vigilance :}
\begin{itemize}
\item \textbf{Majuscule systématique} à tous les noms : \all{Farm, Feld, Vater, Mais, Bohnen, Gemüse, Hühner, Ziegen, Geld, Bruder, Morgen, Vieh, Bank, Kredit, Geld, Werkzeuge, Samen, Produkte, Markt, Traum, Dorf, Arbeit}.
\item Futur I : \all{werden} (2ᵉ position) + infinitif (fin).
\item \all{um ... zu} : virgule avant \all{um}, \all{zu} + infinitif (fin) ; \all{anzubauen} (pas \all{zu anbauen}).
\item Ne pas calquer le français : « créer une ferme » = \all{eine Farm gründen / aufbauen} (pas \all{*kreieren}) ; « parler à » = \all{mit ... sprechen} (pas \all{*sprechen zu}).
\end{itemize}
\end{corrige}

\newpage

\coursec{Fiche bilan}

\begin{popfigure}
\begin{tikzpicture}[
    mindmap,
    concept color=popblue,
    level 1 concept/.append style={level distance=3.5cm, sibling angle=60},
    level 2 concept/.append style={level distance=2.5cm},
    every node/.style={font=\small, align=center},
    branch1/.style={concept color=popgreen},
    branch2/.style={concept color=poporange},
    branch3/.style={concept color=poppurple},
    branch4/.style={concept color=poppink},
    branch5/.style={concept color=popgold},
    branch6/.style={concept color=popblue!70!popgreen}
]
\node[concept, text=white, scale=1.1, align=center]{Leçon ALA08\\Création d'une\\activité agropastorale}
    child[branch1] { node[concept, text=white, align=center]{Texte support\\„Jonas gründet\\eine Farm“} }
    child[branch2] { node[concept, text=white, align=center]{Vocabulaire\\Ferme \& Élevage\\die Farm, das Vieh,\\anbauen, züchten} }
    child[branch3] { node[concept, text=white, align=center]{Grammaire 1\\Futur I\\werden + Infinitif} }
    child[branch4] { node[concept, text=white, align=center]{Grammaire 2\\Finalité\\um ... zu + Inf.} }
    child[branch5] { node[concept, text=white, align=center]{Expression\\Réponses, Résumé,\\Production guidée} }
    child[branch6] { node[concept, text=white, align=center]{Contexte\\Cameroun\\Cacao, Maïs,\\Élevage local} };
\end{tikzpicture}
\legende{Carte mentale récapitulative de la leçon ALA08}
\end{popfigure}

\courssub{Lexique clé (Allemand / Français)}

\begin{center}
\begin{tabularx}{\linewidth}{|X|X|X|}
\hline
\rowcolor{popblueL} \textbf{Deutsch (Artikel, Plural)} & \textbf{Français} & \textbf{Famille / Remarque} \\
\hline
\all{die Farm, die Farmen} & la ferme (exploitation) & \cle{Substantiv}, fém. \\
\all{der Bauer, die Bauern} & le paysan, l'agriculteur & \cle{Substantiv}, masc. (n-décl.) \\
\all{das Vieh (kein Plural)} & le bétail, le cheptel & \cle{Substantiv}, n., collectif \\
\all{das Rind, die Rinder} & le bovin, la vache, le bœuf & \cle{Substantiv}, n. \\
\all{das Schwein, die Schweine} & le porc, le cochon & \cle{Substantiv}, n. \\
\all{das Huhn, die Hühner} & la poule, le poulet & \cle{Substantiv}, n. \\
\all{der Mais (kein Plural)} & le maïs & \cle{Substantiv}, masc. \\
\all{das Feld, die Felder} & le champ & \cle{Substantiv}, n. \\
\all{anbauen, baut an, hat angebaut} & cultiver (quelque chose) & \cle{Verb}, séparable, fort \\
\all{züchten, züchtet, hat gezüchtet} & élever, sélectionner & \cle{Verb}, faible \\
\all{gründen, gründet, hat gegründet} & fonder, créer & \cle{Verb}, faible \\
\all{investieren, investiert, hat investiert} & investir & \cle{Verb}, faible \\
\all{der Ertrag, die Erträge} & le rendement, la récolte & \cle{Substantiv}, masc. \\
\all{der Gewinn, die Gewinne} & le bénéfice, le profit & \cle{Substantiv}, masc. \\
\all{der Kredit, die Kredite} & le crédit, le prêt & \cle{Substantiv}, masc. \\
\hline
\end{tabularx}
\end{center}

\begin{exemplebox}[Le lexique en contexte]
\all{Der Bauer baut Mais auf seinen Feldern an und züchtet Rinder und Schweine.} « Le paysan cultive du maïs dans ses champs et élève des bovins et des porcs. »\\
\all{Mit einem Kredit der Bank will er seine Farm gründen.} « Avec un crédit de la banque, il veut créer sa ferme. »\\
\all{Er hofft auf einen guten Ertrag und einen hohen Gewinn.} « Il espère un bon rendement et un bénéfice élevé. »
\end{exemplebox}

\courssub{Règles de grammaire à connaître par cœur}

\begin{propriete}[Formation du Futur I]
\textbf{Forme :} \all{werden} (conjugué au présent, en 2\textsuperscript{e} position) + \textbf{Infinitif du verbe principal} (en position finale).\\
\textbf{Conjugaison de \all{werden} (Präsens) :}
\begin{center}
\begin{tabular}{|l|l|l|l|}
\hline
\rowcolor{popblueL} Person & Singular & Person & Plural \\
\hline
ich & \all{werde} & wir & \all{werden} \\
du & \all{wirst} & ihr & \all{werdet} \\
er/sie/es & \all{wird} & sie/Sie & \all{werden} \\
\hline
\end{tabular}
\end{center}
\textbf{Exemples :} \all{Ich werde einen Hof gründen.} « Je fonderai une ferme. » --- \all{Wir werden Mais anbauen.} « Nous cultiverons du maïs. » --- \all{Wirst du im Dorf bleiben?} « Resteras-tu au village ? »
\end{propriete}

\begin{propriete}[Emploi du Futur I]
\begin{itemize}
    \item \cle{Prédiction / Supposition} : \all{Es wird morgen regnen.} « Il va pleuvoir / Il pleuvra demain. »
    \item \cle{Intention / Promesse} (souvent avec \all{sicher}, \all{bestimmt}) : \all{Ich werde das Geld sparen.} « Je vais économiser l'argent. »
    \item \cle{Futur proche} : l'allemand préfère souvent le présent + indicateur de temps : \all{Morgen fahre ich nach Douala.} « Demain je vais à Douala. » Le Futur I reste correct : \all{Morgen werde ich nach Douala fahren.}
\end{itemize}
\end{propriete}

\begin{propriete}[La finalité : \all{um ... zu} + Infinitif]
\textbf{Structure :} \all{, um} + [compléments] + \textbf{zu} + \textbf{Infinitif} (en toute fin de phrase).\\
Le sujet de la proposition infinitive est \emph{identique} à celui de la principale.\\
\textbf{Exemples :} \\
\all{Jonas spart Geld, um eine Farm zu gründen.} « Jonas économise de l'argent pour fonder une ferme. » \\
\all{Wir lernen Deutsch, um in Deutschland zu arbeiten.} « Nous apprenons l'allemand pour travailler en Allemagne. » \\
\all{Er kauft Saatgut, um Mais anzubauen.} « Il achète des semences pour cultiver du maïs. » (verbe séparable : \all{zu} entre la particule et le radical)
\end{propriete}

\begin{attention}[Pièges fréquents (Francophones)]
\begin{itemize}
    \item \textbf{Oubli du \all{zu}} : *\all{um eine Farm gründen} $\rightarrow$ \all{um eine Farm \textbf{zu} gründen}.
    \item \textbf{Sujet différent} : si le sujet change, on ne peut pas utiliser \all{um ... zu} ; il faut une subordonnée : \all{Jonas spart Geld, damit seine Kinder zur Schule gehen können.} « Jonas économise pour que ses enfants puissent aller à l'école. »
    \item \textbf{Verbe séparable} : \all{zu} s'insère \emph{entre} la particule et le radical : \all{um anzufangen}, \all{um mitzuarbeiten}, \all{um anzubauen}.
    \item \textbf{Futur I vs Présent} : en allemand, le présent + indicateur de temps (\all{morgen, nächstes Jahr}) suffit souvent pour le futur proche. Le Futur I marque plutôt la \cle{volonté}, la \cle{promesse} ou la \cle{probabilité}.
    \item \textbf{Place de l'infinitif} : *\all{Ich werde gründen eine Farm} $\rightarrow$ \all{Ich werde eine Farm \textbf{gründen}.} L'infinitif se place en toute fin.
    \item \textbf{Faux amis} : \all{der Kredit} = le crédit bancaire (pas « le crédit » au sens de réputation) ; \all{die Farm} = la ferme/exploitation (le mot « ferme » au sens de « solide » se traduit par \all{fest}).
\end{itemize}
\end{attention}

\courssub{Méthodes express}

\begin{methode}[Comprendre un texte agropastoral]
\begin{enumerate}
    \item \textbf{Anticiper} : titre, images $\rightarrow$ hypothèses (vocabulaire attendu : \all{Farm, Vieh, anbauen, Kredit}).
    \item \textbf{Lecture globale} : Qui ? (Jonas) Quoi ? (création d'une ferme) Où ? (Allemagne / Cameroun) Quand ? (projet, futur).
    \item \textbf{Lecture détaillée} : souligner les verbes (Futur I, \all{um ... zu}), les connecteurs (\all{zuerst, dann, deshalb, schließlich}), les chiffres.
    \item \textbf{Glossaire personnel} : noter 5 à 7 mots nouveaux avec article et pluriel.
\end{enumerate}
\end{methode}

\begin{methode}[Traduire Thème / Version (Futur I \& Finalité)]
\begin{enumerate}
    \item Identifier le temps : futur simple français $\rightarrow$ allemand : Futur I \textbf{ou} Présent + adverbe de temps.
    \item Repérer la finalité : « pour + infinitif » / « afin de + infinitif » $\rightarrow$ allemand : \all{um ... zu} (si le sujet est identique).
    \item Vérifier l'ordre des mots : verbe conjugué en 2\textsuperscript{e} position ; \all{werden} conjugué + infinitif principal en fin de phrase ; dans la subordonnée, \all{zu} + infinitif en toute fin.
\end{enumerate}
\end{methode}

\begin{methode}[Produire un court paragraphe (5 à 7 phrases)]
\begin{enumerate}
    \item \textbf{Plan} : 1. Présentation du projet (Qui, Quoi). 2. Moyens (terrain, crédit). 3. Actions (cultures, élevage --- avec \all{um ... zu}). 4. Objectifs (revenus, emplois --- au Futur I).
    \item \textbf{Rédiger} : phrases courtes, liaisons logiques (\all{zuerst, dann, deshalb, schließlich}).
    \item \textbf{Relire} : accords, cas (accusatif pour l'objet direct), place de \all{zu}, majuscules à tous les noms.
\end{enumerate}
\end{methode}

\begin{exoresolu}[Thème : Futur I et finalité]
Traduisez en allemand : « L'année prochaine, Jonas achètera des poules pour vendre des œufs au marché. »
\tcblower
\textbf{Solution détaillée :}\\
1. « L'année prochaine » = \all{nächstes Jahr} (indicateur de temps, placé en tête $\rightarrow$ inversion du sujet).\\
2. « achètera » = futur $\rightarrow$ Futur I : \all{wird ... kaufen} (infinitif en fin de phrase).\\
3. « des poules » = \all{Hühner} (accusatif pluriel).\\
4. « pour vendre » = finalité, même sujet $\rightarrow$ \all{um ... zu verkaufen}.\\
5. « des œufs » = \all{Eier} ; « au marché » = \all{auf dem Markt}.\\
\textbf{Traduction :} \all{Nächstes Jahr wird Jonas Hühner kaufen, um Eier auf dem Markt zu verkaufen.}
\end{exoresolu}

\begin{aretenir}[Checklist avant le Bac]
\begin{itemize}
    \item $\square$ Je sais conjuguer \all{werden} au présent sans hésiter.
    \item $\square$ Je forme le Futur I correctement : \all{werden} (2\textsuperscript{e} place) + infinitif (fin de phrase).
    \item $\square$ Je distingue le Futur I (volonté, promesse, prédiction) et le Présent + indicateur de temps (futur proche planifié).
    \item $\square$ Je construis une phrase de finalité avec \all{um ... zu} + infinitif (sujet identique).
    \item $\square$ Je place \all{zu} correctement devant l'infinitif (verbes simples et inséparables) ou entre la particule et le radical (verbes séparables : \all{anzubauen}).
    \item $\square$ Je connais le vocabulaire de base : \all{die Farm, der Bauer, das Vieh, der Mais, das Feld, anbauen, züchten, gründen, der Kredit, der Gewinn} (genre + pluriel).
    \item $\square$ Je sais répondre en allemand à des questions simples sur un texte (\all{Wer? Was? Wo? Wann? Warum?}).
    \item $\square$ Je sais résumer un texte de 150 mots en 5 à 6 phrases (Présent / Perfekt / Futur I).
    \item $\square$ J'écris une courte production (projet agropastoral) structurée avec des connecteurs logiques.
    \item $\square$ Je vérifie les majuscules sur \textbf{tous} les noms et le point final.
    \item $\square$ Je repère les faux amis : \all{das Gift} (le poison, pas le cadeau), \all{der Chef} (le patron), \all{die Farm} (la ferme, pas « ferme » au sens de solide).
\end{itemize}
\end{aretenir}

\findecours

\end{document}
